\chapter{几何与表面重建}

\paragraph*{本章阅读锚点}
原始 3DGS 的几何来自 SfM 点初始化、位置梯度和透明度混合，它擅长渲染，却不天然等价于一个干净表面。高斯可以漂浮、重叠、拉长，也可能为了拟合高光或纹理在真实表面后方生成伪几何。因此本章的读法是先看几何锚点：SuGaR 把高斯压向表面并提网格，2DGS 把体高斯改成 surfel，DN-Splatter/PGSR/RaDe-GS 改深度和法线监督，GSDF/GSurf/GaussianUDF 用隐式场补连续表面。后续方法大多是在这些锚点上改一个环节：深度是否无偏、法线是否可信、透明/光泽表面如何处理、网格怎样提取、以及曲线和平面这类结构先验如何进入高斯。

\section{深度与法线正则化}

深度/法线路线不一定改变高斯表示本身，它们主要解决“光度损失不等于几何正确”。原始 3DGS 可以用错误几何拟合训练图，特别是在无纹理、反射、透明或稀疏视角区域。下面这些方法通常加入单目深度、法线、置信度、无偏深度或视图一致约束，让高斯更贴近可提取表面。



% METHOD: 2403.17822
\paragraph{DN-Splatter（深度与法线先验）（2024.03）}
DN-Splatter~\cite{2403_17822} 是室内几何正则化的代表性锚点。它继承原始 3DGS 的显式高斯、可微光栅化和 ADC 密度控制，但改变了监督信号：原始 3DGS 主要靠 RGB 重建损失移动高斯，室内无纹理墙面、弱视差区域和物体边界很容易出现漂浮高斯或错误深度；DN-Splatter 把深度、法线和局部平滑写入优化，让高斯不仅能拟合图像，也要接近可提取表面。

它的关键不是把单目深度当成硬标签。方法先让渲染器同时输出颜色和深度，再把传感器深度或 ZoeDepth、Depth Anything 等单目深度与渲染深度对齐；深度损失按当前 RGB 图像梯度调权，平滑区域加强约束，边缘区域降权，避免 iPhone 深度或单目深度在轮廓处的误差把表面拉坏。法线部分则利用高斯协方差中最短轴近似局部表面法线，并与 Omnidata、DSINE 等网络预测的法线对齐；额外的 normal smoothing loss 约束相邻像素的渲染法线，使局部表面不要因为单个高斯自由旋转而变得破碎。

训练流程仍是 Nerfstudio/Splatfacto 式的 3DGS 优化，但初始化可以使用传感器深度反投影得到更密集的起始点云，损失函数在 RGB 重建之外加入梯度感知深度项、法线项和局部平滑项。优化完成后，方法从渲染深度和法线反投影得到 oriented point set，再用泊松表面重建提取 mesh。相对 SuGaR 主要从高斯分布本身推表面，DN-Splatter 的改变是把外部几何先验接进训练；相对普通深度监督，它强调“边界处少信深度、平滑处多信深度”。
\textbf{锚点细化：} DN-Splatter：针对纯深度先验只约束位置不约束朝向的局限，用"深度与法线联合先验"替代纯深度正则化，核心思想是利用预训练单目深度/法线估计器分别约束高斯位置与朝向（法线由高斯协方差矩阵最短轴推导），让高斯从空中漂浮的椭球变成贴合表面的扁平薄片，是 3DGS 几何正则化的里程碑工作。
\textbf{方法信息：} 代码：\url{https://github.com/maturk/dn-splatter}；GitHub 关注度高，约 821 stars；引用很多，Semantic Scholar 约 187 次。



% METHOD: 2406.01467
\paragraph{RaDe-GS（通用 3D 高斯深度光栅化）（2024.06）}
RaDe-GS~\cite{2406_01467} 的锚点价值在于：它不把 3D 高斯改成 2D surfel，而是为通用 3D Gaussian 设计深度和法线渲染。许多几何方法依赖 2DGS，是因为 2D surfel 天然有平面和法线；RaDe-GS 试图保留原始 3DGS 的渲染效率和外观质量，同时得到更准确的几何图。

方法把每个高斯放到 ray space 中分析：对一个像素射线，不再简单取高斯中心深度，而是推导该射线与高斯局部椭球贡献最大位置对应的深度。这个深度可写成高斯中心深度加上与像素偏移相关的校正项，因此同一个投影椭圆内部不同像素可以得到不同深度。法线也从高斯局部几何和射线交互位置推导，而不是把每个高斯的旋转轴粗略混合。训练时，RaDe-GS 可以像 RGB 渲染一样快速产生 depth map 和 normal map，再把深度/法线损失接入 3DGS 优化。它可作为模块接入其他 3DGS 方法，为多视图三角测量或法线正则提供更可信几何；后续 TriaGS 就利用 RaDe-GS 思路得到几何一致深度。
\textbf{实现细节：} DN-Splatter 使用标准 $\alpha$-混合 渲染深度和法线，但"期望深度"本身就有问题。RaDe-GS（2024.06） 修改了原始 3DGS 的深度/法线渲染方式，用精确的封闭形式公式替代期望深度的近似计算，训练流程保持不变。其核心贡献在于提出了精确的高斯深度和法线光栅化方法，揭示了标准 3DGS 深度渲染的固有缺陷并给出封闭形式的修正方案。问题：期望深度忽略了高斯的体积分布。标准 3DGS 只渲染颜色和不透明度，深度图通常通过 $\alpha$-混合 的期望深度 $\hat{z} = \sum_i \omega_i z_i$ 近似计算，其中 $\omega_i = \alpha_i \prod_{j=1}^{i-1}(1-\alpha_j)$ 是混合权重，$z_i$ 是第 $i$ 个高斯中心的深度。RaDe-GS 指出这种做法将每个高斯基元视为一个"点"：只用其中心深度代表整个基元，完全忽略了高斯在三维空间中的体积分布。实际上，不同像素的视线射线与同一个高斯基元相交于不同深度，因此深度应该是像素位置的函数，而非常数。精确的深度光栅化。简单来说，RaDe-GS 发现射线穿过高斯时，不同像素位置看到的"表面"深度其实不同——就像一个椭球气球，从不同位置看气球表面，深度各不相同。RaDe-GS 的关键洞察是：在 3DGS 采用的局部仿射投影模型下，射线与高斯基元的"交点"（定义为使高斯值最大的点）的深度可以封闭形式地计算。具体推导在"光线空间"（ray space）中完成。3DGS 将每个高斯从世界坐标系通过仿射变换投影到屏幕空间，得到 2D 高斯协方差 $\boldsymbol{\Sigma}' = \mathbf{J}\mathbf{W}\boldsymbol{\Sigma}\mathbf{W}^{}\top\mathbf{J}^{}\top$，其中 $\mathbf{W}$ 是世界到相机的旋转矩阵，$\mathbf{J}$ 是透视变换的雅可比矩阵（Jacobian Matrix）。RaDe-GS 进一步在光线空间中考虑第三个坐标维度 $t$（点到相机的距离），定义一个 $1\times 3$ 常数向量： 其中 $\mathbf{v}'=(0,0,1)^{}\top$ 是光线空间中的射线方向常数。对每个高斯，$\hat{\mathbf{q}}$ 可预计算。
\textbf{锚点细化：} RaDe-GS：针对 $\alpha$-混合 期望深度将高斯基元视为"点"忽略体积分布的局限，用"光线空间中的封闭形式深度与法线"替代期望深度近似，核心思想是在仿射投影模型下推导射线与高斯"交点"的解析公式——每个像素的精确深度由高斯中心深度加像素偏移修正项构成，且射线-高斯交点构成平面，平面法线即可作为像素法线直接光栅化。
\textbf{方法信息：} 代码：\url{https://github.com/HKUST-SAIL/RaDe-GS}；GitHub 关注度很高，约 687 stars；引用很多，Semantic Scholar 约 130 次。



% METHOD: 2406.05774
\paragraph{VCR-GauS（视图一致深度-法线正则）（2024.06）}
VCR-GauS~\cite{2406_05774} 解决法线监督只更新旋转参数的问题。普通做法把高斯法线和预测法线对齐，但位置和尺度未必随之改善；外部单目法线估计器在多视图之间也可能不一致。VCR-GauS 用渲染深度的梯度推导法线，使法线损失能影响位置、尺度和其他几何参数。

它提出 D-Normal 公式：先用 scale regularization 把 3D 高斯压成接近平面的形状，并把高斯表示为平面点和法线；对每条像素射线，计算射线与该平面的交点深度，而不是取高斯中心深度；随后由深度梯度得到法线。由于这个交点深度同时依赖位置、尺度和法线，D-Normal 损失能回传到更多几何参数。多视图法线预测不一致处，方法根据渲染法线与预测法线的差异计算置信度，降低不可靠区域的监督权重。它还通过 densification/splitting 控制高斯大小和分布，减少表面 bump。相对 DN-Splatter 的外部法线对齐，VCR-GauS 更强调法线应从深度几何中反推。
\textbf{实现细节：} 单独的深度或法线正则化各有盲点：深度约束位置但不约束朝向，法线约束朝向但梯度无法传到位置。联合正则化旨在打通这两个维度。VCR-GauS（2024.06） 在原始 3DGS 的训练环节引入"交点深度"替代高斯中心深度，使法线监督的梯度能够传导到高斯位置参数，光度损失保持不变。动机。法线仅与高斯旋转参数相关，法线监督无法有效更新高斯位置——即"知道表面朝向不对"却"无法移动高斯到正确位置"，需要一座桥梁打通法线与位置的梯度通路。所谓"交点深度"，是指将扁平高斯视为平面、求光线与平面的交点深度。它同时依赖于高斯的法线（决定平面朝向）和位置（决定平面在哪里），因此法线监督的梯度可以"顺道"推动位置更新。VCR-GauS 指出 3DGS 中仅监督法线不足以驱动表面重建，原因是法线仅与高斯的旋转参数相关，法线监督无法有效更新高斯的位置。NeuS-based 方法能将法线正则化传导至位置更新（法线即 SDF 的梯度），但 3DGS 缺乏这种机制。为此，VCR-GauS 提出交点深度（intersection depth）替代高斯中心深度：先将高斯挤压为扁平形状（通过尺度正则化 $\mathcal{L}_s = \|(s_1, s_2, s_3)\|_1$ 使最小尺度分量趋零），再将扁平高斯视为平面 $(p, n)$，求光线与平面的交点深度 $d(n, p) = r_z \cdot (n \cdot p) / (n \cdot r)$，其中 $r$ 为光线方向、$r_z$ 为其 $z$ 分量。交点深度同时依赖于法线 $n$ 和位置 $p$，从而打通了法线监督到位置更新的梯度通路。简言之，以前法线梯度只能更新高斯的旋转，现在通过"交点深度"这个桥梁，法线梯度也能推动高斯移动到正确位置。
\textbf{锚点细化：} VCR-GauS：针对法线仅与高斯旋转参数相关导致法线监督无法更新位置的局限，用"交点深度"替代高斯中心深度作为正则化目标，核心思想是先把高斯压扁为平面，再求光线与平面的交点深度——该深度同时依赖法线（决定平面朝向）和位置（决定平面位置），从而打通法线监督到位置更新的梯度通路。
\textbf{设计动机与边界：} 相对原始 3DGS 以独立高斯和图像重建误差为主的设定，VCR-GauS 把问题落在以下可实现的改动上：VCR-GauS解决两个关键问题：(1) 传统法线监督仅有效更新旋转参数，对位置等其他几何参数影响有限。解决方法是引入D-Normal公式，法线从渲染深度的梯度推导而非直接混合3D高斯法线，深度计算为射线与压缩高斯的交点而非高斯中心位置。(2) 多视图法线预测不一致导致严重重建伪影。引入置信度项降低不一致区域的法线监督权重。额外贡献：尺度正则化压缩高斯椭球为平面；自适应密度化和分裂策略控制高斯大小和分布。整体框架实现法线、深度和位置的联合优化。其设计目的不是单独增加一个模块，而是用上述表示、监督或渲染约束消除原流程在目标任务中的具体失配。提出深度-法线正则化器，将法线与几何参数直接耦合，实现法线正则化对位置等其他几何参数的完整更新。引入置信度项缓解多视图法线预测不一致问题。提出密度化和分裂策略正则化高斯的大小和分布。在保持100+FPS渲染速度和有竞争力的外观质量下，实现更好的重建质量。适用时还应保留论文给出的边界：依赖外部法线估计器的质量；对于高度反射或透明表面仍具挑战。
\textbf{方法信息：} 代码：\url{https://github.com/HLinChen/VCR-GauS}；GitHub 关注度很高，约 170 stars；引用很多，Semantic Scholar 约 62 次。



% METHOD: 2409.07456
\paragraph{Self-Evolving DS-GS（渲染立体对自进化深度）（2024.09）}
Self-Evolving Depth-Supervised 3DGS~\cite{2409_07456} 的关键想法是让 3DGS 自己产生逐步变好的深度监督。固定单目深度或 SfM 深度在训练早期能提供几何方向，但它们的误差不会随着 3DGS 变好而更新；如果一直相信同一张深度图，后期可能限制几何细节。

方法先系统比较 SfM 深度、单目深度、深度补全和多视图立体四类先验，然后把最有效的立体深度做成训练闭环。训练过程中，当前 3DGS 定期从相邻虚拟相机渲染左右立体对；这些图像送入 RAFT-Stereo 一类预训练 stereo 网络，得到新的深度估计，再作为额外深度损失回到 3DGS 优化。由于虚拟立体对来自当前模型，模型越清晰，后续 stereo 深度也越可靠，形成“渲染更好图像--估计更好深度--继续约束几何”的循环。相对 CDGS 的置信度加权固定深度，Self-Evolving DS-GS 的差异是深度监督本身随训练状态更新。
\textbf{方法信息：} 代码：\url{https://github.com/sadrasafa/stereogs}。
\textbf{设计动机与边界：} 相对原始 3DGS 以独立高斯和图像重建误差为主的设定，Self-Evolving Depth-Supervised 3D Gaussian Splatting from Rendered Stereo Pairs 把问题落在以下可实现的改动上：工作分为两部分：1）系统性研究四种深度先验整合策略在GS优化中的效果：SfM深度、单目深度估计、深度补全和多视图立体；2）提出Self-Evolving策略：在GS训练过程中，模型定期渲染虚拟立体对，送入预训练立体网络（如RAFT-Stereo）获取深度估计，作为额外监督信号。随着GS训练进行，渲染质量提升，立体网络得到的深度估计也愈加准确，形成正向循环。无需外部深度传感器。其设计目的不是单独增加一个模块，而是用上述表示、监督或渲染约束消除原流程在目标任务中的具体失配。全面分析深度先验在GS优化中的整合方式，提出Self-Evolving DS-GS策略：GS模型在训练过程中自身渲染虚拟立体对，通过现成立体网络提取深度线索，动态利用这些深度信号持续改进场景表示。首次系统评估GS模型的深度精度。适用时还应保留论文给出的边界：依赖立体网络的泛化能力；对于纹理极少的场景效果受限；额外渲染步骤增加训练时间。
% METHOD: 2502.07840
\paragraph{TranSplat（透明物体表面嵌入）（2025.02）}
TranSplat~\cite{2502_07840} 面向机器人场景中的透明物体深度补全。普通 RGB-D 或 3DGS 管线在这里会同时遇到两个问题：深度传感器经常在玻璃、透明塑料上失效，RGB 又受反射和折射影响，直接把透明物体当朗伯表面会让高斯 opacity 退化，表面深度变得不连续。

它的增量不是先改高斯形状，而是先为透明表面构造一个更稳定的 2D 代理信号。第一阶段，方法用 Surface Embedding 数据训练表面嵌入模型，再用 latent diffusion/ControlNet 结构从裁剪 RGB 图生成表面嵌入图；这个嵌入更接近物体真实表面颜色化后的几何线索，受视角和光照影响小于原始 RGB。第二阶段，3DGS 优化同时使用 RGB 图和表面嵌入图：RGB 负责外观一致性，表面嵌入负责把透明区域的高斯拉回连续表面，并缓解透明处 opacity 坍塌到零的问题。最终输出用于透明物体的稠密深度补全和抓取。相对 TSGS 的去光照和首表面深度，TranSplat 更依赖学习到的表面嵌入先验，并把它接入 Gaussian Splatting 优化流程。
\textbf{实现细节：} 前述所有正则化方法都假设场景表面是不透明的——高斯的不透明度越高越好。透明物体打破了这一假设：表面的不透明度本应很低，导致标准深度渲染失效。TranSplat（2025.02） 修改了原始 3DGS 的渲染通道，为每个高斯同时维护 RGB 颜色和表面嵌入两组属性，实现双通道渲染，训练流程保持不变。TranSplat 专门解决透明物体的深度重建问题。透明物体作为非朗伯表面，对光照和视角变化高度敏感，直接用 RGB 图像进行 3DGS 优化会面临两个核心困难：其一，透明表面的光度一致性在不同视角间极不可靠，导致深度图充满噪声和伪影；其二，由于透明物体"不可见"，标准 3DGS 训练中代表透明表面的高斯不透明度 $\alpha$ 会坍缩为零，深度渲染出现大量孔洞。TranSplat 的核心洞察是：既然 RGB 图像本身不可靠，就需要一种对光照和视角不敏感的替代特征来监督透明表面的重建——即表面嵌入（surface embedding）。
\textbf{方法信息：} 官方代码：\url{https://github.com/jeongyun0609/TranSplat}；GitHub 约 42 stars。



% METHOD: 2502.14684
\paragraph{CDGS（置信度感知深度正则）（2025.02）}
CDGS~\cite{2502_14684} 仍沿用“给 3DGS 加深度监督”的路线，但它修正的是深度先验的可信度问题。单目深度在大平面上有用，在物体边缘、反射区域或低纹理区域却容易给错深度；如果整张图用同一个权重监督，几何会被错误先验拉偏。

它的流程先用 Depth Anything V2 预测稠密深度，再用 SfM 稀疏点深度求尺度和偏移，把单目深度对齐到重建场景坐标系。随后方法同时看深度自身的对齐残差和 RGB 图像特征，生成逐像素置信图。训练 3DGS 时，颜色损失仍是标准 L1/SSIM，深度正则只在置信度高的位置给强约束；在边缘、反射和深度预测不稳定区域，损失权重自动降低。相对 DN-Splatter 的图像梯度启发式，CDGS 的改变是把“哪里该相信深度”做成多线索置信估计；相对 GVGS 的可见性感知约束，它仍主要服务于单目深度正则。
\textbf{方法信息：} 代码：\url{https://github.com/zqlin0521/cdgs-release}。
\textbf{设计动机与边界：} 相对原始 3DGS 以独立高斯和图像重建误差为主的设定，CDGS 把问题落在以下可实现的改动上：CDGS包含两个核心组件：深度细化与对齐、置信度感知深度正则化。首先使用Depth Anything V2模型获取初始深度图，通过梯度下降方法将其与SfM稀疏深度对齐。对于每张对齐后的深度图，通过分析深度图和对应RGB图像的特征生成置信图。在优化过程中，除了原始3DGS的L1和SSIM损失外，引入额外的深度正则化项。该正则化包含深度归一化和基于置信图的动态损失调整机制。置信图综合考虑深度对齐质量、图像特征等因素，使深度监督在可靠区域增强、不可靠区域减弱。方法从2D和3D两个角度全面评估，并分析了整个优化过程的性能变化。其设计目的不是单独增加一个模块，而是用上述表示、监督或渲染约束消除原流程在目标任务中的具体失配。提出CDGS方法，通过多线索置信图自适应调整深度监督强度，在NVS质量上提升达2.31 dB PSNR，仅需50%训练迭代即可达到原始3DGS相当的F-score。适用时还应保留论文给出的边界：依赖预训练的单目深度估计模型质量；置信图生成增加了额外的计算开销；在极端光照或弱纹理场景下的鲁棒性需进一步验证。
% METHOD: 2503.06587
\paragraph{Unbiased Depth for 2DGS（无偏深度与光泽孔洞）（2025.03）}
Introducing Unbiased Depth into 2D Gaussian Splatting~\cite{2503_06587} 是对 2DGS 深度定义的修正。2DGS 已经把基元改成表面 surfel，但在光泽表面上仍会出现孔洞：高光颜色随视角突变，低阶 SH 难以拟合，优化器就可能在真实表面后方放置额外高斯去解释反射，最终让同一条射线上的深度贡献分裂。

方法做了两处替换。第一，用 depth convergence loss 替代原来的 depth distortion loss，直接约束同一像素射线附近贡献高斯的中心深度收敛到一起，减少后方伪高斯造成的深度分叉。第二，深度渲染不再取累积 opacity 到达阈值的中位点，而是综合该射线上所有相交 surfel 的贡献，得到更少偏置的深度估计。相对 PGSR 的无偏平面深度，它专门解释 2DGS 在光泽区域的孔洞来源，并把修正接在 2DGS 自身的深度损失和深度估计公式上。
\textbf{实现细节：} 纯深度先验只约束高斯的位置，不约束其朝向和形状。法线先验填补了这一空白。DN-Splatter（2024.03） 在原始 3DGS 的训练环节引入深度与法线联合先验，利用预训练单目深度/法线估计器为每个视角生成伪标注，并将其作为几何正则化项加入训练，光度损失保持不变。动机。标准 3DGS 仅靠光度损失无法约束高斯的几何形状，导致高斯退化成飘浮在空中的椭球而非贴合表面的扁平薄片——深度约束位置、法线约束朝向，两者缺一不可。DN-Splatter 是最早将深度和法线先验系统引入 3DGS 的方法之一，也是本节的里程碑工作。DN-Splatter 的核心思路是利用预训练的单目深度估计器（如 ZoeDepth、DepthAnything）和法线估计器（Omnidata）为每个训练视角生成伪深度图和法线图，将它们作为几何正则化项加入训练。深度正则化。在 DN-Splatter 中，深度正则化的实现方式是以标准 $\alpha$-混合 方式渲染深度图 $\hat{D} = \sum_{i \in N} d_i \alpha_i \prod_{j=1}^{i-1}(1-\alpha_j)$，其中 $d_i$ 是第 $i$ 个高斯在视图空间中的 $z$ 深度。然后用梯度感知的深度损失将渲染深度与估计深度对齐： 其中 $g_{\mathrm{rgb}} = \exp(-\nabla I)$ 是由当前 RGB 图像梯度驱动的边缘感知权重。直观来说，在图像梯度大的区域（物体边缘），深度监督被自动削弱；在平坦无纹理区域，深度监督被增强，因为这些区域仅靠光度损失很难约束几何。此外，对数惩罚 $\log(1+\cdot)$ 比线性或二次惩罚产生更平滑的重建。对于没有传感器深度的数据集，DN-Splatter 通过将单目深度与稀疏 SfM 点对齐来消除尺度歧义：对每帧求解线性回归 $\hat{a}, \hat{b} = \operatorname{argmin}_{a,b} \sum_{ij} \|(a \cdot D_{\mathrm{mono},ij} + b) - D_{\mathrm{sparse},ij}\|_2^{}2$。法线正则化。DN-Splatter 的法线直接从高斯的几何参数推导。
\textbf{设计动机与边界：} 相对原始 3DGS 以独立高斯和图像重建误差为主的设定，Introducing Unbiased Depth into 2D Gaussian Splatting for High-accuracy Surface Reconstruc 把问题落在以下可实现的改动上：本文首先分析了2DGS在光泽表面产生孔洞的机理：光泽表面在特定视角产生高光，球谐函数无法拟合剧烈变化，优化过程在高光后方创建高斯以适配反射跳变，导致深度不连续和孔洞。提出两个改进：(1)深度收敛损失——替代原始深度失真损失（使用透射率权重导致约束不足），新损失直接约束沿光线方向相邻高斯中心之间的距离，强约束深度连续性，惩罚不连续的深度分布；(2)无偏深度判据——原始2DGS使用累积不透明度达到0.5的中位点作为表面深度，在光泽区域会偏向高光高斯。新方法综合考虑光线上所有相交高斯的贡献，加权计算无偏深度估计。两个改进配合使用显著改善了光泽和半光泽区域的重建完整性。为了消除原始 3DGS 在该目标任务中的具体失配，方法把上述表示、监督或渲染约束作为可学习和可执行的改动，而不是单独堆叠模块。发现2DGS在光泽表面产生孔洞的原因是反射不连续性引入深度偏差，提出深度收敛损失替代深度失真损失以强约束深度连续性，并修正深度判据全面考虑光线上的所有相交高斯，显著提升重建表面完整性和精度。适用时还应保留论文给出的边界：对镜面反射等极端高光效果仍可能有限；深度收敛损失的超参数需根据场景调节；计算开销略有增加。
\textbf{方法信息：} 代码：\url{https://github.com/xiaoxinyyx/unbiased_surfel}。
% METHOD: 2504.12799
\paragraph{TSGS（透明表面深度与去光照先验）（2025.04）}
TSGS~\cite{2504_12799} 针对透明玻璃器皿这类最难几何。透明表面会把背景和反射一起混入图像，原始 alpha 深度常被背景高斯主导，透明度也容易坍塌。TSGS 的改变是把几何学习和外观优化分开。

第一阶段用去光照图像和法线先验抑制镜面/反射干扰，专注学习透明表面几何。第二阶段冻结已建立的透明度，使用原始图像和各向异性球面高斯优化镜面外观，避免外观拟合破坏几何。第三阶段提出首表面深度提取：通过累积透射率滑动窗口定位最可能表面，而不是用普通 alpha 混合深度。相对 TranSplat 透明物体深度补全，TSGS 更强调 3DGS 内部的几何-外观分阶段训练。
具体流程上，第一阶段用去光照图像和法线先验训练几何，尽量让高斯先学到透明物体的首表面。第二阶段在保持几何的前提下优化外观：冻结已经形成的透明度结构，用原始图像和各向异性球面高斯（ASG）建模镜面反射，使外观分量不再反向破坏表面位置。第三阶段重新定义深度提取，不直接用标准 alpha 混合深度，而是在累积透射率上做滑动窗口，定位最可能的首表面区域，再计算加权平均深度。相对 TranSplat 的表面嵌入先验，TSGS 更强调 3DGS 内部的“先几何、后外观、再首表面深度”流程，并配套构建了透明实验室器皿数据集 TransLab。
\textbf{实现细节：} TranSplat 和 TRAN-D 都在渲染环节引入额外通道，但没有专门处理"透明表面的深度渲染本身有偏"这一根本问题。TSGS（2025.04） 修改了原始 3DGS 的训练流程，将优化分为几何学习、外观精调和首表面提取三个独立阶段，并在推理时用滑动窗口提取首表面深度，光度损失保持不变。TSGS 系统分析了透明表面重建中"透明-深度困境"的根源。标准 3DGS 通过 $\alpha$-混合 渲染深度 $D_{\mathrm{standard}} = \sum_{i} T_i \alpha_i z_i$，其权重 $\alpha_i T_i$ 本质上是为外观渲染优化的——透明表面为了正确混合前后景的光线传输，需要较低的不透明度，而精确的首表面深度提取则需要不透明度在真实首表面处急剧集中。这种需求矛盾使得标准深度渲染对透明物体产生严重偏差。此前 PGSR 的无偏深度 $D_{\mathrm{unbiased}}(u) = D(u) N_w(u) \cdot v_u$ 虽改进了标准混合深度，但其 $D$ 和 $N_w$ 仍通过 $\alpha$-混合 聚合所有相交高斯的贡献，无法隔离首表面。TSGS 提出三阶段策略来解耦几何与外观。
\textbf{方法信息：} 代码：\url{https://github.com/longxiang-ai/tsgs}；GitHub 关注度较高，约 122 stars。



% METHOD: 2505.05591
\paragraph{QuickSplat（学习式初始化与密集化）（2025.05）}
QuickSplat~\cite{2505_05591} 针对大规模室内表面重建的训练速度。原始 2DGS 依赖 SfM 点云初始化和手工 densification；在弱纹理、遮挡或未充分观测区域，它往往需要较长优化才能补齐表面。QuickSplat 的改变是把“初始高斯放在哪里、训练中该如何新增高斯”改成可学习模块。

方法从 SfM 点云出发，先用稀疏 3D CNN 在体素空间提取场景特征，再用小 MLP 解码 2DGS 的位置、尺度、法线、opacity 等初始参数。为了补足 SfM 点云没有覆盖的区域，它还预测额外高斯作为密集初始化。训练过程中，更新网络预测高斯参数的修正向量，密集化网络根据渲染残差和梯度特征决定是否生成新高斯，替代原始 2DGS 中固定阈值的 clone/split 规则。最终仍用收敛 2DGS 渲染深度并做 TSDF 融合提取网格。相对 DN-Splatter 或 CDGS 的外部几何监督，QuickSplat 主要学习初始化和 densification 策略本身。
\textbf{实现细节：} DN-Splatter 使用标准 $\alpha$-混合 渲染深度和法线，但"期望深度"本身就有问题。RaDe-GS（2024.06） 修改了原始 3DGS 的深度/法线渲染方式，用精确的封闭形式公式替代期望深度的近似计算，训练流程保持不变。其核心贡献在于提出了精确的高斯深度和法线光栅化方法，揭示了标准 3DGS 深度渲染的固有缺陷并给出封闭形式的修正方案。问题：期望深度忽略了高斯的体积分布。标准 3DGS 只渲染颜色和不透明度，深度图通常通过 $\alpha$-混合 的期望深度 $\hat{z} = \sum_i \omega_i z_i$ 近似计算，其中 $\omega_i = \alpha_i \prod_{j=1}^{i-1}(1-\alpha_j)$ 是混合权重，$z_i$ 是第 $i$ 个高斯中心的深度。RaDe-GS 指出这种做法将每个高斯基元视为一个"点"：只用其中心深度代表整个基元，完全忽略了高斯在三维空间中的体积分布。实际上，不同像素的视线射线与同一个高斯基元相交于不同深度，因此深度应该是像素位置的函数，而非常数。精确的深度光栅化。简单来说，RaDe-GS 发现射线穿过高斯时，不同像素位置看到的"表面"深度其实不同——就像一个椭球气球，从不同位置看气球表面，深度各不相同。RaDe-GS 的关键洞察是：在 3DGS 采用的局部仿射投影模型下，射线与高斯基元的"交点"（定义为使高斯值最大的点）的深度可以封闭形式地计算。具体推导在"光线空间"（ray space）中完成。3DGS 将每个高斯从世界坐标系通过仿射变换投影到屏幕空间，得到 2D 高斯协方差 $\boldsymbol{\Sigma}' = \mathbf{J}\mathbf{W}\boldsymbol{\Sigma}\mathbf{W}^{}\top\mathbf{J}^{}\top$，其中 $\mathbf{W}$ 是世界到相机的旋转矩阵，$\mathbf{J}$ 是透视变换的雅可比矩阵（Jacobian Matrix）。RaDe-GS 进一步在光线空间中考虑第三个坐标维度 $t$（点到相机的距离），定义一个 $1\times 3$ 常数向量： 其中 $\mathbf{v}'=(0,0,1)^{}\top$ 是光线空间中的射线方向常数。对每个高斯，$\hat{\mathbf{q}}$ 可预计算。
\textbf{方法信息：} 代码：\url{https://github.com/liu115/QuickSplat}。
% METHOD: 2509.18090
\paragraph{GeoSVR（稀疏体素表面正则）（2025.09）}
GeoSVR~\cite{2509_18090} 虽基于稀疏体素 rasterization，但它与 3DGS 几何路线密切相关：二者都追求快速显式表示下的表面准确性。稀疏体素不像 3DGS 那样天然来自 SfM 点云初始化，早期几何更容易扭曲；GeoSVR 的改变是给稀疏体素加入不确定性感知深度约束和局部表面一致性正则。

它以 SVRaster 的稀疏体素八叉树为基础。Voxel Uncertainty Depth Constraint 使用 Depth Anything 等单目深度作为场景约束，但不把深度当成全局可靠标签，而是为体素学习或估计不确定性：深度先验与渲染几何一致的区域给强监督，不一致或边缘区域降低权重，避免错误单目深度把体素拉歪。Sparse Voxel Surface Regularization 则在相邻微小体素之间施加几何一致性，使局部体素更容易形成锐利连续表面。八叉树自适应细分负责覆盖复杂区域。相对 DN-Splatter/CDGS 在高斯上做深度正则，GeoSVR 是稀疏体素分支上的对应方案，也为后续 structured priors 方法提供了对照。
\textbf{实现细节：} PAGaS 通过极端的参数约束来保障几何质量，但这种方法只适用于深度图精炼而非全局优化。GeoSVR（2025.09） 回到正则化路线，但换了一个底层表示：建立在稀疏体素光栅化（Sparse Voxel Rasterization, SVRaster）之上。SVRaster 用八叉树组织稀疏体素，每个体素在八个角点存储密度值，通过三线性插值建模体内部的连续密度场。与 3DGS 基于高斯泼溅的表示不同，SVRaster 的体素具有极强的局部性：单个体素只通过少数角点与邻居连接，缺乏全局结构约束，优化过程容易产生严重的几何畸变。此外，之前基于 SDF 或 3DGS 的方法能从结构化的初始化（如零水平集或 SfM 点云）中受益，而 SVRaster 采用常量初始化，几乎没有任何先验约束，几何收敛面临更大挑战。GeoSVR 针对这两个问题，提出了基于体素层级的几何不确定性建模。直觉是：粗糙体素（低分辨率）覆盖大范围空间但精度低，应更多依赖外部深度先验；精细体素（高分辨率）已有足够细节，应更多信任自身优化结果。
\textbf{设计动机与边界：} 相对原始 3DGS 以独立高斯和图像重建误差为主的设定，GeoSVR 把问题落在以下可实现的改动上：GeoSVR基于SVRaster（稀疏体素rasterization）框架。核心解决两个挑战：(1) 体素不确定性深度约束（VUDC）：针对稀疏体素缺乏点云初始化等结构先验导致的几何扭曲问题，利用单目深度先验（如Depth Anything）提供场景约束，同时引入体素不确定性建模避免深度先验误差导致的几何退化。体素级不确定性允许在可靠区域强约束、不确定区域弱约束。(2) 稀疏体素表面正则化（SVSR）：针对单个体素影响范围局部导致表面形成困难的问题，设计正则化增强微小体素间的几何一致性，促进形成锐利准确的表面。基于八叉树的自适应体素细分策略保持覆盖完整性。其设计目的不是单独增加一个模块，而是用上述表示、监督或渲染约束消除原流程在目标任务中的具体失配。提出GeoSVR，探索稀疏体素在表面重建中的潜力。设计体素不确定性深度约束（VUDC）最大化单目深度先验效果同时引入体素不确定性避免质量退化；提出稀疏体素表面正则化（SVSR）增强微小体素的几何一致性并促进锐利表面形成，在几何精度、细节保持和重建完整性上全面超越现有方法。适用时还应保留论文给出的边界：八叉树细分的最大层级可能限制极细细节的表达；体素化分辨率与计算成本的权衡。
\textbf{方法信息：} 代码：\url{https://github.com/fictionarry/geosvr}；GitHub 关注度较高，约 190 stars。



% METHOD: 2510.11473
\paragraph{VA-GS（视图对齐几何增强）（2025.10）}
VA-GS~\cite{2510_11473} 不改 3DGS 表示，而是把几何改进集中到单视图和多视图对齐损失上。原始 3DGS 的颜色损失能得到好看的图像，但边界、遮挡和多视角几何关系不一定稳定；VA-GS 的目标是让每个视图的几何线索与跨视图关系同时参与优化。

单视图部分加入边缘感知图像重建损失，让高斯在物体轮廓和纹理边缘处保留更尖锐结构；法线一致性也按边缘加权，避免边界附近被平滑法线过度拉扯。多视图部分使用可见性感知光度对齐：只在跨视图可见关系可靠的位置比较颜色和结构，遮挡或不可见区域不强行对齐。同时，深度特征嵌入提供跨视图一致的几何描述，使不同视角下的深度和法线更容易闭合。相对 PGSR 的无偏深度公式，VA-GS 更偏“损失组织”路线：它不重新定义深度渲染，而是让可见性、边缘和跨视图特征共同约束高斯位置。
\textbf{实现细节：} MVN-DG 关注多视角几何一致性，但未针对光照变化做特殊处理。VA-GS（2025.10） 在原始 3DGS 的训练环节引入视图对齐策略，同时添加边缘感知梯度监督和多视角光度/特征对齐约束，光度损失保持不变。为此，VA-GS 在单视角和多视角两个维度分别设计约束。多视角方面包含两项损失：光度对齐损失利用单应矩阵将参考视角的局部面片 $P_r$（由法线 $n_r$ 和距离 $d_r$ 定义）扭曲到源视角得到 $P_s$，通过归一化互相关 $C(\cdot)$ 衡量面片间的一致性，并引入遮挡感知权重 $\psi(p_r) = 1/\exp(\epsilon(p_r))$（$\epsilon$ 为重投影误差）剔除被遮挡像素的干扰： 特征对齐损失用预训练网络提取的特征图代替原始 RGB 图，以余弦相似度度量跨视角特征一致性，增强了光照变化和低纹理场景下的鲁棒性。最终损失为五项的加权和 $\mathcal{L} = \mathcal{L}_I + \lambda_1 \mathcal{L}_{nc} + \lambda_2 \mathcal{L}_{ns} + \lambda_3 \mathcal{L}_p + \lambda_4 \mathcal{L}_f$。
\textbf{方法信息：} 代码：\url{https://github.com/leoqli/va-gs}。
\textbf{设计动机与边界：} 相对原始 3DGS 以独立高斯和图像重建误差为主的设定，VA-GS 把问题落在以下可实现的改动上：VA-GS从单视图和多视图两个层面增强几何表示。单视图对齐包括：(1)边缘感知图像重建损失：在标准L1+D-SSIM损失基础上加入图像梯度监督，鼓励模型保留尖锐结构和边界细节；(2)基于法线的几何对齐：使用边缘感知法线一致性损失，在边界区域利用图像梯度作为几何边缘代理，通过下加权边界区域的损失贡献避免模糊法线方向导致的几何不准。多视图对齐包括：(1)可见性感知光度对齐损失：建模遮挡关系，鼓励高斯基元间准确的空间关系；(2)深度特征嵌入：利用深度特征（如DINO特征）实现跨视图一致性约束，增强几何在不同视角和光照条件下的鲁棒性。整体框架在不修改3DGS基本表示的前提下，通过精心设计的损失函数组合提升几何质量。其设计目的不是单独增加一个模块，而是用上述表示、监督或渲染约束消除原流程在目标任务中的具体失配。提出通过视图对齐（VA）增强3DGS几何表示的新方法。引入边缘感知图像线索改进表面边界描绘、可见性感知光度对齐损失建模遮挡并鼓励准确空间关系、基于法线的约束优化空间取向、以及深度图像特征嵌入实现跨视图一致性。在标准基准上实现表面重建和新视角合成的双SOTA。适用时还应保留论文给出的边界：多个损失项的超参数调优较复杂；深度特征提取增加了计算开销；对极度缺乏纹理的区域仍可能产生几何伪影。
% METHOD: 2510.16837
\paragraph{2DGS-R（法线一致性的分层训练）（2025.10）}
2DGS-R~\cite{2510_16837} 重新审视 2DGS 的 normal consistency。法线一致性对几何很重要，但会损失渲染质量；2DGS-R 的改变是分阶段训练，让几何先稳定，再有选择地补外观容量。

方法的第一阶段保持 2DGS 的深度失真和法线一致性损失，目标是先得到可用表面。论文分析指出，normal consistency 会同时改变 2D 高斯的覆盖面积和 opacity：覆盖更大、opacity 更高通常有利于 mesh，但会压低 PSNR。第二阶段只对渲染残差高或细节不足的区域做 in-place cloning，也就是在原位置复制 2D 高斯并给它额外外观容量，而不是像常规 densification 那样大范围分裂。第三阶段冻结所有高斯的 opacity，只微调位置、尺度、颜色等属性，让渲染质量回升但不把前一阶段稳定下来的几何重新打散。相对 PGSR 或 VCR-GauS 调整几何监督，2DGS-R 关注的是几何与外观的训练顺序：先牺牲一点外观换稳定几何，再局部补外观容量。
\textbf{实现细节：} 单独的深度或法线正则化各有盲点：深度约束位置但不约束朝向，法线约束朝向但梯度无法传到位置。联合正则化旨在打通这两个维度。VCR-GauS（2024.06） 在原始 3DGS 的训练环节引入"交点深度"替代高斯中心深度，使法线监督的梯度能够传导到高斯位置参数，光度损失保持不变。动机。法线仅与高斯旋转参数相关，法线监督无法有效更新高斯位置——即"知道表面朝向不对"却"无法移动高斯到正确位置"，需要一座桥梁打通法线与位置的梯度通路。所谓"交点深度"，是指将扁平高斯视为平面、求光线与平面的交点深度。它同时依赖于高斯的法线（决定平面朝向）和位置（决定平面在哪里），因此法线监督的梯度可以"顺道"推动位置更新。VCR-GauS 指出 3DGS 中仅监督法线不足以驱动表面重建，原因是法线仅与高斯的旋转参数相关，法线监督无法有效更新高斯的位置。NeuS-based 方法能将法线正则化传导至位置更新（法线即 SDF 的梯度），但 3DGS 缺乏这种机制。为此，VCR-GauS 提出交点深度（intersection depth）替代高斯中心深度：先将高斯挤压为扁平形状（通过尺度正则化 $\mathcal{L}_s = \|(s_1, s_2, s_3)\|_1$ 使最小尺度分量趋零），再将扁平高斯视为平面 $(p, n)$，求光线与平面的交点深度 $d(n, p) = r_z \cdot (n \cdot p) / (n \cdot r)$，其中 $r$ 为光线方向、$r_z$ 为其 $z$ 分量。交点深度同时依赖于法线 $n$ 和位置 $p$，从而打通了法线监督到位置更新的梯度通路。简言之，以前法线梯度只能更新高斯的旋转，现在通过"交点深度"这个桥梁，法线梯度也能推动高斯移动到正确位置。
\textbf{方法信息：} Intel 参与了这篇论文；当前未确认完整公开代码。
% METHOD: 2601.17720
\paragraph{Structured Priors for Sparse-Voxel Reconstruction（结构化稀疏体素先验）（2026.01）}
Advancing Structured Priors for Sparse-Voxel Surface Reconstruction~\cite{2601_17720} 继续稀疏体素几何路线，并借鉴 3DGS 的快速收敛与视觉基础模型先验。方法用 VGGT 伪深度和置信图生成逐图像稀疏体素八叉树，再通过拓扑对齐融合成统一结构。

虽然它不是纯 3DGS 表示，而是基于 SVRaster 的 sparse voxel octree，但与本章关系在于：它把 3DGS 里“快速显式优化、用几何先验约束表面”的思路搬到稀疏体素上。第一步是 LoD unprojection：对每个输入图像，用 VGGT 估计伪深度、置信图和相机，再把像素反投影成体素；像素在三维中覆盖面积越大，就分配越粗的 octree level，细节区域则保留更细体素。第二步是 topology-aligned fusion：不同视图生成的八叉树分辨率不同，方法先建立共同拓扑，再聚合颜色、opacity 和置信度，形成全局一致 SVO。第三步是 refined depth geometry supervision：训练时不只依赖重投影光度误差，而是把多视角深度候选变成逐射线的直接深度正则，增强深度一致性且尽量不把边缘抹平。相对 GeoSVR，它更强调初始化阶段就把体素放在合理位置和合理尺度上，而不是从均匀网格慢慢删改。
\textbf{实现细节：} GeoSVR 从正则化角度改善几何，而初始化质量同样关键。Advancing Structured Priors（2026.01） 修改了原始 3DGS 的初始化环节，用基础视觉模型生成的 LoD 体素替代 SfM 稀疏点云作为高斯起点，训练流程保持不变。其核心改进是用层次细节（Level-of-Detail, LoD）体素初始化替代标准 3DGS 的 SfM 点云初始化。标准 3DGS 依赖 COLMAP 等工具从 SfM 恢复的稀疏点云初始化高斯位置和尺度，但 SfM 点云分布不均匀——纹理丰富的区域点密，弱纹理区域点稀甚至缺失——且含有噪声和离群点。Advancing Structured Priors 改用基础视觉模型（Foundation Vision Model, VGGT）为每个视角估计伪深度图和置信度图，再通过 LoD 反投影算法生成多分辨率体素。对每个像素 $p=(u,v)$，利用伪深度 $\tilde{D}(u,v)$ 将其反投影到三维空间，并通过采样像素中心和四个半偏移邻域计算其三维足迹面积 $A_{i,p}$，据此确定体素层级： 其中 $S$ 是场景尺寸，$s_l$ 是层级 $l$ 的体素边长。三维足迹越大的像素对应越粗糙的体素（低 $l$ 值），足迹越小的像素获得越精细的体素（高 $l$ 值），实现分辨率与空间覆盖的自适应匹配。生成每视图八叉树后，通过拓扑对齐融合算法将 $N$ 个八叉树合并为全局八叉树：对同一空间区域，若某视图已细分到更精细层级，其他视图的粗糙体素也递归细分为八个子体素，属性从父节点继承。融合后的体素属性（球谐系数）通过对所有贡献视图取均值聚合。不透明度低于阈值 $\sigma_{\mathrm{prune}}$ 的体素被剪枝移除。在优化阶段，Advancing Structured Priors 进一步提出精炼深度几何监督：对每个像素在渲染深度附近生成候选深度集（空间邻域 + 随机扰动），通过多视角 NCC 评分选择最优深度，以 $L_1$ 损失直接约束几何，权重设为 $\lambda_{\mathrm{photo}} = 0.1$、$\lambda_{\mathrm{reg}} = 1.0$、$\lambda_{\mathrm{refD}} = 0.05$。
\textbf{技术实现补充：} GeoSVR 从正则化角度改善几何，而初始化质量同样关键。Advancing Structured Priors（2026.01） 修改了原始 3DGS 的初始化环节，用基础视觉模型生成的 LoD 体素替代 SfM 稀疏点云作为高斯起点，训练流程保持不变。其核心改进是用层次细节（Level-of-Detail, LoD）体素初始化替代标准 3DGS 的 SfM 点云初始化。标准 3DGS 依赖 COLMAP 等工具从 SfM 恢复的稀疏点云初始化高斯位置和尺度，但 SfM 点云分布不均匀——纹理丰富的区域点密，弱纹理区域点稀甚至缺失——且含有噪声和离群点。Advancing Structured Priors 改用基础视觉模型（Foundation Vision Model, VGGT）为每个视角估计伪深度图和置信度图，再通过 LoD 反投影算法生成多分辨率体素。对每个像素 $p=(u,v)$，利用伪深度 $\tilde{D}(u,v)$ 将其反投影到三维空间，并通过采样像素中心和四个半偏移邻域计算其三维足迹面积 $A_{i,p}$，据此确定体素层级： 其中 $S$ 是场景尺寸，$s_l$ 是层级 $l$ 的体素边长。三维足迹越大的像素对应越粗糙的体素（低 $l$ 值），足迹越小的像素获得越精细的体素（高 $l$ 值），实现分辨率与空间覆盖的自适应匹配。生成每视图八叉树后，通过拓扑对齐融合算法将 $N$ 个八叉树合并为全局八叉树：对同一空间区域，若某视图已细分到更精细层级，其他视图的粗糙体素也递归细分为八个子体素，属性从父节点继承。融合后的体素属性（球谐系数）通过对所有贡献视图取均值聚合。不透明度低于阈值 $\sigma_{\mathrm{prune}}$ 的体素被剪枝移除。在优化阶段，Advancing Structured Priors 进一步提出精炼深度几何监督：对每个像素在渲染深度附近生成候选深度集（空间邻域 + 随机扰动），通过多视角 NCC 评分选择最优深度，以 $L_1$ 损失直接约束几何，权重设为 $\lambda_{\mathrm{photo}} = 0.1$、$\lambda_{\mathrm{reg}} = 1.0$、$\lambda_{\mathrm{refD}} = 0.05$。
\textbf{方法信息：} NVIDIA 参与了这篇论文；当前未确认完整公开代码。
\textbf{设计动机与边界：} 相对原始 3DGS 以独立高斯和图像重建误差为主的设定，Advancing Structured Priors for Sparse-Voxel Surface Reconstruction 把问题落在以下可实现的改动上：方法基于SVRaster（稀疏体素八叉树SVO）而非纯3DGS，但利用3DGS的快速收敛和几何先验：(1)LoD体素初始化：利用预训练3D视觉基础模型(VGGT)获取伪深度和置信图，通过细节层次反投影算法生成逐图像SVO，使用拓扑对齐融合算法合并为统一八叉树，不确定性感知不透明度场分配初始体素不透明度；(2)精细深度几何监督(RDGS)：将多视角深度线索转化为直接逐射线深度正则化，增强深度一致性且不模糊边缘。其设计目的不是单独增加一个模块，而是用上述表示、监督或渲染约束消除原流程在目标任务中的具体失配。结合3DGS和稀疏体素光栅化的互补优势：提出基于细节层次(LoD)的体素初始化方法，将体素放置在合理位置并设置适当细节层级；提出精细深度几何监督将多视角线索转化为逐射线深度正则化。在标准基准上实现几何精度、细结构恢复和更完整表面的提升。适用时还应保留论文给出的边界：依赖3D基础模型的伪深度质量；体素八叉树的分辨率受内存限制；对动态场景的扩展性未验证。

\section{网格提取与优化}

网格路线的目标是把高斯变成传统图形管线可编辑、可碰撞、可绑定的表面。SuGaR 是关键锚点：先让高斯对齐表面，再提取网格。后续方法有的从高斯采样点云，有的用 opacity field 或 TSDF，有的直接优化三角形/网格，让输出不再是松散高斯集合。
% METHOD: 2311.12775
\paragraph{SuGaR（表面对齐高斯与网格重建）（2023.11）}
SuGaR~\cite{2311_12775} 是本章最重要的网格锚点之一。它指出若高斯要提网格，就必须让高斯扁平、均匀并贴合场景表面，而不是任由光度损失把高斯堆在体积中。原始 3DGS 的高斯可能漂浮、重叠、尺度很小，直接把它们相加成密度场再 Marching Cubes，通常只能得到破碎或近乎空的 level set。

SuGaR 的流程分三段。第一段先按普通 3DGS 训练，让高斯获得可渲染的位置、尺度、颜色和 opacity。第二段加入表面对齐正则：方法在“高斯足够扁、相邻高斯不过度重叠、opacity 更接近二值”的假设下，从高斯密度推导近似距离场，并最小化真实密度与理想表面对齐密度之间的差异；这样高斯会被压到场景表面附近，并沿表面更均匀分布。第三段不在稀疏密度场上直接 Marching Cubes，而是利用训练视角深度图在可见 level set 附近采样点，再用 Poisson reconstruction 生成三角网格。可选 refinement 会把高斯绑定到 mesh 表面，联合优化 mesh 和表面高斯：mesh 提供可编辑、可绑定、可雕刻的几何骨架，高斯继续承担高频纹理和新视角渲染。SuGaR 的意义在于建立了“3DGS 先渲染优化，再表面对齐、mesh 化和绑定”的基本路线。
\textbf{实现细节：} SuGaR（2023.11） 修改了原始 3DGS 的高斯优化环节，用表面对齐正则化替代自由分布的高斯基元，渲染管线保持不变。SuGaR 是从 3DGS 中提取 mesh 的里程碑工作。动机。标准 3DGS 优化后高斯基元散布在空间中，彼此大量重叠，没有明确的表面概念——直接从这样的"概率云"中提取网格几乎不可能，需要先将高斯拉到表面上。SuGaR 的核心思想是：通过正则化使高斯基元对齐到场景表面，然后用 Poisson 表面重建（一种从带法线的点集推断封闭三角网格的经典算法）从对齐的高斯中提取 mesh。SuGaR 首先定义了一个从高斯基元推导的密度函数。直觉上，如果将所有高斯基元视为一团团三维的"概率云"，那么空间中任意一点的密度就是所有云团在该点的叠加——表面附近密度高、远离表面密度低。对空间中任意点 $\mathbf{p}$，将所有高斯加权求和得到密度值： 其中 $\boldsymbol{\mu}_g$、$\boldsymbol{\Sigma}_g$、$\alpha_g$ 分别为第 $g$ 个高斯的中心、协方差和混合权重。如果高斯已经很好地对齐到表面且充分扁平，那么在每个空间位置附近，贡献最大的高斯 $g^{*}$（马氏距离最小者）的主导作用远超其他高斯，密度函数可近似为一个理想形式。基于此观察，SuGaR 构造了一个理想符号距离函数（SDF）。直觉上，SDF 就像一个"距离指示器"——表面上的点距离为零，表面外侧为正、内侧为负。设 $s_{g^{*}}$ 是 $g^{*}$ 的最小缩放因子，则理想 SDF 定义为： 这个公式的含义是：当高斯完美对齐时，$f(\mathbf{p})=0$ 的零交叉面恰好对应场景表面——离表面越远，$f$ 值越大；在表面上，$f$ 恰好为零。SuGaR 的正则化目标就是让实际 SDF 估计 $\hat{f}$ 逼近理想 SDF $f$。
\textbf{锚点细化：} SuGaR：针对 3DGS 优化后高斯基元散布空间无明确表面概念的局限，用"表面对齐正则化 + Poisson 网格提取"替代直接网格提取，核心思想是定义从高斯加权求和的密度函数 $d(\mathbf{p})$，构造理想 SDF $f(\mathbf{p}) = s_{g^{*}} - \sqrt{2\ln d(\mathbf{p})}$，通过 align 与 normal 两项损失迫使高斯贴合表面，再从训练视角深度图采样搜索密度等值面点执行 Poisson 重建（一种从带法线的散乱点集推断封闭水密三角网格的经典算法），是从 3DGS 提取 mesh 的里程碑工作。
\textbf{方法信息：} 代码：\url{https://github.com/Anttwo/SuGaR}；GitHub 关注度很高，约 3447 stars；引用很多，Semantic Scholar 约 836 次。



% METHOD: 2404.01810
\paragraph{GS2Mesh（渲染立体对转网格）（2024.04）}
GS2Mesh~\cite{2404_01810} 不直接从高斯参数提几何，而是利用 3DGS 能渲染任意视角这个能力。它人为从训练相机生成立体对：左图为原视角，右图为水平基线偏移的新视角，然后用预训练 stereo 网络估计深度。

流程可以理解为“先把 3DGS 变成可控相机，再把 stereo 变成几何提取器”。方法先用 COLMAP 位姿训练普通 3DGS，使其能够在训练轨迹附近稳定渲染；随后为每个训练视角合成一个基线很小、极线对齐的新视角，与原视角组成 stereo pair；预训练立体匹配网络从这些成对图像预测深度，再结合遮挡检查、深度一致性过滤和可选 SAM 掩码去掉不可靠区域。最后，所有视角的深度被送入 TSDF fusion 得到 mesh。相对 SuGaR/GOF 从高斯密度或 opacity 场中直接抽 surface，GS2Mesh 把几何判断交给通用 stereo 先验；相对 Self-Evolving DS-GS，它的 stereo 主要服务于最终 mesh 转换，而不是在 3DGS 训练中反复生成监督。
\textbf{实现细节：} GS2Mesh（2024.04） 的核心思想是：不直接从高斯基元提取几何，而是将 3DGS 渲染的立体图像对输入预训练立体匹配网络，借其大规模数据训练的强几何先验来重建 mesh——即用"立体视觉先验"替代"高斯基元直接重建"作为几何推断路径，从而绕开 3DGS 体积元素带来的噪声和伪影。GS2Mesh 利用预训练立体匹配模型从 3DGS 渲染的立体图像对中提取高质量 mesh。GS2Mesh 的出发点是一个关键观察：标准 3DGS 本身的几何质量有限——高斯基元是体积元素而非表面元素，从中直接提取 mesh 往往充满噪声和伪影。但如果用 3DGS 渲染出的高质量图像（而非高斯基元本身）作为输入，再通过成熟的立体视觉管线重建几何，就能借助立体匹配网络在大规模数据上学习的强几何先验来弥补 3DGS 的几何缺陷。GS2Mesh 的流程分为三个阶段。第一阶段是立体对渲染：对每个训练视角位姿 $(R_L, T_L)$，将左相机设为训练位姿本身（确保渲染质量最高），右相机通过水平基线 $b$ 偏移构造： 这样产生的立体对天然满足立体标定条件。基线 $b$ 的设置需要在两个矛盾因素间权衡：基线越大，立体匹配的深度精度越高（"甜点区"越宽）；但基线越大，3DGS 新视角渲染的质量越差（偏离训练视角越远，渲染噪声越大）。第二阶段是立体深度估计：GS2Mesh 用预训练的立体匹配网络（实验中 DLNR + Middlebury 预训练权重效果最佳）估计每对立体图的视差图。深度误差与视差误差的关系为 $\sigma(Z) \propto \sigma(d) \cdot Z^{}2 / (f_x \cdot B)$，即深度误差与距离的平方成正比——远处的深度估计天然不如近处精确。为滤除不可靠估计，GS2Mesh 应用两类掩码：遮挡掩码（左右视差一致性检查，剔除仅单侧可见的区域）和深度范围掩码（仅保留 $4b \leq Z \leq 20b$ 的有效范围）。被遮挡区域由相邻视角的估计填补。
\textbf{锚点细化：} GS2Mesh：针对高斯基元本身几何质量有限直接提取 mesh 噪声大的局限，用"立体匹配网络从渲染图像提取几何"替代直接从高斯提取，核心思想是用 3DGS 渲染立体对（左相机为训练位姿、右相机水平基线偏移），用预训练立体匹配网络估计视差图，深度误差与距离平方成正比，经遮挡掩码与深度范围掩码过滤后通过 TSDF 融合（截断符号距离函数融合，将多视角深度图加权平均为全局一致体素网格的经典去噪步骤）提取网格——3DGS 只负责"画图"，几何重建完全交给立体视觉。
\textbf{方法信息：} 代码：\url{https://github.com/yanivw12/gs2mesh}；GitHub 关注度很高，约 353 stars；引用很多，Semantic Scholar 约 83 次。



% METHOD: 2404.10772
\paragraph{Gaussian Opacity Fields（连续不透明度场）（2024.04）}
Gaussian Opacity Fields~\cite{2404_10772} 直接把 3DGS 扩展成可查询连续不透明度场。普通 3DGS 是离散高斯集合，网格提取常依赖 Poisson 或 TSDF；GOF 的改变是从 ray-Gaussian intersection 出发，为任意 3D 点定义 view-independent opacity field，再从其 level set 提取几何。

给定一条射线，GOF 在高斯局部坐标中求该射线上高斯响应最大的交点，并用这个交点定义该高斯沿射线的不透明度贡献。对任意三维点，方法不只看一个视角，而是在所有可见训练视图中计算 opacity，并取跨视图稳定的估计，使 opacity 变成主要由空间位置决定的连续场。法线由射线-高斯交点平面近似，因而可以接入深度失真和法线一致性正则。网格提取使用 Marching Tetrahedra，四面体网格由高斯中心和缩放自适应诱导，而非均匀网格。相对 3DGS-to-PC 的后处理采样，GOF 把几何场定义进高斯模型本身，也减少了 Poisson/TSDF 与训练目标不一致的问题。
\textbf{实现细节：} Gaussian Opacity Fields (GOF)（2024.04） 修改了原始 3DGS 的渲染环节，用显式的光线-高斯交点计算取代二维投影，从而保留完整的三维不透明度信息，使其可直接用于表面提取。GOF 是另一项里程碑工作。动机。标准 3DGS 将三维高斯投影到二维屏幕后即丢失三维信息，无法从投影结果反推空间中任意一点的不透明度——因此无法定义连续的表面场来提取网格。标准 3DGS 的渲染管线将三维高斯投影到二维屏幕空间后计算像素颜色，一旦投影完成，三维信息就丢失了，无法从投影结果反推空间中任意一点的不透明度。具体而言，给定相机中心 $\mathbf{o}$ 和射线方向 $\mathbf{d}$，射线上的点为 $\mathbf{x} = \mathbf{o} + t\mathbf{d}$。直觉上，要计算射线与高斯的交点，最自然的思路是把射线变换到高斯的局部坐标系下——在那里高斯退化为一个标准椭球，射线变成一条直线，交点计算大大简化。GOF 先将射线变换到每个高斯 $k$ 的归一化局部坐标系下： 在此归一化坐标系中，高斯值沿射线退化为一个一维高斯函数 $G_k^{\mathrm{1D}}(t) = \exp\!\bigl(-\frac{1}{2}(At^{}2 + 2Bt + C)\bigr)$，其中 $A = \mathbf{d}_g^{}\top\mathbf{d}_g$，$B = \mathbf{p}_g^{}\top\mathbf{d}_g$。对 $t$ 求极值可得交点深度 $t^{*} = -B/A$，该交点处的高斯值即为高斯 $k$ 对这条射线的最大贡献。GOF 的关键创新在于定义了逐点不透明度。直觉上，从相机出发沿射线行进时，每个高斯像一个"半透明的团雾"——越靠近高斯中心，雾越浓；穿过后，浓度不再增加（因为你已经穿过了最浓处）。在此基础上，多个高斯沿射线累积不透明度的方式与标准体渲染完全一致： 由于一个三维点可被多个训练视角看到，GOF 取所有视角中的最小值定义该点的全局不透明度 $O(\mathbf{x}) = \min_{(\mathbf{o},\mathbf{d})} \tau(\mathbf{o}, \mathbf{d}, t)$，这就是所谓的高斯不透明度场。
\textbf{锚点细化：} GOF：针对 3DGS 投影到 2D 后丢失三维信息无法定义连续表面场的局限，用"光线-高斯交点的三维不透明度场"替代 2D 投影，核心思想是在每个高斯归一化局部坐标系中求射线-高斯交点，定义逐点不透明度 $\tau_k(t)$——前向递增到交点最大值后保持不变（一个高斯对前方点贡献递增、对后方点贡献恒定），取多视角最小值定义全局不透明度场 $O(\mathbf{x})$，用 Marching Tetrahedra（用四面体而非立方体单元划分空间以提取等值面，支持非均匀分辨率）提取自适应分辨率 mesh。
\textbf{方法信息：} 代码：\url{https://github.com/autonomousvision/gaussian-opacity-fields}；GitHub 关注度很高，约 1011 stars；引用很多，Semantic Scholar 约 304 次。



% METHOD: 2501.07478
\paragraph{3DGS-to-PC（概率采样点云/网格转换）（2025.01）}
3DGS-to-PC~\cite{2501_07478} 是实用转换工具路线。SuGaR、GOF 这类方法会在训练或提取阶段重新定义几何约束；3DGS-to-PC 不改已有 3DGS，而是把一个训练好的 \texttt{.ply} 或 \texttt{.splat} 场景转换成传统点云/网格，方便进入 Meshlab、Open3D 或标准三维处理流程。

它把每个 3D 高斯视为概率密度函数，按高斯尺度分配采样数量，在椭球内部随机采样点；采样点若超过马氏距离阈值就丢弃，避免大协方差高斯产生离中心过远的离群点。颜色也不能直接读取 SH 系数，因为 SH 颜色依赖观察方向。方法会回到训练相机视图，计算每个高斯对渲染像素的贡献，把贡献最大的可见颜色作为采样点颜色。这样输出的点云既保留高斯空间分布，又有稳定的 RGB；用户还可以控制目标点数和采样尺度，在密度与文件大小之间取舍。最后可选用泊松表面重建把点云转成彩色网格。相对 SuGaR，它不学习新的表面对齐表示，更像把已有 3DGS 变成可交换、可后处理的几何资产。
\textbf{实现细节：} 3DGS-to-PC（2025.01） 提出将 3DGS 场景转换为密集点云或 mesh 的通用框架。3DGS-to-PC 的核心问题是：直接从高斯分布中采样三维点是容易的（高斯本身就是概率分布），但如何确保采样出的点云颜色与实际渲染结果一致却非常困难——每个高斯的颜色只是球谐系数的零阶分量，而最终像素颜色是多个高斯 $\alpha$-混合 的结果，单个高斯的颜色不能代表它所在位置的视觉外观。3DGS-to-PC 的流程分为三个阶段。第一阶段是颜色校正：直觉上，一个像素的最终颜色是多个高斯叠加的结果，每个高斯的"贡献"取决于它的 alpha 值和累积透射率。3DGS-to-PC 首先用原始相机位姿渲染场景，对每个像素计算每个高斯的贡献权重： 其中 $\alpha_i$ 是高斯 $i$ 在该像素的 alpha 值，$T_i$ 是累积透射率。对每个高斯，找到它贡献最大的那个像素（即 $C_i$ 最大的像素），将该像素的渲染颜色赋给该高斯作为校正后颜色。贡献为零的高斯（从未被任何视角渲染到）被直接移除。这样，每个高斯的颜色就不再是孤立的球谐零阶分量，而是反映了它实际参与渲染的视觉外观。第二阶段是概率采样：3DGS-to-PC 根据每个高斯的体积分配采样点数量。直觉上，体积越大的高斯覆盖的空间范围越广，需要更多的采样点来充分表达其内部结构。高斯体积的计算公式为： 其中 $s_i$ 是各维度的尺度参数。体积越大的高斯获得越多的采样点（因为它们通常代表大面积表面），用户可控制总点数来调节点云密度。采样点从高斯的多变量正态分布中随机抽取，但为了防止采样点落在高斯椭球外过远的位置，3DGS-to-PC 计算每个采样点的马氏距离： 默认阈值为 $D_M \leq 2$（即 $2\sigma$ 椭球内），超出阈值的点被丢弃并重新采样，最多重试 5 次。马氏距离的关键优势是尺度不变性——无论高斯大小如何，阈值含义一致。第三阶段是可选 mesh 生成：3DGS-to-PC 先用表面高斯筛选过滤掉内部高斯——仅保留最大渲染贡献高于全体高斯均值的高斯作为"表面高斯"，再用统计离群点移除（Open3D）清理噪声，最后用 Poisson 表面重建生成三角网格。法线取自高斯最小缩放轴方向（与 SuGaR 一致）。
\textbf{设计动机与边界：} 相对原始 3DGS 以独立高斯和图像重建误差为主的设定，3DGS-to-PC 把问题落在以下可实现的改动上：3DGS-to-PC框架包含三个核心组件：(1)概率点采样：将每个高斯视为3D密度函数，概率性地采样点，采样数量与高斯大小成比例，确保整个场景均匀覆盖。使用马氏距离阈值过滤极端离群点，防止采样点偏离高斯中心过远。(2)颜色重计算：由于高斯使用球谐函数根据视角调整颜色，直接采样会得到不正确或不稳定的颜色。通过自定义渲染方法，将每个高斯分配其在所有视角中对渲染像素贡献最大的颜色。(3)网格生成：支持通过泊松曲面重建从表面高斯采样点生成彩色网格。框架接受.ply或.splat格式的3DGS数据，以及COLMAP格式的相机位姿，输出标准.ply点云文件。用户可自定义生成点数和采样策略。其设计目的不是单独增加一个模块，而是用上述表示、监督或渲染约束消除原流程在目标任务中的具体失配。提出3DGS-to-PC框架，通过概率采样从每个高斯中生成密集点云，使用马氏距离阈值防止极端离群点，并通过自定义图像渲染方法重新计算高斯颜色。支持泊松曲面重建生成网格，实现3DGS到点云/网格的高精度无损转换。适用时还应保留论文给出的边界：采样质量和数量受原始高斯质量限制；对透明和半透明区域的处理仍有挑战；网格质量取决于泊松重建参数。
\textbf{方法信息：} 代码：\url{https://github.com/lewis-stuart-11/3dgs-to-pc}；GitHub 关注度高，约 789 stars。



% METHOD: 2501.18152
\paragraph{StructuredField（四面体结构场）（2025.01）}
StructuredField~\cite{2501_18152} 与 SuGaR 的方向相反。SuGaR 先训练高斯，再把高斯表面对齐并提取网格；StructuredField 先保留结构化四面体网格，再把四面体重参数化为可渲染高斯，目标是在可编辑拓扑和高质量渲染之间取得平衡。

方法把每个四面体的局部几何映射成 3D 高斯的位置、旋转和尺度，使结构化网格也能进入 Gaussian rasterization 管线。为了表达多尺度细节，它在四面体内部做层次化隐式细分，生成保形子四面体，而不是只依赖原始粗网格。优化时加入无反转同胚约束，要求四面体变形保持正雅可比，避免翻转、自交和局部拓扑破坏。渲染损失通过高斯形式回传，几何约束则维持网格结构。相对“先高斯后网格”的路线，StructuredField 更适合需要连通结构、变形控制或传统图形管线接口的场景。
\textbf{实现细节：} StructuredField（2025.01） 统一结构化四面体网格与高保真可微渲染，解决点云表示（3DGS）缺乏结构、无法直接用于动画/仿真等下游工程应用的根本问题。核心表示是四面体网格：每个顶点携带可学习的高斯基元属性（位置、协方差、球谐系数、不透明度），四面体内部通过层次隐式细分（hierarchical implicit subdivision）实现多尺度细节——粗四面体捕获大尺度结构，细分子四面体补充高频纹理，类似 Instant-NGP 的多分辨率哈希编码但施加在结构化网格上。关键几何保证是方向保持同胚约束（orientation-preserving homeomorphism）：约束顶点的可行变形空间，严格防止自相交和单元翻转（inversion），确保优化过程中四面体网格始终保持数学有效的流形结构。高斯属性通过四面体顶点的重参数化与几何参数绑定，图像损失可直接端到端优化网格顶点位置。最终输出的四面体网格既支持高质量实时渲染，又天然兼容动画、物理仿真和变形等传统图形学管线。StructuredField 将高斯基元嵌入四面体网格。Gaussian Mesh Renderer (GMR) 则反向操作：将已有的 mesh 三角面片转换为高斯基元，使 mesh 可以利用 3DGS 的高效光栅化器进行可微渲染。
\textbf{方法信息：} 代码：\url{https://github.com/kevin2000a/structuredfield}。
\textbf{设计动机与边界：} 相对原始 3DGS 以独立高斯和图像重建误差为主的设定，StructuredField 把问题落在以下可实现的改动上：StructuredField以四面体网格为基础表示，将四面体的几何属性重参数化为3D高斯的参数（位置、旋转、缩放），使四面体网格可直接进行可微高保真渲染。核心创新包括：(1)层次化隐式细分：通过隐式函数在四面体内部进行自适应细分，生成保形的子四面体结构以捕获多尺度细节。(2)无反转同胚约束：提出新的约束函数保证四面体在优化过程中始终保持正值雅可比行列式，确保无反转和无自交。(3)端到端优化：使用图像损失直接优化四面体网格参数，无需后处理即可得到高质量结构化网格。其设计目的不是单独增加一个模块，而是用上述表示、监督或渲染约束消除原流程在目标任务中的具体失配。StructuredField用结构化四面体网格表示重建物体，将四面体的几何属性重参数化为3D高斯参数实现可微高保真渲染。引入层次化隐式细分策略确保保形网格结构，并提出无反转同胚约束保证四面体网格在优化中保持无反转和无自交。实现高质量结构化网格与3DGS可比的重建质量。适用时还应保留论文给出的边界：四面体网格的分辨率受限于初始网格质量；优化过程中网格质量的维护增加计算开销；对初始四面体化有依赖
% METHOD: 2506.06988
\paragraph{Hybrid Mesh-Gaussian（室内混合表示）（2025.06）}
Hybrid Mesh-Gaussian Representation~\cite{2506_06988} 的动机是：室内场景里大量墙面、地板和桌面适合纹理网格，只有复杂细节需要高斯。全部用高斯会数量大、渲染慢；全部用网格又难处理复杂几何。

方法以 PGSR/3DGS 提取的初始 mesh 为起点，先做 mesh pruning。它用高斯尺度、mesh 法线变化、相邻三角面夹角和三角形尺寸判断哪些区域几何复杂或不可靠；这类区域的 mesh 被删掉，交给高斯继续表示。平坦但纹理丰富的区域则保留 mesh，并给它优化 texture map，因为墙面、地板上的高频纹理用贴图通常比堆大量高斯更高效。联合优化时，warm-up 阶段先让高斯稳定；之后引入 transmittance-aware texture loss：若某像素前方高斯透射率低，说明该像素被高斯物体挡住，就降低 mesh 贴图监督，避免把前景颜色错误烘焙到背景 mesh 上。渲染时，纹理网格作为不透明背景，高斯负责复杂几何和前景细节。相对 SuGaR 的单网格输出，它接受 mesh 与 GS 长期共存，目标是减少室内场景中的冗余高斯并提高 FPS。
\textbf{实现细节：}
\textbf{方法信息：} Huawei 参与了这篇论文；当前未确认完整公开代码。
\textbf{设计动机与边界：} 相对原始 3DGS 以独立高斯和图像重建误差为主的设定，Hybrid Mesh-Gaussian Representation for Efficient Indoor Scene Reconstruction 把问题落在以下可实现的改动上：混合表示方法包含三个步骤：(1) 网格修剪与精炼：首先从3DGS提取网格，修剪几何复杂区域的网格（为高斯腾出空间），精炼纹理丰富区域的网格以优化纹理贴图。通过评估每个三角面片的几何复杂度来决定保留或移除。(2) 联合优化：提出warm-up策略和透射感知纹理损失来平衡3DGS和网格的贡献。透射感知损失根据高斯的透射率动态调整纹理监督权重，避免高斯和网格之间的冲突。(3) 渲染管线：纹理网格（不透明度为1）作为背景，高斯表示复杂几何和细节，通过混合渲染实现高效推理。其设计目的不是单独增加一个模块，而是用上述表示、监督或渲染约束消除原流程在目标任务中的具体失配。提出混合Mesh-GS表示用于室内场景：用纹理网格处理纹理丰富的平坦区域，用高斯处理复杂几何区域。通过网格修剪、联合优化和透射感知监督策略，在保持渲染质量的同时显著减少高斯数量并提升FPS。适用时还应保留论文给出的边界：网格提取质量影响整体效果；网格修剪阈值需要手动调整；主要针对室内场景设计，室外场景泛化性未验证。
% METHOD: 2509.25122
\paragraph{Triangle Splatting+（不透明三角形 splatting）（2025.09）}
Triangle Splatting+~\cite{2509_25122} 直接把可微 splatting 图元改成三角形。3DGS 输出不是标准图形基元，Triangle Splatting 的早期版本又像三角形汤；Triangle Splatting+ 通过共享顶点让三角面片有连通性，并在训练结束时收敛为不透明三角形。

它从一个共享顶点集合出发，每个顶点带三维位置、颜色和 opacity，三角形只保存三个顶点索引。投影到图像后，方法用二维三角形 SDF 和平滑窗口函数计算像素覆盖，再用类似 splatting 的方式做可微混合；三角形内部颜色由重心坐标插值。训练早期允许半透明和软边界，以便梯度能稳定进入顶点；后期逐步推动 opacity 到 1，并让三角形边界收敛成硬边界，使最终输出可以被标准图形引擎当作普通不透明三角形渲染。共享顶点让相邻面片可以合并和简化，缓解早期 Triangle Splatting 的孤立三角形汤问题。相对 MeshSplatting，它仍更接近 splatting 优化和半连通三角表示；MeshSplatting 进一步用受限 Delaunay 约束形成更强的 mesh 连通结构。
\textbf{实现细节：} Triangle Splatting+（2025.09） 修改了原始 3DGS 的基元类型，用不透明三角面片替代高斯核，渲染管线的 $\alpha$-混合 框架保持不变。该方法将基元从高斯核替换为不透明三角面片，实现可直接导入游戏引擎的端到端网格优化。其表示层采用顶点共享设计：定义全局顶点集 $V=\{v_i \in \mathbb{R}^{}3\}_{i=1}^{N}$，每个顶点携带位置 $(x_i, y_i, z_i)$、颜色 $c_i$ 和不透明度 $o_i$；三角面片 $T_m=(i,j,k)$ 由三个顶点索引定义，面片不透明度取顶点均值 $o_{T_m}=\mathrm{mean}(o_i, o_j, o_k)$，面片内颜色通过重心坐标插值 $c_{T_m}=\lambda_i c_i + \lambda_j c_j + \lambda_k c_k$。共享顶点使得相邻面片自然连通，梯度从所有相连面片回传到同一顶点。渲染时，3D 顶点经针孔相机模型投影到 2D，利用带符号距离场定义窗口函数 $I(p)$：在三角面片内心处值为 1、边界处为 0，由全局平滑参数 $\sigma$ 控制过渡宽度。像素颜色按深度顺序 $\alpha$-混合： \begin{equation*} C(p) = \sum\_{n=1}\^{}{N} c\_{T\_n}\, o\_{T\_n}\, I(p) \prod\_{i=1}\^{}{n-1}(1-o\_{T\_i}\, I(p)). \end{equation*} 训练结束后 $\sigma \to 0$、$o \to 1$，渲染公式退化为 $C(p) = c_{T_n}\, I(p)$，即每个像素仅需一次求值，实现零过度绘制的实时渲染。训练策略是获得高质量不透明网格的关键。平滑参数退火：$\sigma$ 从 1.0 线性退火至 0.0001，前期提供充分的梯度流、后期收敛到锐利的三角面片。不透明度调度：重定义 Sigmoid 映射为 $\phi(x) = O_t + (1-O_t)\,\mathrm{sigm}(x)$，其中不透明度下界 $O_t$ 在训练中从 0 线性增至 1，逐步把半透明面片推向完全不透明。
\textbf{锚点细化：} Triangle Splatting+：针对高斯核无法直接导入游戏引擎的局限，用"不透明三角面片"替代高斯核，核心思想是定义全局顶点集 $V$ 共享给所有三角面片 $T_m=(i,j,k)$，面片不透明度取顶点均值、面片内颜色重心插值，配合平滑参数退火（$\sigma$ 从 1.0 线性退至 0.0001）和不透明度调度（下界 $O_t$ 从 0 增至 1）将半透明软边界转为不透明锐利三角形，训练结束渲染退化为单次求值实现零过度绘制。
\textbf{方法信息：} Adobe 参与了这篇论文；代码：\url{https://github.com/trianglesplatting/triangle-splatting}；GitHub 关注度很高，约 1172 stars。



% METHOD: 2511.09944
\paragraph{TSPE-GS（半透明表面概率深度）（2025.11）}
TSPE-GS~\cite{2511_09944} 处理半透明表面提网格。TSGS 主要寻找透明物体的首表面；半透明材料更复杂，同一条视线可能同时看到外表面、内表面和背景。标准 3DGS 的单一 alpha 混合深度会把这些贡献压成一个平均值，导致 TSDF 融合时表面位置模糊。

TSPE-GS 把深度从单峰估计改成概率深度提取。具体流程是在渲染阶段沿射线均匀采样透射率，分析每个像素的 opacity 和深度分布，允许出现多个候选表面峰值；这样外表面和内部结构不会被一个平均深度合并。随后方法把这些概率深度作为几何输入逐步送入 TSDF 融合，并在融合时约束外表面和内部表面分别成形。相对 TSGS 的累积透射率窗口，TSPE-GS 更强调半透明多层结构；相对普通 TSDF 融合，它提供的是多模态深度输入而不是单张深度图。
\textbf{实现细节：} TSGS 聚焦于单层透明表面，但半透明物体（如薄纱、烟雾）存在多个表面层，深度不再唯一。TSPE-GS（2025.11） 在原始 3DGS 的训练环节引入基于透射率的概率深度框架——将透射率视为深度的累积分布函数，从中提取多峰深度并通过渐进式 TSDF 逐层重建——光度损失保持不变。TSPE-GS 解决的是半透明表面的多层重建问题。半透明物体（如薄纱、烟雾）的深度不唯一——同一射线上可能有多个表面。直觉上，如果沿射线逐步前进，每次遇到一个半透明层，透射率 $T$ 就下降一些；在完全不透明的表面处 $T$ 骤降至零。因此 TSPE-GS 的核心洞察是：沿射线的透射率 $T$ 本身可以充当深度 $d$ 的累积分布函数（CDF），即 $P(d_{p,g} < d) = 1 - T_d$，对应的概率密度函数为 $p(d) = -\mathrm{d}T_d / \mathrm{d}d \approx \alpha_d T_d$——恰与体积渲染中期望深度的权重一致。在不透明场景下，该 PDF 呈单峰高斯形态，峰值即表面位置；当存在半透明层时，PDF 自然呈现多峰混合高斯形态，每个峰对应一个可见表面深度 $\{\mu_k\} = \{d \mid p'(d) = 0, \, p"(d) < 0\}$。
\textbf{方法信息：} 代码：\url{https://github.com/nortonii/tspe-gs}。



% METHOD: 2512.06818
\paragraph{MeshSplatting（不透明连通网格可微渲染）（2025.12）}
MeshSplatting~\cite{2512_06818} 扩展 Triangle Splatting，把输出从三角形汤推进到连通网格。它通过 Restricted Delaunay Triangulation 在优化中维持网格连通性，并持续精炼顶点位置。

它的训练先从 SfM 点云初始化一组三角形 soup，让每个三角形像 splat 一样独立适配图像；随后根据优化后的顶点构建受限 Delaunay 三角剖分，只保留与当前表面相交的 Delaunay faces，把松散三角形逐步转成共享顶点的连通 mesh。顶点携带位置、颜色和 opacity，多个三角形共享同一顶点时，颜色和透明度梯度会在相邻面片之间传播，因此输出比独立三角形汤更平滑。后期训练会推动 opacity 收敛到不透明，使最终网格能直接使用标准引擎的深度缓冲、遮挡剔除和碰撞管线。相对 Triangle Splatting+ 的半连通共享顶点，MeshSplatting 更强调受限 Delaunay 带来的连通结构；相对 Gaussian Mesh Renderer，它不是把已有 mesh 转成高斯代理，而是端到端优化出可用 mesh。
\textbf{实现细节：} MeshSplatting（2025.12） 基于 Triangle Splatting 的顶点共享表示，通过两阶段优化将无结构三角面片汤（triangle soup）转换为连通网格。阶段一——三角面片汤优化：从 SfM 稀疏点云初始化等边三角形（大小正比于最近邻距离的均值），初始不透明度设为 0.28。此阶段不施加任何连通性约束，每个三角形自由移动，行为类似基于点的泼溅，能快速覆盖可见场景。与原始 Triangle Splatting 不同的是，MeshSplatting 在面片内部插值顶点属性（颜色、不透明度），而非保持均匀。阶段二——网格创建与精调：在三角面片汤收敛后，对其执行受限 Delaunay 三角剖分（restricted Delaunay triangulation）。该操作先计算标准 Delaunay 四面体化，再识别对偶 Voronoi 边与输入面片汤表面相交的四面体面。关键特性：不引入新顶点、不修改已有顶点位置，仅建立连通关系，保留阶段一优化得到的空间精度和外观。获得连通性后继续优化，共享顶点的梯度由所有相邻面片累积，确保每个顶点根据所有相邻面片一致更新。与 Triangle Splatting+ 类似，MeshSplatting 也采用不透明度调度（$O_t$ 从 0 线性增至 1）和平滑参数 $\sigma$ 的全局退火策略（1.0 至 0.0001），将半透明软边界逐步转为不透明锐利三角形。致密化同样基于 MCMC 采样加中点细分，但连通性使得每次细分仅新增 6 个顶点（连通时），而非无结构时的 12 个。剪枝与损失函数。
\textbf{锚点细化：} MeshSplatting：基于 Triangle Splatting 的顶点共享表示，通过\textbf{两阶段优化}将无结构三角面片汤（triangle soup）转换为连通网格。
\textbf{设计动机与边界：} 相对原始 3DGS 以独立高斯和图像重建误差为主的设定，MeshSplatting 把问题落在以下可实现的改动上：MeshSplatting基于Triangle Splatting扩展：（1）背景：Triangle Splatting用三角形作为可微分喷射原语，每个三角形由三个顶点、颜色、平滑度和不透明度定义，通过2D有符号距离场的sigmoid窗口函数实现可微分光栅化。但输出三角形汤（无连通性）且三角形不完全不透明。（2）顶点共享三角形表示：定义网格顶点V={v\_i}包含3D位置、颜色和不透明度，三角形由顶点三元组定义，不透明度和颜色通过顶点属性双线性插值获得。这避免了独立三角形的冗余，实现跨三角形的光滑属性插值。（3）连通网格生成：通过限制Delaunay三角剖分（Restricted Delaunay Triangulation）在优化过程中维持网格连通性，持续精炼顶点位置而非生成三角形汤。（4）不透明度优化：设计优化策略使三角形收敛为不透明状态，实现标准游戏引擎中深度缓冲和遮挡剔除的使用。包含致密化/剪枝策略和训练损失。其设计目的不是单独增加一个模块，而是用上述表示、监督或渲染约束消除原流程在目标任务中的具体失配。提出首个端到端网格优化方法，通过限制Delaunay三角剖分生成连通网格（非三角形汤），在Mip-NeRF360上比SOTA MiLo提升+0.69dB PSNR，训练速度快2倍，内存减少2倍，实现与游戏引擎的无缝实时交互。适用时还应保留论文给出的边界：网格表示可能丢失3DGS的毛发等复杂细节；三角剖分更新策略需要仔细平衡；对极薄结构可能仍有困难。
\textbf{方法信息：} Adobe 参与了这篇论文；代码：\url{https://github.com/meshsplatting/mesh-splatting}；GitHub 关注度很高，约 653 stars。



% METHOD: 2602.11577
\paragraph{LeafFit（植物叶片模板网格）（2026.02）}
LeafFit~\cite{2602_11577} 是植物资产的专门网格化方法。植物 3DGS 往往包含大量重复叶片，用普通点云或 mesh 提取会存储大、编辑难。LeafFit 利用叶片形状重复性，把非结构化 3DGS 转成共享模板和变形参数。

流程先从训练好的植物 3DGS 中分割单片叶子；自动分割不稳时，系统允许用户用少量交互修正。随后选择一个代表性叶片组，把它转成薄而锐利的模板网格，并把每个目标叶片表示为模板到目标的可微 Moving Least Squares 变形。优化时，模板顶点通过 MLS 控制点映射到目标叶片，渲染或几何误差回传到变形参数；运行时则在顶点着色器里实时评估 MLS，只需存储共享模板和每片叶子的变形参数。相对通用网格提取，LeafFit 的关键不是追求任意场景的完整 mesh，而是利用植物“许多叶片形状相近”的实例重复结构，把 GB 级高斯资产压成可编辑的实例化网格。
\textbf{实现细节：} GS2Mesh（2024.04） 的核心思想是：不直接从高斯基元提取几何，而是将 3DGS 渲染的立体图像对输入预训练立体匹配网络，借其大规模数据训练的强几何先验来重建 mesh——即用"立体视觉先验"替代"高斯基元直接重建"作为几何推断路径，从而绕开 3DGS 体积元素带来的噪声和伪影。GS2Mesh 利用预训练立体匹配模型从 3DGS 渲染的立体图像对中提取高质量 mesh。GS2Mesh 的出发点是一个关键观察：标准 3DGS 本身的几何质量有限——高斯基元是体积元素而非表面元素，从中直接提取 mesh 往往充满噪声和伪影。但如果用 3DGS 渲染出的高质量图像（而非高斯基元本身）作为输入，再通过成熟的立体视觉管线重建几何，就能借助立体匹配网络在大规模数据上学习的强几何先验来弥补 3DGS 的几何缺陷。GS2Mesh 的流程分为三个阶段。第一阶段是立体对渲染：对每个训练视角位姿 $(R_L, T_L)$，将左相机设为训练位姿本身（确保渲染质量最高），右相机通过水平基线 $b$ 偏移构造： 这样产生的立体对天然满足立体标定条件。基线 $b$ 的设置需要在两个矛盾因素间权衡：基线越大，立体匹配的深度精度越高（"甜点区"越宽）；但基线越大，3DGS 新视角渲染的质量越差（偏离训练视角越远，渲染噪声越大）。第二阶段是立体深度估计：GS2Mesh 用预训练的立体匹配网络（实验中 DLNR + Middlebury 预训练权重效果最佳）估计每对立体图的视差图。深度误差与视差误差的关系为 $\sigma(Z) \propto \sigma(d) \cdot Z^{}2 / (f_x \cdot B)$，即深度误差与距离的平方成正比——远处的深度估计天然不如近处精确。为滤除不可靠估计，GS2Mesh 应用两类掩码：遮挡掩码（左右视差一致性检查，剔除仅单侧可见的区域）和深度范围掩码（仅保留 $4b \leq Z \leq 20b$ 的有效范围）。被遮挡区域由相邻视角的估计填补。
\textbf{方法信息：} 代码：\url{https://github.com/netbeifeng/leaf_fit}。
\textbf{设计动机与边界：} 相对原始 3DGS 以独立高斯和图像重建误差为主的设定，LeafFit 把问题落在以下可实现的改动上：LeafFit管线包含：1)叶片分割：从非结构化3DGS中分割独立叶片，支持可选的用户交互作为后备纠错；2)模板叶片选择与转换：选择代表性叶片组，转换为薄而锐利的网格作为模板；3)可微分MLS变形拟合：使用Moving Least Squares将模板网格变形拟合到所有其他叶片，变形参数化确保局部形状保持；4)高效运行时渲染：通过顶点着色器实时评估MLS变形，最小化存储需求（所有叶片共享同一模板网格，仅存储变形参数）。其设计目的不是单独增加一个模块，而是用上述表示、监督或渲染约束消除原流程在目标任务中的具体失配。提出利用叶片形状重复性的管线，从非结构化3DGS中分割叶片，通过可微分MLS变形将模板叶片网格拟合到所有其他叶片，运行时通过顶点着色器高效评估变形，显著减少存储并支持参数级编辑。适用时还应保留论文给出的边界：对非典型叶片形状（如卷曲、受损叶片）拟合质量可能下降；分割步骤对密集重叠叶片可能有困难；仅适用于植物类结构。
% METHOD: 2602.14493
\paragraph{Gaussian Mesh Renderer（网格到高斯解析转换）（2026.02）}
Gaussian Mesh Renderer~\cite{2602_14493} 反过来使用 3DGS 的渲染器。SuGaR、MeshSplatting 一类方法通常问“怎样从高斯得到 mesh”；GMR 问的是“已有三角网格能否借用高斯 rasterizer 获得更平滑的可微梯度”。它面向固定拓扑 mesh 优化，把传统 SoftRas/Nvdiffrast 边界附近梯度稀疏的问题改写成高斯覆盖度优化问题。

对每个三角形，GMR 在面片平面内建立局部正交坐标系，把三个顶点投到局部二维坐标，解析计算均匀三角形分布的一阶矩和二阶矩；均值成为高斯中心，二维协方差再沿局部切向基嵌入三维，法向方向只给一个很小厚度。这样得到的高斯参数都是三角形顶点的可微函数，颜色、轮廓和 mask 损失通过 3DGS/gsplat 光栅化回传时，不只在三角形边界有梯度，而是在整个投影 footprint 内形成更平滑的更新。训练仍优化 mesh 顶点，并可叠加颜色、silhouette、边长和 Laplacian 平滑损失；相对 Triangle Splatting+ 直接把三角形作为 splat 图元，GMR 更像把每个三角形“高斯化”为一个可微代理。
\textbf{方法信息：} 代码：\url{https://github.com/huntorochi/gaussian-mesh-renderer}。
\textbf{设计动机与边界：} 相对原始 3DGS 以独立高斯和图像重建误差为主的设定，Gaussian Mesh Renderer for Lightweight Differentiable Rendering 把问题落在以下可实现的改动上：GMR的核心是网格到高斯的解析转换。对每个三角形面片：首先在面片平面建立局部正交坐标系（原点在顶点v\_i，x轴指向v\_j-v\_i方向，y轴通过Gram-Schmidt正交化获得，法向量为隐含z轴）；然后计算面片上均匀分布的2D协方差，通过三角形二阶矩E[x²]、E[xy]、E[y²]获得；最终将2D协方差嵌入3D空间构建完整的3D高斯。高斯的均值设为三角形质心。这种解析转换保证了梯度从渲染损失到网格顶点的平滑流动，因为高斯参数是三角形顶点的可微函数。渲染过程复用3DGS的光栅化管线。其设计目的不是单独增加一个模块，而是用上述表示、监督或渲染约束消除原流程在目标任务中的具体失配。提出GMR，将每个网格三角形解析转换为各向异性高斯基元，保持结构保真度并实现梯度流动。相比传统网格渲染器获得更平滑的梯度，尤其在有限内存的小批量设置下实现更好的优化。实现了3DGS与网格表示的紧密集成。适用时还应保留论文给出的边界：高斯近似对较大或极端形状的三角形精度可能下降；未探索与神经辐射场的结合；渲染质量可能略低于原生3DGS。
% METHOD: 2603.24725
\paragraph{Confidence-Based Mesh Extraction（置信度自监督提网格）（2026.03）}
Confidence-Based Mesh Extraction~\cite{2603_24725} 在 3DGS 提网格中加入每高斯可学习置信度。不同高斯对光度和几何监督的可信程度不同，统一加权会让漂浮或不稳定高斯污染表面。

方法给每个高斯增加一个可学习 confidence，用它在光度监督和几何监督之间动态分配权重：如果某个高斯主要解释视角相关外观或不稳定区域，就不应把它的几何信号强行用于表面；如果它在多视角下稳定，则可以更强地约束法线和表面。围绕这个 confidence，方法加入逐 primitive 的颜色方差和法线方差惩罚，要求同一高斯内部的颜色与法线不要剧烈变化，从而减少由高光、反射或错误外观解释引起的碎网格。外观模型上，它还拆开 D-SSIM 的组成项，让结构相似性监督更细粒度地作用于渲染误差。相对 SuGaR，它不依赖多视图采样再 Poisson；相对 ARGS，它的核心不是有效秩或 SDF 正则，而是让模型自己学习“哪些高斯的几何可被信任”。
\textbf{实现细节：} Confidence-Based Mesh Extraction（ConfMesh）（2026.03） 针对 3DGS mesh 提取中的一个核心矛盾：在视角依赖外观强烈的区域（如反射面、树叶），光度损失占据主导地位，大梯度幅值导致过度致密化——小高斯被反复克隆，用几何"伪造"视角依赖外观，导致提取的 mesh 充满伪影。——小高斯被反复克隆。在这些高基元密度区域，3DGS 尤其容易用几何来"伪造"视角依赖的外观，导致提取的 mesh 充满伪影。此前方法在平衡光度与几何损失时捉襟见肘：统一加权无法适应不同区域的可靠性差异。ConfMesh 的核心设计是可学习的逐基元置信度。每个高斯基元被赋予一个可学习标量 $\mathrm{conf}_i$（初始化为 0），通过指数激活 $\mathrm{conf}_{a_i} = \exp(\mathrm{conf}_i)$ 得到正值置信度，初始值恰好为 1。置信度图通过 $\alpha$-混合 渲染到像素空间： 即每个像素的置信度是射线上各高斯置信度的加权混合，权重与颜色渲染完全一致。渲染出的置信度图能有效识别反射表面、细薄植被、稀疏观测区域——这些区域的置信度自动降低，但仍然保持良好的重建质量。置信度感知损失。
\textbf{设计动机与边界：} 相对原始 3DGS 以独立高斯和图像重建误差为主的设定，Confidence-Based Mesh Extraction from 3D Gaussians 把问题落在以下可实现的改动上：核心是自监督置信度框架：为每个高斯分配可学习置信度值，动态平衡光度监督和几何监督的权重。扩展部分包括：(1)逐原语颜色和法线方差损失：惩罚每个高斯原语内部的颜色和法线不一致性，改善表面提取质量。(2)改进的外观模型：解耦D-SSIM损失的各个项，提供更精细的外观监督。方法避免了多视图技术、迭代网格提取或大型预训练模型的额外开销，保持了3DGS的固有高效性。其设计目的不是单独增加一个模块，而是用上述表示、监督或渲染约束消除原流程在目标任务中的具体失配。提出基于自监督置信度的3DGS网格提取框架。可学习置信度值动态平衡光度和几何监督。引入逐原语颜色和法线方差惩罚损失，以及改进的外观模型（解耦D-SSIM损失项）。在保持高效的同时实现无界网格的SOTA结果。适用时还应保留论文给出的边界：论文key\_sections为空，方法细节有限；置信度框架的泛化能力待验证。
\textbf{方法信息：} Huawei 参与了这篇论文；代码：\url{https://github.com/r4dl/CoMe}。



\section{SDF与隐式表面引导}

SDF/UDF 路线试图补足高斯的离散性：高斯渲染快，但表面可能有洞和噪声；SDF 连续平滑，适合 Marching Cubes，却训练慢。相关方法要么让 SDF 引导高斯，要么让高斯监督 SDF，要么把两者交替增强。


% METHOD: 2403.16964
\paragraph{GSDF（3DGS 与 SDF 双分支）（2024.03）}
GSDF~\cite{2403_16964} 把高效渲染的 GS 分支和负责表面的 SDF 分支放在一起。它不强制高斯变扁，而是通过互相引导让两种表示对齐。

GSDF 使用双分支框架。GS 分支以 3DGS/Scaffold-GS 为基础，保留快速 rasterization 和高质量外观；SDF 分支借鉴 NeuS/Neuralangelo 这类隐式表面训练，负责提供连续零水平集。二者不是简单并排训练，而是互相交换几何信号：GS 渲染出的深度图会引导 SDF 分支沿射线在表面附近更密集采样，减少无效体采样；SDF 分支估计的几何结构反过来参与高斯密度控制，使高斯分裂、保留和约束更贴近真实表面；两条分支还通过深度、法线等属性做细粒度一致性正则。相对 2DGS 直接把图元改成 surfel，GSDF 保留 3D 高斯的外观容量；相对 3DGSR 的轻量 SDF-to-opacity 连接，GSDF 更像一个完整的双系统互引导方案。
\textbf{实现细节：} GSDF（2024.03） 在原始 3DGS 的训练过程中引入了并行的 SDF 隐式场分支，用 NeuS + 多分辨率哈希网格作为独立的几何重建分支，与基于 Scaffold-GS 的渲染分支协同训练，光栅化渲染管线保持不变。动机。单独优化 3DGS 会产生悬浮的浮动物（floaters），单独优化 SDF 会在表面留下孔洞——两种表示各有盲区，而两者恰好互补。其出发点是一个基本观察：单纯优化 3DGS 会产生悬浮在空中的浮动物（floaters），单纯优化 SDF 会在学习到的表面上留下孔洞，两种表示各有盲区。GSDF 的解决方案是构建双分支架构，3DGS 分支负责高效渲染，SDF 分支负责几何重建，两者通过三条互相引导机制协同工作。第一条机制是深度引导的射线采样（GS $\to$ SDF）。直觉上，SDF 分支需要知道"表面大概在哪里"才能在正确的位置密集采样，但 SDF 本身在训练初期并不准确。GSDF 的巧妙之处在于：3DGS 虽然几何不精确，但渲染的深度图至少给出了表面位置的粗略估计——这个估计"便宜"且"即时"。在标准 NeuS 中，SDF 分支需要自己预测 SDF 值来指导射线采样，这陷入了"先有鸡还是先有蛋"的困境：SDF 还没学好时采样位置不准确，采样不准又导致 SDF 学不好。GSDF 直接利用 3DGS 分支渲染的深度图来引导采样。具体地，对于从相机中心 $\mathbf{o}$ 沿方向 $\mathbf{v}$ 发出的射线，3DGS 分支渲染的深度值为 其中 $N$ 是射线遇到的高斯数量，$d_i$ 是第 $i$ 个高斯到相机的距离，$\alpha_i$ 是其不透明度。虽然高斯渲染的深度不如传感器深度精确，但计算代价极低且无需任何额外硬件。获得 $D$ 后，SDF 分支不再在整个射线范围内均匀采样，而是将采样范围集中在 $\mathbf{o} + D\mathbf{v}$ 附近。
\textbf{实现补充：} 实现上还要核对输入、输出、参数、特征、初始化、训练损失、正则、渲染或推理接口，以及优化中的剪枝、分裂、采样、编码或解码步骤；这些抓手决定该方法到底改了 3DGS 管线的哪一环。
\textbf{锚点细化：} GSDF：针对单独优化 3DGS 产生 floaters、单独优化 SDF 表面留孔洞的局限，用"3DGS-SDF 双分支协同训练"替代单分支优化，核心思想是通过三条互相引导机制形成正向循环——3DGS 渲染深度引导 SDF 的射线采样避开空旷区域、SDF 零水平集指导致密化与剪枝让 3DGS 高斯合理分布、两分支深度/法线图互监督保证一致性。
\textbf{方法信息：} 代码：\url{https://github.com/city-super/GSDF}；GitHub 关注度很高，约 468 stars；引用很多，Semantic Scholar 约 108 次。



% METHOD: 2403.17888
\paragraph{2DGS（几何准确的平面高斯）（2024.03）}
2D Gaussian Splatting~\cite{2403_17888} 是几何章节的核心锚点。原始 3DGS 的各向异性椭球更像一个局部体积密度，渲染外观很强，但同一个 3D 高斯从不同视角投影时不一定对应同一片真实表面，深度图和法线图会出现多视角不一致。2DGS 的改变是把基元从 3D 体高斯改成 2D 定向平面高斯盘，也就是可优化的 surfel：每个基元有中心、两个切向尺度、朝向、透明度和颜色，天然带有表面法线。

渲染时，2DGS 不再只把 3D 协方差投影成屏幕椭圆，而是对像素射线和高斯平面做 ray-splat intersection，在局部平面坐标中计算 2D 高斯权重，再沿射线按 alpha compositing 混合。这个透视正确的交叉过程使同一片高斯盘在不同视角下仍然表示同一局部表面。为了让训练稳定，方法继续使用 3DGS 的可微光栅化框架和自适应密度控制，但把 densification 的梯度处理改成适合 2D surfel 的形式；训练中渲染颜色、深度、法线和 depth distortion map。

两个正则项决定了它为什么能提几何。Depth distortion loss 惩罚同一像素射线上的贡献分散，迫使可见基元在深度上集中成薄表面，减少漂浮中间高斯；normal consistency loss 让渲染法线和渲染深度梯度导出的表面法线一致，避免 surfel 朝向与实际深度面脱节。mesh 提取时，2DGS 渲染训练视角深度图，再用 TSDF 融合得到表面；后续许多逆渲染、重光照和网格提取工作把 2DGS 作为几何基础，原因正是它把“可看见的高斯”改成了“更像表面的高斯”。
\textbf{锚点细化：} 2DGS：针对 3D 椭球没有明确定义表面法线导致无法提取一致几何的局限，用"2D 面元（surfel）"替代 3D 椭球，核心思想是将椭球坍缩为盘状面元去掉第三缩放轴，每个 surfel 自然代表一个表面片段，配合基于齐次坐标的精确射线-片元交叉（避免仿射近似误差）和深度扭曲+法线一致性正则化，在 DTU 和 Tanks and Temples 上显著超越所有 3DGS 方法成为表面重建的新基线。
\textbf{实现细节：} 第一阶段：灰度重建。标准 3DGS 使用 3D 高斯基元，但 SEM 图像缺乏纹理特征，导致 3D 高斯的几何拟合质量差，产生大量飘浮伪影（floater）。ArCSEM 选择 2DGS（2024.03） 替代 3DGS 作为骨干——2D 高斯圆盘天然贴合表面，在无纹理场景下能获得更精确的几何和深度图。相机标定方面，SEM 实际使用正交投影（平行电子束），但为兼容现有新视角合成管线，ArCSEM 以透视投影近似，估计约 50,000 像素的超大焦距，相机远离场景内容。COLMAP 等传统 SfM 工具在 SEM 图像上严重失效，最终使用 RealityCapture 以共享内参+逐图畸变系数的策略获得高精度标定。此外，SEM 图像间存在显著光照变化，ArCSEM 引入仿射颜色变换（Affine Color Transformation, ACT）来补偿：对每个高斯的 SH 解码亮度 $L$，在光栅化前施加逐图像的仿射变换 $L' = WL + b$，其中 $W = \{w_1, w_2, w_3\}$ 和 $b = \{b_1, b_2, b_3\}$ 是可学习参数。渲染新视角时，取训练视角的权重和偏置均值，获得稳定一致的照明。第二阶段：着色传播。灰度模型只提供几何，艺术家仅对极少视角手动着色。ArCSEM 的关键问题是如何将稀疏着色信息传播到所有训练视角。传播机制由三个损失函数协同驱动： （1）伪颜色监督。利用灰度模型的深度图，将有色视角的像素反投影到 3D 空间，对无色视角的每个像素，搜索最近的 3D 有色点作为伪颜色 $\hat{c}_{pc}$，通过 L1 损失监督渲染结果： 若某像素在给定半径内找不到有色点，则将其排除在损失之外。此外，背景区域的伪颜色噪声严重，ArCSEM 使用 SAM（Segment Anything Model）生成前景掩模，仅对前景计算伪颜色损失。（2）模板对应模块（Template-based Correspondence Module, TCM）。伪颜色只考虑空间邻近度，无法为被遮挡区域推断颜色。
\textbf{方法信息：} 官方代码：\url{https://github.com/hbb1/2d-gaussian-splatting}；GitHub 关注度高，约 3224 stars；引用很多，Semantic Scholar 约 1410 次。



% METHOD: 2404.00409
\paragraph{3DGSR（SDF 到 opacity 转换）（2024.04）}
3DGSR~\cite{2404_00409} 在 3DGS 内部集成轻量神经 SDF。核心桥梁是 SDF-to-opacity 可微转换函数，使高斯 opacity 受 SDF 值控制，从而把高斯分布推向表面附近。

方法先用多分辨率 hash grid 加一个轻量 MLP 表示 SDF，再研究两种连接方式。紧耦合把 SDF 值通过可微 bell-shaped 函数转换成对应高斯的 opacity，直觉是越接近表面的点应越不透明；这样外观损失会直接更新 SDF，但高频纹理也可能把 SDF 拉出噪声。松耦合则把 SDF 零水平集作为几何教师，只约束高斯靠近表面并让高斯最短轴方向对齐表面法线，给 3DGS 留出表达外观细节的自由度。由于高斯只在离散位置给 SDF 监督，3DGSR 还从 SDF 做体渲染得到深度/法线，再与 3DGS 导出的几何属性对齐，把监督扩展到高斯未覆盖区域。最终从 SDF 零水平集提取网格。相对 GSDF 的双分支互引导，它更强调一个轻量、可微的 SDF-高斯连接函数；论文给出代码发布地址，但当前对应仓库还不可访问，需要后续跟踪。
\textbf{方法信息：} 引用很多，Semantic Scholar 约 116 次；论文声明代码将发布，但当前未确认完整公开代码。
\textbf{设计动机与边界：} 相对原始 3DGS 以独立高斯和图像重建误差为主的设定，3DGSR 把问题落在以下可实现的改动上：方法核心是在3D Gaussians中集成轻量级神经隐式SDF场，实现表面与外观的统一优化和对齐。包含两个关键设计：表面与Gaussians对齐和体渲染正则化。对齐策略方面，探索了紧密耦合（SDF直接控制Gaussian不透明度的可微转换函数）和松散耦合（SDF表面约束Gaussian分布和方向）两种方式。SDF-to-opacity转换函数使Gaussian不透明度受SDF值控制，确保Gaussian分布在表面附近。体渲染正则化通过从神经隐式SDF渲染深度和法线，与从3DGS导出的几何属性对齐，为SDF提供离散Gaussian未覆盖区域的稠密监督信号。训练过程利用3DGS的渲染损失优化Gaussians，同时通过一致性正则化优化SDF，无需耗时体渲染颜色预测分支。最终从SDF零水平集提取高质量表面。其设计目的不是单独增加一个模块，而是用上述表示、监督或渲染约束消除原流程在目标任务中的具体失配。提出3DGSR方法，在3D Gaussians中集成隐式符号距离场(SDF)实现统一优化。引入可微SDF到不透明度转换函数连接SDF和3D Gaussians；通过体渲染一致性正则化为SDF提供稠密监督信号。实现高质量3D表面重建同时保持3DGS的渲染质量和效率，训练过程比隐式表面方法快20倍以上。适用时还应保留论文给出的边界：需要额外MLP网络存储SDF；紧密耦合策略可能限制Gaussian在非表面区域的表达能力；对薄结构和遮挡边界的重建仍有挑战；初始化质量影响最终结果。
% METHOD: 2410.14189
\paragraph{Pulled-Gaussian SDF（牵引到零水平集）（2024.10）}
Neural SDF Inference through Splatting 3D Gaussians Pulled on Zero-Level Set~\cite{2410_14189} 的核心是可微牵引操作。离散高斯可能偏离真实表面，直接从它们推断连续 SDF 会不稳定；方法用神经 SDF 计算高斯到零水平集的最近点，把高斯拉回表面。

在每次优化中，方法查询高斯中心处的 SDF 值和梯度，并沿 SDF 梯度方向把中心投影到近似零水平集位置，可写成“中心位置减去距离值乘以归一化梯度”的形式。渲染时使用被牵引后的高斯进行 splatting，RGB 重建损失会同时更新高斯属性和 SDF 网络；几何损失则要求被牵引高斯、SDF 零水平集和局部法线保持一致。这样，高斯不只是监督 SDF 的稀疏点，而是在训练过程中不断把 SDF 表面附近的空间拉紧。最终仍用 Marching Cubes 从 SDF 中提取网格。相对 3DGSR 的 SDF-to-opacity，它显式改变高斯参与渲染的位置；相对 GSDF 的双分支互引，它更强调一个可微的“拉到表面”算子。
\textbf{实现细节：} Neural Pulling GS（2024.10） 提出一种 GS-SDF 双向融合机制：通过可微牵引（pulling）操作将 3D 高斯中心动态对齐到神经 SDF 的零水平集，同时利用牵引后的高斯盘面约束 SDF 学习。方法包含一个 MLP 表示的 SDF 场 $f$（8 层、256 隐藏单元，初始化为球体）和一组可学习的 3D 高斯 $\{g_j\}$，二者通过五个协同损失联合优化。牵引操作。给定高斯中心 $\boldsymbol{\mu}_j$，牵引操作将其投影到 SDF 的零水平集上： 直观理解：在 $\boldsymbol{\mu}_j$ 处查询 SDF 值 $s_j$ 和梯度方向，沿梯度方向移动 $s_j$ 的距离，高斯中心就被"拉"到了表面上。梯度归一化保证移动距离恰好等于 SDF 值，这是牵引操作收敛的理论基础。牵引后的高斯 $g_j'$ 与原始高斯 $g_j$ 共享除中心外的所有属性（颜色、不透明度、形状），因此可以直接送入 3DGS 的可微光栅化进行渲染。薄盘约束与切平面对齐。标准 3D 高斯是椭球体，过于"臃肿"不适合表示薄表面。Neural Pulling GS 引入薄盘约束，惩罚三个轴方差中的最小值： 这一约束将椭球体"压扁"成薄盘，使高斯天然贴合表面。盘面法线由最小方差对应轴的方向 $\mathbf{n}_j$ 给出。基于薄盘形状，切平面约束进一步对齐高斯法线与牵引中心处的 SDF 梯度： 即确保高斯盘面是零水平集在 $\boldsymbol{\mu}_j'$ 处的切平面。盘面牵引与 SDF 推理。与传统的点到点牵引不同，Neural Pulling GS 将空间中随机采样的查询点 $q$ 牵引到最近的高斯盘面（而非盘面中心），以应对高斯分布稀疏且不均匀的挑战： 其中 $q'=q - s \cdot \nabla f(q)/|\nabla f(q)|$ 是牵引后的查询点，$g^{*}$ 是距 $q$ 最近的牵引高斯。将牵引目标从"中心点"改为"盘面"，可以有效覆盖高斯中心之间的几何空隙。
\textbf{锚点细化：} Neural Pulling GS：针对离散高斯与连续 SDF 表面存在本质差异的局限，用"可微牵引操作"替代独立的 SDF-GS 优化，核心思想是将高斯中心沿 SDF 梯度方向移动 $|SDF|$ 距离投影到零水平集（$\boldsymbol{\mu}' = \boldsymbol{\mu} - s \cdot \nabla f / |\nabla f|$），引入薄盘约束与切平面对齐使椭球贴合表面，并将空间随机点牵引到最近高斯盘面（而非盘面中心）覆盖中心间几何空隙。
\textbf{设计动机与边界：} 相对原始 3DGS 以独立高斯和图像重建误差为主的设定，Neural Signed Distance Function Inference through Splatting 3D Gaussians Pulled on Zero-Le 把问题落在以下可实现的改动上：方法核心是将离散的3D高斯体与连续的神经SDF通过可微操作桥接：(1)给定3D高斯体，通过神经SDF计算其到零水平集的最近点，实现可微牵引操作将高斯体对齐到表面；(2)将对齐后的高斯体通过可微分光栅化渲染；(3)同时通过牵引操作将高斯体邻近空间拉向零水平集，逐步精化表面附近的符号距离场；(4)联合优化RGB渲染损失和几何约束损失，同时优化3D高斯体属性和神经SDF网络参数。最终通过Marching Cubes从SDF提取网格表面。其设计目的不是单独增加一个模块，而是用上述表示、监督或渲染约束消除原流程在目标任务中的具体失配。提出可微牵引操作将3DGS与神经SDF无缝融合，通过零水平集对齐和多视图一致性约束，恢复更精确、平滑且完整的表面适用时还应保留论文给出的边界：牵引操作增加计算开销；依赖SDF网络的容量和初始化；对弱纹理区域仍有挑战
\textbf{方法信息：} 代码：\url{https://github.com/wen-yuan-zhang/GS-Pull}；GitHub 关注度很高，约 153 stars；引用很多，Semantic Scholar 约 53 次。



% METHOD: 2411.15723
\paragraph{GSurf（高斯直接监督 SDF）（2024.11）}
GSurf~\cite{2411_15723} 试图去掉隐式表面方法中的冗余体渲染分支，直接用 2DGS 监督 SDF。GSDF/NeuSG 一类方法通常并行维护一个 SDF/NeuS 体渲染分支，几何质量好但训练慢；GSurf 的改变是把 2DGS 的 surface-aligned 高斯盘作为几何教师，让 SDF MLP 直接学习连续表面。

它先用 2DGS 管线得到贴近表面的高斯盘、法线和渲染深度，再训练一个 MLP 查询任意空间点的 SDF 和颜色。关键正则是不透明度熵损失：高斯 opacity 如果长期半透明，往往说明它在表面外形成体积雾，对渲染有利但对几何有害；熵正则鼓励 opacity 走向接近 0 或 1 的确定状态，并推动有用高斯盘紧贴目标表面。MLP 训练时整合高斯中心、法线、局部几何等线索，把离散 2D 高斯的表面信息变成连续 SDF。相对 3DGSR 的 SDF-to-opacity，GSurf 的连接方向相反：不是让 SDF 管高斯 opacity，而是让高斯直接监督 SDF。
\textbf{方法信息：} 代码：\url{https://github.com/xubaixinxbx/gsurf}。
\textbf{设计动机与边界：} 相对原始 3DGS 以独立高斯和图像重建误差为主的设定，GSurf 把问题落在以下可实现的改动上：方法基于2DGS（2D高斯盘）渲染管线，三个核心设计：(1) 直接高斯监督SDF：不同于先验工作（如NeuSG、GSDF）在3DGS旁并行维护体渲染分支（导致2-16小时训练），GSurf直接用2DGS渲染监督SDF训练，消除体渲染的冗余。SDF由MLP参数化，在每个采样点查询SDF值和颜色。(2) 不透明度熵正则化：原始3DGS中离表面半透明高斯体对渲染有利但损害3D重建。设计熵损失鼓励高斯体不透明度趋近1，迫使高斯盘紧密贴合目标表面，减少冗余透明基元。(3) 几何线索整合：MLP整合来自高斯体的几何信息（位置、法线等），辅助SDF预测。颜色由MLP预测而非球谐系数。整体训练效率高，2DGS渲染提供快速前向和梯度传播。其设计目的不是单独增加一个模块，而是用上述表示、监督或渲染约束消除原流程在目标任务中的具体失配。提出GSurf，首个端到端从高斯基元直接学习SDF的方法。核心创新：(1) SDF直接在3DGS监督下训练，消除冗余体渲染分支；(2) 基于熵的不透明度正则化，减少冗余透明高斯体，迫使高斯体对齐表面；(3) MLP整合几何线索。SDF的连续平滑特性解决了3DGS因噪声或缺失深度数据产生的孔洞问题。训练和渲染速度更快，重建质量可比VolSDF和NeuS等隐式表面方法。适用时还应保留论文给出的边界：依赖2DGS渲染管线的质量；不透明度熵正则化的强度需要调优；对弱纹理区域可能产生过平滑结果。
% METHOD: 2503.05182
\paragraph{MGSR（2D/3D 双高斯互增强）（2025.03）}
MGSR~\cite{2503_05182} 把 2DGS 和 3DGS 做成互相监督的双分支。2DGS 的表面几何更稳定，但对复杂反射、透射和光照变化的外观表达有限；3DGS 的外观容量强，却容易用错误高斯去解释反射，从而污染几何。MGSR 的目标是让一个分支负责“表面可信”，另一个分支负责“外观可分解”。

3DGS 分支在普通颜色 SH 之外加入反射 opacity、反射置信度和反射 SH，把观测颜色拆成透射/固有外观与反射分量。2DGS 分支则维持 surface-aligned 的几何，并渲染深度图监督 3DGS，使反射分解不能随意改变物体表面。反过来，3DGS 分支可以渲染去反射后的图像，作为 2DGS 的更干净外观监督，帮助 2DGS 在变化光照下仍保持稳定表面。训练通常先分别 warm-up 两个分支，再交替优化，最后用 2DGS 的深度或 TSDF 融合提取网格。相对 GSDF 的 SDF/GS 互引，MGSR 的互补对象变成了 2D 高斯表面分支和 3D 高斯反射分解分支。
\textbf{实现细节：} GSurf 证明了 2DGS 的光栅化足以替代体渲染来监督 SDF 学习，但仅在单一光照条件下验证。MGSR（2025.03） 将 2DGS 的表面重建能力扩展到多光照条件——用 3DGS 分支处理光照分解（透射/反射分离），2DGS 分支负责精确几何，两个分支通过深度图交叉监督互增强。核心观察是：3DGS 的体积高斯表达力强，擅长渲染但几何不准；2DGS 的 surfel 天然对齐表面，擅长几何但渲染能力受限——两者恰好互补。此外，多光照场景中的透明/反射物体使问题进一步复杂化，单一分支无法同时处理光照分离和几何重建。3DGS 分支：光照分解。MGSR 在标准 3DGS 基础上引入三个反射相关参数——反射不透明度 $\alpha_{\mathrm{ref}}$、反射置信度 $\rho \in [0,1]$、反射球谐系数 $C_{\mathrm{ref}}^k$——将场景建模为透射分量和反射分量的叠加。渲染时，沿射线累积反射置信度得到像素级反射权重 $W = \sum_{i=1}^K \rho_i \alpha_i \prod_{j=1}^{k-1}(1-\alpha_j)$，再分别通过 $\alpha$-混合 计算透射颜色 $C_{\mathrm{tran}}$ 和反射颜色 $C_{\mathrm{ref}}$，最终合成颜色为 $C = C_{\mathrm{tran}} + W \cdot C_{\mathrm{ref}}$。
\textbf{方法信息：} 代码：\url{https://github.com/tsingyuanchou/mgsr}。
\textbf{设计动机与边界：} 相对原始 3DGS 以独立高斯和图像重建误差为主的设定，MGSR 把问题落在以下可实现的改动上：框架包含两个分支和互增强机制：(1) 3DGS分支：引入反射相关参数（反射不透明度、反射置信度、反射SH），将场景建模为透射和反射分量，实现光照分解。通过几何引导的光照分解模块利用2DGS分支的深度信息提升渲染性能。(2) 2DGS分支：将2D高斯盘与表面对齐，即使在显著变化的光照条件下也能捕获初始且相对准确的几何信息。(3) 互增强交替优化：两分支先独立热身，然后3DGS分支由2DGS分支生成的深度图监督，2DGS分支由3DGS分支渲染的无反射图像监督。使用早停策略减少计算成本。最终使用TSDF提取重建网格。其设计目的不是单独增加一个模块，而是用上述表示、监督或渲染约束消除原流程在目标任务中的具体失配。提出2DGS和3DGS双分支互增强框架MGSR，解决渲染质量与表面重建精度之间的矛盾。2DGS分支提供精确几何信息，3DGS分支通过几何引导的光照分解实现多种光照下的真实感渲染。两分支交替优化、互相监督。适用时还应保留论文给出的边界：双分支训练增加计算成本；对极端镜面反射处理有限；交替优化的超参数敏感。
% METHOD: 2503.19458
\paragraph{GaussianUDF（开放表面无符号距离）（2025.03）}
GaussianUDF~\cite{2503_19458} 面向开放表面，如薄片、布料边界或非封闭物体。SDF 默认空间有内外符号，适合闭合曲面；开放表面没有明确内外，强行学习有符号距离会让边界附近不稳定。GaussianUDF 因此把目标从 SDF 改成 UDF，即只预测到表面的无符号距离。

方法先把薄而平的 2D 高斯过拟合到可见表面上，让这些高斯成为显式几何代理。之后并不直接把高斯当最终输出，而是从高斯的局部分布、法线和梯度关系中构造自监督信号，训练连续 UDF 网络。近表面区域由贴合表面的高斯提供强约束；离表面较远的位置则通过梯度推理和一致性约束传播距离信息，缓解 UDF 零水平集附近梯度方向不唯一的问题。相对 GSurf 的 SDF，它处理无内外符号的开放表面；相对 SurfaceSplat 的交替增强，它更专注于从显式高斯代理推断连续 UDF 表示。
\textbf{实现细节：} 前述方法均假设场景表面是封闭的（有明确的内外之分），但许多真实物体（如布料、叶片、开放曲面）不具有封闭表面。GaussianUDF（2025.03） 将 SDF 替换为 UDF（Unsigned Distance Function）——用无符号距离场替代有符号距离场，去掉内外判断只保留到表面的距离——以 2DGS 薄盘作为开放表面的离散近似，通过三步渐进策略桥接离散高斯与连续 UDF 的鸿沟。核心挑战在于：3D 高斯是离散的显式表示，而 UDF 是连续的隐式场——两者的本质差异使得从高斯推理 UDF 极其困难。过拟合高斯到表面。GaussianUDF 采用 2DGS 的扁平化高斯盘表示场景——这些盘足够薄，天然适合近似开放表面。除了 2DGS 本身的薄盘特性，GaussianUDF 还通过投影损失将高斯中心拉向 UDF 的零水平集：用梯度场将高斯中心 $\boldsymbol{\mu}_i$ 投影到表面 $\boldsymbol{\mu}_i' = \boldsymbol{\mu}_i - d_i \cdot \nabla f(\boldsymbol{\mu}_i) / |\nabla f(\boldsymbol{\mu}_i)|$（$d_i = f(\boldsymbol{\mu}_i)$ 为无符号距离），然后以 $\boldsymbol{\mu}_i'$ 为目标直接优化高斯位置 $L_{\mathrm{proj}} = \|\boldsymbol{\mu}_i' - \boldsymbol{\mu}_i\|_2$。
\textbf{方法信息：} 代码：\url{https://github.com/lisj575/gaussianudf}。
% METHOD: 2507.15602
\paragraph{SurfaceSplat（SDF 与 3DGS 互补迭代）（2025.07）}
SurfaceSplat~\cite{2507_15602} 把 SDF 和 3DGS 做成互补循环。稀疏视角下，3DGS 的全局几何容易散，SDF 虽平滑但细节不足；方法先用 Voxurf/NeuS 类 SDF 拟合全局结构，再从网格表面采样点初始化 3DGS。

第一步，SurfaceSplat 用 SDF 方法在稀疏图像上先恢复一个粗但全局一致的网格，并清理表面后从 mesh 上均匀采样点云，替代 COLMAP 在弱纹理或低重叠区域可能缺失的初始化点。第二步，用这些点训练 3DGS，让高斯补回 SDF 容易丢掉的局部纹理和细节。第三步，训练好的 3DGS 渲染额外新视角图像，把原始稀疏输入扩展成更密集的监督，再返回去训练更高分辨率的 SDF。实际使用中可只跑一个循环，以控制训练时间。相对 GSDF 的联合双分支，SurfaceSplat 更像交替增强管线：SDF 给 3DGS 全局形状，3DGS 给 SDF 补视角和细节。
\textbf{实现细节：} GaussianUDF 解决了开放表面问题，但所有前述方法都假设有充足的训练视角。在稀疏视角条件下，无论是 SDF 还是 3DGS 都力不从心。SurfaceSplat（2025.07） 专门针对稀疏视角场景设计了循环优化框架——用 SDF 生成的 mesh 初始化 3DGS，再用 3DGS 渲染的新视角图像回过头来增强 SDF 训练——让两种表示互补增强。其出发点是一个关键观察：SDF 方法凭借稠密表示能很好地捕获全局结构但缺乏细节，3DGS 方法凭借稀疏表示擅长保留局部高频细节但全局几何差——两者恰好互补。SDF 改善 3DGS（初始化阶段）。标准 3DGS 用 COLMAP 稀疏点云初始化，但在稀疏视角下 COLMAP 经常失败或只产生极少且不准确的点。SurfaceSplat 改用 SDF 生成的 mesh 来初始化 3DGS：先用 Voxurf 的粗阶段（$96^{}3$ 体素网格，约 15 分钟）在稀疏输入图像上训练 SDF 场 $V^{(sdf)}: \mathbb{R}^{}3 \to \mathbb{R}$，通过可微体渲染和图像重建损失 $L = L_{\mathrm{recon}} + L_{\mathrm{TV}}(V^{(sdf)}) + \lambda L_{\mathrm{smooth}}(V^{(sdf)})$ 监督。为消除稀疏视角导致的漂浮伪影，引入法线一致性损失 $L_{\mathrm{normal}} = \|\hat{\mathbf{N}}(r) - \tilde{\mathbf{N}}(r)\|_1$，用 Metric3D 预测的单目法线 $\hat{\mathbf{N}}$ 监督体渲染法线 $\tilde{\mathbf{N}}$。然后用 Marching Cubes 提取 mesh，通过连通分量聚类保留最大组件并移除漂浮碎片，再将 mesh 投影到训练视角生成深度图，从可见深度区域均匀采样点（每帧 $5k$ 点，融合后约 $50k$ 点）作为 3DGS 的初始化点云。对于背景区域，仍保留 COLMAP 稀疏点补充。这种基于 mesh 的初始化保证点云均匀分布在可见表面上，远优于 COLMAP 的稀疏不可靠结果。3DGS 增强 SDF（视角扩展阶段）。SDF 重建的核心瓶颈是训练视角不足导致的监督信号匮乏。SurfaceSplat 利用初始化良好的 3DGS 渲染额外的新视角图像来扩充训练数据。3DGS 在 $7k$ 次迭代（约 5 分钟）内即可收敛，然后从新相机位姿渲染图像。新位姿通过两种策略生成：相机位置扰动 $c_m^{*} = c + \Delta p$（在原位姿附近添加微小偏移）和相机位姿插值（用三次样条在原始相机轨迹间平滑插值位置和朝向）。
\textbf{方法信息：} ByteDance 参与了这篇论文；项目仓库 \url{https://github.com/aim-uofa/SurfaceSplat} 已存在，但当前只看到 README 和图片，完整代码仍需持续跟踪。



\section{其他几何重建方法}

几何重建的主线不只有法线正则和 SDF/UDF。很多方法从更具体的失败模式切入：漂浮高斯贡献小但不一定低透明度，曲线结构不能用普通面片概括，开放表面没有内外符号，或者局部几何需要跨形状泛化。本节把这些较窄但可信的路线放在一起，重点看它们如何改变高斯的筛选、代理几何或查询方式。



% METHOD: 2406.07499
\paragraph{TrimGS（贡献度高斯修剪）（2024.06）}
TrimGS~\cite{2406_07499} 是几何章节的贡献度修剪锚点。SuGaR、2DGS、DN-Splatter 这类方法主要通过表面正则或几何先验把高斯推向正确表面；TrimGS 从相反方向切入：如果某些高斯长期只制造漂浮物、模糊纹理或错误几何，就不应只靠 opacity 阈值等待它们消失，而应根据它们对最终渲染的真实贡献被主动修剪。

它首先重新定义“一个高斯有没有用”。原始 3DGS 的剪枝主要看不透明度，但高不透明度的 floater 也可能对少数视角造成伪影，低不透明度的小高斯却可能承载高频细节。TrimGS 因此在 CUDA rasterization 中累计每个高斯对像素的 $\alpha$-混合权重：单视图贡献来自该高斯覆盖像素上的 blending weight，并按投影面积归一化；多视图贡献则取若干高贡献视角的平均。训练过程中每隔固定迭代遍历训练视图，计算多视图贡献，移除最低贡献的一部分高斯。这一机制能识别“看起来 opacity 不低、实际几乎没为正确图像服务”的错误基元。

第二个关键是尺度控制。论文的分析认为，过大的高斯虽然容易获得较高贡献，但会模糊高频区域，而且尺度梯度噪声更大，不利于精细几何优化。TrimGS 因此修改 densification：当高斯最大尺度超过场景相关阈值时，把它拆成多个更小高斯并缩小尺度，使场景由更细的基元表达。它还可叠加 normal consistency regularization，并能接到原始 3DGS 或 2DGS 后面形成 Trim3DGS/Trim2DGS。实现上，官方代码修改了 3DGS/2DGS 的 differentiable rasterization kernel 来统计贡献度，训练通常在已有模型上做约 7K iteration 的快速 refinement，每 1000 iteration 修剪一次；mesh 仍可沿用 2DGS 的 median depth + TSDF 融合流程。
\textbf{锚点细化：} Trim 3DGS：针对以不透明度作为剪枝判据无法识别高不透明度 floaters 与重要低不透明度小高斯的局限，用"渲染贡献度驱动的渐进修剪"替代不透明度剪枝，核心思想是一个高斯"有没有用"应看它在最终渲染图像中实际贡献了多少像素颜色——贡献定义为该高斯在所有覆盖像素上的 $\alpha$-混合 权重归一化累积，配合尺度控制避免退化大高斯主导渲染。
\textbf{方法信息：} 代码：\url{https://github.com/yuxueyang1204/trimgs}；GitHub 关注度较高，约 263 stars；引用很多，Semantic Scholar 约 50 次。



% METHOD: 2410.07266
\paragraph{Spiking GS（脉冲神经元剪枝与截断）（2024.10）}
Spiking GS~\cite{2410_07266} 延续 TrimGS 的“清理无效高斯”问题，但判断对象不再只是渲染贡献度，而是高斯 opacity 和高斯函数尾部。原始 3DGS 中有两类几何噪声：一类是低 opacity 高斯，数量多但贡献小；另一类是每个高斯分布的长尾，虽然中心在表面附近，尾部仍会与其他高斯重叠，导致体积感和深度偏差。

方法把 forward-integrate-and-fire（FIF）脉冲神经元接到两个位置。全局 FIF 作用于高斯 opacity，学习一个全局阈值 $\tau^o$，低于阈值的高斯在前向中被置零，相当于可学习剪枝。局部 FIF 作用于单个高斯的函数值，为每个高斯学习截断阈值 $\tau_i^p$，把低函数值尾部截掉，只保留对局部表面最有用的主体部分。由于阈值操作不可导，训练用代理梯度近似回传；同时修改 densification，用 opacity 阈值和尺度感知 cloning 避免遮挡区域被剪得过空。相对 TrimGS 的多视图贡献度修剪，Spiking GS 更像在高斯函数形式上直接压掉无效 opacity 和长尾积分。
\textbf{实现细节：} SolidGS 从核函数形状入手改善几何，Spiking GS（2024.10） 则从另一个角度切入：在原始 3DGS 的训练环节引入脉冲神经元机制，用可学习阈值自动"熄灭"低不透明度高斯和高斯尾部，渲染流程保持不变。该方法从脉冲神经网络（Spiking Neural Network, SNN）汲取灵感，针对标准 3DGS 的一个结构性问题：高斯的低不透明度区域会产生大量冗余和伪影。标准 3DGS 通过定期剪枝（移除不透明度低于阈值 $\tau_{\alpha}$ 的高斯）和不透明度重置（将高斯的不透明度重置为低值 $\alpha_{\mathrm{reset}}$）来控制高斯密度，但这种启发式策略难以有效管理表面附近漂浮的低不透明度高斯（Low-Opacity Gaussians, LOGs）以及高斯尾部重叠（LOTs）问题。Spiking GS 引入了两种 FIF（Fast-Sigmoid Integrate-and-Fire）脉冲神经元。第一种是全局 FIF 神经元，作用于不透明度： 其中 $\tau^{\alpha}$ 是所有高斯共享的可学习阈值，在训练中通过正则化项 $L_{\alpha} = \lambda_{\alpha} \cdot \frac{1}{\tau^{\alpha}}$ 被逐步推高（$\lambda_{\alpha} = 2 \times 10^{-5}$）。直觉上，随着阈值升高，越来越多低不透明度的高斯被"熄灭"并随后剪枝，LOGs 被自动清理。Spiking GS 还将不透明度重置值从原来的 $\alpha_{\mathrm{reset}}$ 替换为更高的 $\tau^{\alpha}$，间接引导 3DGS 生成更少低不透明度高斯。第二种是局部 FIF 神经元，作用于高斯的表示函数本身： 其中每个高斯拥有独立的可学习截断阈值 $\tau_i^p$，通过 $L_p = \lambda_p \cdot \frac{1}{\bar{\tau}^p}$ 正则化（$\lambda_p = 2 \times 10^{-5}$）。这一机制将高斯尾部截断为零，消除了相邻高斯尾部重叠导致的渲染偏差。此外，Spiking GS 还引入了尺度损失 $L_s$ 来约束高斯的尺寸，并在盲区（约束不足区域）设计了基于尺度的克隆策略来补充高斯。总损失为 $L = L_c + L_d + L_N + \lambda_{\alpha} L_{\alpha} + \lambda_p L_p + \lambda_s L_s + \lambda_t L_t$，其中 $\lambda_s = 5 \times 10^{-4}$，$\lambda_t = 1.0$。这种脉冲神经元机制使最终表示更加紧凑，同时改善了表面对齐质量。
\textbf{方法信息：} 代码：\url{https://github.com/zju-bmi-lab/spikinggs}。



% METHOD: 2505.04668
\paragraph{SGCR（球面高斯曲线重建）（2025.05）}
SGCR~\cite{2505_04668} 把 3DGS 从表面重建转向 3D 特征曲线重建。传统曲线方法通常先从多视角图像提取边缘，再建立跨视角匹配或三维点云，最后拟合曲线；误差会在多个后处理步骤中累积。SGCR 的做法是保留 splatting 的可微多视角监督，但把高斯简化成适合边缘的球面基元。

它把标准 3DGS 的各向异性椭球改成固定半径的各向同性球面高斯，去掉旋转、协方差和 SH，只保留位置、半径和单通道灰度强度。训练时，多视角 2D 边缘图作为监督，球面高斯被 splat 到图像平面形成渲染边缘图，并与图像边缘对齐；由于监督来自所有训练视图，球面高斯逐渐落到 3D 曲线附近。第二阶段再从这些对齐的球面高斯中提取参数化曲线。相对 CurveGaussian，SGCR 的表示更简单、优化更像“边缘点云生成”，而曲线参数化仍是后处理步骤。
\textbf{实现细节：} 如果说 GaussianPU 关心的是"如何增加高斯"，Trim 3DGS 则关注相反的问题——"如何删除多余的高斯"。标准 3DGS 以不透明度作为剪枝判据，但这一判据存在明显漏洞：高不透明度的漂浮高斯（floaters）会被保留，而重要的低不透明度小高斯反而被误删。Trim 3DGS（2024.06） 的核心贡献是提出一个更忠实的"贡献"度量。直觉。一个高斯"有没有用"，不应该看它自身有多亮，而应该看它在最终渲染的每张图片中实际贡献了多少像素颜色。一个覆盖半幅画面的漂浮高斯虽然不透明度高，但它的颜色会被前方更近的高斯遮挡，对最终图像几乎无贡献；而一个微小的表面高斯虽然不透明度低，却可能是某个细节区域的唯一贡献者。贡献度量。基于上述直觉，高斯 $G$ 对第 $k$ 个视角的贡献定义为其 $\alpha$-混合 权重的归一化累积： 其中 $P_k$ 是 $G$ 的 2D 投影区域，$i(p)$ 是 $G$ 在像素 $p$ 处的深度排序。直觉上，这就是 $G$ 在所有它覆盖的像素上贡献的颜色比例之和。归一化因子 $|P_k|$ 消除大高斯因面积大而获得不公平高贡献的偏差。多视角总贡献为 $C = \sum_k C_k$。渐进修剪与尺度控制。修剪策略是渐进式的——每隔一定迭代步数，移除贡献最低的一批高斯，使剩余高斯逐步向表面聚集。此外，Trim 3DGS 引入尺度控制策略：限制高斯的最大尺度，避免退化的大高斯主导渲染且难以优化。配合法线正则化损失（将渲染深度图梯度法线与高斯最短轴对齐），Trim 3DGS 在不牺牲渲染质量的前提下显著改善了几何精度。
\textbf{方法信息：} 代码：\url{https://github.com/martinyxr/sgcr}。
\textbf{设计动机与边界：} 相对原始 3DGS 以独立高斯和图像重建误差为主的设定，SGCR 把问题落在以下可实现的改动上：SGCR方法包含两阶段：(1) 球面高斯生成：将标准3DGS的各向异性椭球体简化为固定半径r0的各向同性球体，去除协方差矩阵(缩放和旋转属性)和球谐系数，颜色改为单通道灰度值。通过基于视角的渲染损失优化，将渲染的2D边缘图与从图像提取的边缘图对比，无需任何3D监督；(2) 参数化曲线提取：提出优化算法从对齐的球面高斯中直接提取精确参数化曲线。球面高斯作为中间表示桥接3D描述和2D指导，解决了多视角边缘对应关系的难题。其设计目的不是单独增加一个模块，而是用上述表示、监督或渲染约束消除原流程在目标任务中的具体失配。提出球面高斯(Spherical Gaussians)表示和SGCR优化算法，首次直接从多视角图像通过2D边缘图监督重建精确3D参数化曲线，无需3D标注。适用时还应保留论文给出的边界：固定半径球面高斯可能无法精确表示细长结构特征；对2D边缘检测质量敏感；复杂拓扑结构的曲线提取仍有挑战。
% METHOD: 2506.21401
\paragraph{CurveGaussian（参数曲线重建）（2025.06）}
Curve-Aware Gaussian Splatting~\cite{2506_21401} 继续曲线重建路线，但比 SGCR 更早引入参数曲线本身。SGCR 先用球面高斯形成 3D 边缘代理，再提曲线；CurveGaussian 则让曲线和边缘高斯从训练一开始就耦合，使输出天然是可编辑的参数曲线。

方法把一条三次或一阶贝塞尔曲线关联到一组边缘高斯。曲线提供连续几何骨架，高斯负责可微 2D 边缘渲染；训练时，渲染边缘与图像边缘图对齐，梯度同时更新高斯和曲线控制点。为了处理拓扑变化，它在线性化、合并、分裂和剪枝之间切换：过长或弯曲不足的曲线段会拆分，冗余段会合并，低贡献段被删掉。相对 SGCR 的“球面高斯到曲线”两阶段流程，CurveGaussian 的核心变量就是参数曲线，高斯更像让曲线可被多视角图像监督的渲染载体。
\textbf{实现细节：} EdgeGaussians 先重建边缘点云再拟合曲线，是典型的两阶段方法。CurveGaussian（2025.06） 将其合并为单阶段端到端 3D 参数曲线重建，突破两阶段范式的误差累积瓶颈。直觉。传统方法先建点云再拟合曲线是"先画点再连线"，CurveGaussian 直接优化"线"的参数，渲染梯度自然驱动曲线对齐。曲线-高斯双向耦合。核心是每条参数曲线（三次/一阶 B\'{e}zier 曲线）在参数 $t \in [0,1]$ 上采样 $N_g$ 个点生成边缘高斯，高斯的中心、缩放和旋转由曲线控制点通过解析公式决定；渲染误差经高斯反传至曲线控制点。为处理复杂拓扑，CurveGaussian 设计了自适应控制策略：迭代 3k 步后启动曲线线性化（将弯曲角度超过 $20°$ 的曲线分裂），7k 步后执行合并和分裂操作。
\textbf{方法信息：} 代码：\url{https://github.com/zhirui-gao/curve-gaussian}。
\textbf{设计动机与边界：} 相对原始 3DGS 以独立高斯和图像重建误差为主的设定，Curve-Aware Gaussian Splatting for 3D Parametric Curve Reconstruction 把问题落在以下可实现的改动上：核心创新：(1)CurveGaussian表示：建立参数曲线与边缘导向高斯组件的双向耦合机制。每条参数曲线（三次/一阶贝塞尔曲线）与一组边缘高斯对应，使参数曲线可通过高斯泼溅实现可微渲染，同时保持几何连续性。(2)单阶段优化：直接从2D边缘图优化3D参数曲线参数，消除传统两阶段方法（先重建边缘点云，再拟合参数曲线）的误差累积。优化变量大幅减少。(3)自适应拓扑优化：训练过程中通过线性化、合并、分裂和剪枝操作动态调整曲线结构，处理复杂拓扑。使用三次贝塞尔曲线（4个控制点）和一阶贝塞尔曲线。其设计目的不是单独增加一个模块，而是用上述表示、监督或渲染约束消除原流程在目标任务中的具体失配。提出首个端到端单阶段3D参数曲线重建框架，通过参数曲线与边缘导向高斯的双向耦合机制（CurveGaussian）实现可微渲染，配合自适应拓扑优化，消除两阶段方法的误差累积。适用时还应保留论文给出的边界：参数曲线类型限制为贝塞尔曲线；对高度复杂自由曲线的表示能力有限；依赖边缘图质量。
% METHOD: 2508.09830
\paragraph{RayletDF（光线段距离场）（2025.08）}
RayletDF~\cite{2508_09830} 从点云或 3DGS 高斯出发做可泛化表面重建。SDF/UDF 通常在三维坐标上密集查询网络，表达连续但推理成本高；RayletDF 把查询单位改成局部光线段 raylet，让模型直接预测这段光线到表面的距离。

具体流程有三部分。首先在表面附近采样 raylet，并从周围点云或高斯中抽取局部几何特征；raylet 的起点靠近表面，方向给出一个局部查询坐标。然后距离场网络预测 raylet distance，也就是该光线段起点到表面命中点的有符号距离。最后，多 raylet 融合器把来自不同方向和邻域的命中点聚合成稳定表面。它的关键假设是局部 raylet 模式比全局坐标场更容易跨形状泛化，因此可在未见形状上单次前向恢复表面。相对 GaussianUDF，它不是从高斯代理学习全局 UDF，而是把高斯/点云转成局部光线段查询。
\textbf{实现细节：} 触觉和偏振都为特定区域提供了额外约束，RayletDF（2025.08） 则关注另一个维度：泛化能力。它提出一种可泛化的表面重建方法，可直接从点云或 3DGS 高斯中恢复精确表面，无需逐场景优化。核心表示：光线段。核心表示是"光线段"（raylet）——由起点 $p_l$ 和方向单位向量 $d_l$ 参数化的局部光线片段。直觉上，一个 raylet 就是一条从已知位置出发的短光线，网络的任务是预测这条光线走多远才能碰到表面。网络直接预测从光线段起点到表面交点的距离：$(d_l, s_l) = \mathrm{MLPs}(\phi_l, p_l, d_l)$，其中 $\phi_l$ 是 SparseConv 提取的局部几何特征，$s_l$ 是置信度分数。三模块管线。
\textbf{技术实现补充：} 触觉和偏振都为特定区域提供了额外约束，RayletDF（2025.08） 则关注另一个维度：泛化能力。它提出一种可泛化的表面重建方法，可直接从点云或 3DGS 高斯中恢复精确表面，无需逐场景优化。核心表示：光线段。核心表示是"光线段"（raylet）——由起点 $p_l$ 和方向单位向量 $d_l$ 参数化的局部光线片段。直觉上，一个 raylet 就是一条从已知位置出发的短光线，网络的任务是预测这条光线走多远才能碰到表面。网络直接预测从光线段起点到表面交点的距离：$(d_l, s_l) = \mathrm{MLPs}(\phi_l, p_l, d_l)$，其中 $\phi_l$ 是 SparseConv 提取的局部几何特征，$s_l$ 是置信度分数。三模块管线。管线包含三个模块：(1) Raylet 特征提取器——用 SparseConv 获取场景逐点特征，对查询光线段通过 KNN 找 $K$ 个最近邻点，拼接相对位置、距离和特征向量构成局部描述子；(2) Raylet 距离场——8 层 MLP 预测距离值和置信度，无输出层激活函数以支持正负距离；(3) 多光线段混合器——沿每条光线采样 $T$ 个光线段，用 softmax 归一化的置信度加权融合多次估计：$D = \sum_{t=1}^{T}\omega_{l_t}(\|p_{\mathrm{cam}} - p_{l_t}\| + d_{l_t})$，相当于让多个局部预测投票决定最终距离。训练时的采样策略为：基于每个点最近邻距离构建虚拟球体，选择光线穿过的球体的垂足作为光线段起点；若输入为 3DGS 高斯，则利用 $\alpha$-混合 权重最高的交点作为起点。RayletDF 在大规模场景上训练后可直接泛化到新场景，作为 3DGS 的后处理工具提取高质量表面。
\textbf{方法信息：} 代码：\url{https://github.com/vlar-group/rayletdf}。
\textbf{设计动机与边界：} 相对原始 3DGS 以独立高斯和图像重建误差为主的设定，RayletDF 把问题落在以下可实现的改动上：RayletDF流水线包含三个关键模块：(1) 光线段特征提取器：在表面附近采样光线段(raylet)，提取细粒度局部几何特征。光线段定义为光线单位段，起点位于形状表面附近；(2) 光线段距离场预测器：预测光线段距离（表面命中点与光线段起点之间的有符号距离），直接从查询光线预测表面点；(3) 多光线段融合器：聚合多个光线段预测结果重建精确表面点。核心思想是光线段距离场的局部几何模式具有良好的跨形状泛化性，使得方法可以在未见数据集上单次前向推理成功恢复3D表面。不同于坐标方法的密集网络评估需求，该方法更高效。其设计目的不是单独增加一个模块，而是用上述表示、监督或渲染约束消除原流程在目标任务中的具体失配。提出光线段距离场(raylet distance field)新技术，定义光线段为光线单位段，光线段距离为表面命中点与光线段起点之间的有符号距离。设计包含光线段特征提取器、距离预测器和多光线段融合器的流水线，实现从点云或3DGS高斯的可泛化表面重建，单次前向传播即可在未见数据集上恢复3D表面。适用时还应保留论文给出的边界：对点云密度和质量有一定要求；极端稀疏或噪声较大的输入可能影响重建质量；光线段采样的参数需要调整。

\section{多视角几何与平面约束}

多视角几何路线强调显式利用视图间一致性，而不是只相信单目深度或单视图法线。原始 3DGS 主要依靠光度误差驱动，遇到弱纹理、反光或遮挡时容易在深度和法线上漂移；平面方法则利用室内/建筑场景的大平面先验，把高斯聚合成解析平面或让深度在视图间闭合。读这一节时，应关注每个方法把几何约束接在渲染深度、法线、一致性损失还是平面参数上。
% METHOD: 2312.00206
\paragraph{SparseGS（几何、法线、表面与网格）（2023.11）}
SparseGS~\cite{2312_00206} 关注的问题是：把 3DGS 用于几何、表面、深度或网格重建 SparseGS 是宽召回补全中发现的论文，主要根据标题、摘要和 arXiv 类别判定其与 3DGS 的关系。
相对原始 3DGS，SparseGS 的变化不是只替换损失函数，而是把表示、初始化、约束或网络预测路径重新组织到任务流程中：该工作被召回是因为论文文本包含 Gaussian/3D Gaussian 相关表述；当前分类先确定其任务域和与 3DGS 的相关性，后续证据流水线再补代码、企业和引用信息。
\textbf{实现细节：} 具体实现可以按输入、核心模块和渲染/优化路径来读。输入端首先取得任务所需的图像、相机、文本、时间序列、语义或几何线索，并把这些线索转换成可以约束高斯中心、尺度、旋转、不透明度、颜色或附加特征的条件。核心结构中，方法使用的主要机制是：从摘要层面看，SparseGS 围绕“把 3DGS 用于几何、表面、深度或网格重建”展开。若其正文或代码后续可读，应进一步核对它具体改动的是表示、初始化、渲染、优化、损失、动态建模、压缩或应用接口中的哪一环。 训练或推理时，这些模块通常先给出初始化或中间特征，再通过可微 Gaussian rasterization、重建损失、感知/语义约束、几何正则或任务损失回传到高斯参数；如果方法包含前馈网络，则网络输出的是高斯属性、变形场、压缩码、语义特征或控制变量，而不是只输出最终 2D 图像。渲染路径仍以显式高斯投影、深度排序和 alpha 合成为基础，只是在采样、属性预测、局部编辑、动态变形、压缩编码或下游任务接口处加入额外约束。这样写入本章时，它与前后方法的关系是：相对原始 3DGS 和本章已有锚点方法，它可以理解为 3DGS extends\_or\_uses 宽召回分类阶段的关系判断，后续需要按论文方法部分复核。
\textbf{证据与边界：} 本轮召回分类尚未逐项核对实验表格；后续证据补全阶段需要读取论文源码/PDF或项目页确认数据集、指标和对比方法。 该条来自宽召回批次；源码下载受 arXiv e-print 网络限制，部分论文暂时只能依据标题、摘要和类别分类，需要在入可信书前做人工/LLM精读复核。 本书只采用本地 paper pack、数据库分类、源码/引用登记和论文摘要中可核验的信息；没有本地证据支撑的速度、排行榜或开源完整性不在正文中强化。\textbf{方法信息：} 代码：\url{https://github.com/ForMyCat/SparseGS}；引用信号：194 citations (6.2/month, domain-relative cutoff 4.3/month, semantic\_scholar)；代码关注：formycat/sparsegs has 228 stars (7.3/month, domain-relative cutoff 37.1/month)。
\textbf{设计动机与边界：} 相对原始 3DGS 以独立高斯和图像重建误差为主的设定，SparseGS 把问题落在以下可实现的改动上：从摘要层面看，SparseGS 围绕“把 3DGS 用于几何、表面、深度或网格重建”展开。若其正文或代码后续可读，应进一步核对它具体改动的是表示、初始化、渲染、优化、损失、动态建模、压缩或应用接口中的哪一环。其设计目的不是单独增加一个模块，而是用上述表示、监督或渲染约束消除原流程在目标任务中的具体失配。该工作被召回是因为论文文本包含 Gaussian/3D Gaussian 相关表述；当前分类先确定其任务域和与 3DGS 的相关性，后续证据流水线再补代码、企业和引用信息。适用时还应保留论文给出的边界：该条来自宽召回批次；源码下载受 arXiv e-print 网络限制，部分论文暂时只能依据标题、摘要和类别分类，需要在入可信书前做人工/LLM精读复核。
% METHOD: 2312.00846
\paragraph{NeuSG（几何、法线、表面与网格）（2023.12）}
NeuSG~\cite{2312_00846} 关注的问题是：把 3DGS 用于几何、表面、深度或网格重建 NeuSG 是宽召回补全中发现的论文，主要根据标题、摘要和 arXiv 类别判定其与 3DGS 的关系。
相对原始 3DGS，NeuSG 的变化不是只替换损失函数，而是把表示、初始化、约束或网络预测路径重新组织到任务流程中：该工作被召回是因为论文文本包含 Gaussian/3D Gaussian 相关表述；当前分类先确定其任务域和与 3DGS 的相关性，后续证据流水线再补代码、企业和引用信息。
\textbf{实现细节：} 具体实现可以按输入、核心模块和渲染/优化路径来读。输入端首先取得任务所需的图像、相机、文本、时间序列、语义或几何线索，并把这些线索转换成可以约束高斯中心、尺度、旋转、不透明度、颜色或附加特征的条件。核心结构中，方法使用的主要机制是：从摘要层面看，NeuSG 围绕“把 3DGS 用于几何、表面、深度或网格重建”展开。若其正文或代码后续可读，应进一步核对它具体改动的是表示、初始化、渲染、优化、损失、动态建模、压缩或应用接口中的哪一环。 训练或推理时，这些模块通常先给出初始化或中间特征，再通过可微 Gaussian rasterization、重建损失、感知/语义约束、几何正则或任务损失回传到高斯参数；如果方法包含前馈网络，则网络输出的是高斯属性、变形场、压缩码、语义特征或控制变量，而不是只输出最终 2D 图像。渲染路径仍以显式高斯投影、深度排序和 alpha 合成为基础，只是在采样、属性预测、局部编辑、动态变形、压缩编码或下游任务接口处加入额外约束。这样写入本章时，它与前后方法的关系是：相对原始 3DGS 和本章已有锚点方法，它可以理解为 3DGS extends\_or\_uses 宽召回分类阶段的关系判断，后续需要按论文方法部分复核。
\textbf{证据与边界：} 本轮召回分类尚未逐项核对实验表格；后续证据补全阶段需要读取论文源码/PDF或项目页确认数据集、指标和对比方法。 该条来自宽召回批次；源码下载受 arXiv e-print 网络限制，部分论文暂时只能依据标题、摘要和类别分类，需要在入可信书前做人工/LLM精读复核。 本书只采用本地 paper pack、数据库分类、源码/引用登记和论文摘要中可核验的信息；没有本地证据支撑的速度、排行榜或开源完整性不在正文中强化。\textbf{方法信息：} 公开链接：\url{https://github.com/cvpr-org/author-kit}，完整性按本地证据保守处理；引用信号：185 citations (5.9/month, domain-relative cutoff 4.3/month, semantic\_scholar)。
\textbf{设计动机与边界：} 相对原始 3DGS 以独立高斯和图像重建误差为主的设定，NeuSG 把问题落在以下可实现的改动上：从摘要层面看，NeuSG 围绕“把 3DGS 用于几何、表面、深度或网格重建”展开。若其正文或代码后续可读，应进一步核对它具体改动的是表示、初始化、渲染、优化、损失、动态建模、压缩或应用接口中的哪一环。其设计目的不是单独增加一个模块，而是用上述表示、监督或渲染约束消除原流程在目标任务中的具体失配。该工作被召回是因为论文文本包含 Gaussian/3D Gaussian 相关表述；当前分类先确定其任务域和与 3DGS 的相关性，后续证据流水线再补代码、企业和引用信息。适用时还应保留论文给出的边界：该条来自宽召回批次；源码下载受 arXiv e-print 网络限制，部分论文暂时只能依据标题、摘要和类别分类，需要在入可信书前做人工/LLM精读复核。
% METHOD: 2401.16416
\paragraph{Endo-4DGS（几何、法线、表面与网格）（2024.01）}
Endo-4DGS~\cite{2401_16416} 关注的问题是：把 3DGS 用于几何、表面、深度或网格重建 Endo-4DGS 是宽召回补全中发现的论文，主要根据标题、摘要和 arXiv 类别判定其与 3DGS 的关系。
相对原始 3DGS，Endo-4DGS 的变化不是只替换损失函数，而是把表示、初始化、约束或网络预测路径重新组织到任务流程中：该工作被召回是因为论文文本包含 Gaussian/3D Gaussian 相关表述；当前分类先确定其任务域和与 3DGS 的相关性，后续证据流水线再补代码、企业和引用信息。
\textbf{实现细节：} 具体实现可以按输入、核心模块和渲染/优化路径来读。输入端首先取得任务所需的图像、相机、文本、时间序列、语义或几何线索，并把这些线索转换成可以约束高斯中心、尺度、旋转、不透明度、颜色或附加特征的条件。核心结构中，方法使用的主要机制是：从摘要层面看，Endo-4DGS 围绕“把 3DGS 用于几何、表面、深度或网格重建”展开。若其正文或代码后续可读，应进一步核对它具体改动的是表示、初始化、渲染、优化、损失、动态建模、压缩或应用接口中的哪一环。 训练或推理时，这些模块通常先给出初始化或中间特征，再通过可微 Gaussian rasterization、重建损失、感知/语义约束、几何正则或任务损失回传到高斯参数；如果方法包含前馈网络，则网络输出的是高斯属性、变形场、压缩码、语义特征或控制变量，而不是只输出最终 2D 图像。渲染路径仍以显式高斯投影、深度排序和 alpha 合成为基础，只是在采样、属性预测、局部编辑、动态变形、压缩编码或下游任务接口处加入额外约束。这样写入本章时，它与前后方法的关系是：相对原始 3DGS 和本章已有锚点方法，它可以理解为 3DGS extends\_or\_uses 宽召回分类阶段的关系判断，后续需要按论文方法部分复核。
\textbf{证据与边界：} 本轮召回分类尚未逐项核对实验表格；后续证据补全阶段需要读取论文源码/PDF或项目页确认数据集、指标和对比方法。 该条来自宽召回批次；源码下载受 arXiv e-print 网络限制，部分论文暂时只能依据标题、摘要和类别分类，需要在入可信书前做人工/LLM精读复核。 本书只采用本地 paper pack、数据库分类、源码/引用登记和论文摘要中可核验的信息；没有本地证据支撑的速度、排行榜或开源完整性不在正文中强化。\textbf{方法信息：} 代码：\url{https://github.com/lastbasket/endo-4dgs}；引用信号：62 citations (2.1/month, domain-relative cutoff 4.3/month, semantic\_scholar)；代码关注：lastbasket/endo-4dgs has 105 stars (3.6/month, domain-relative cutoff 37.1/month)。
\textbf{原始方法补充：} Endo-4DGS 之所以不把内窥镜序列当作静态 3DGS，是因为单目视角、组织形变和器械遮挡会让仅靠颜色重建的高斯漂移。输入由单目图像、Depth-Anything 估计深度和手术器械掩码组成；它把静态高斯 $\mathcal{G}$ 与时间变形 $\Delta\mathcal{G}=\mathcal{F}(\mathcal{G},t)$ 相加，其中两层多分辨率 HexPlane 编码 $x,y,z,t$ 的六个平面，随后用五个小型 MLP 分别预测位置、旋转、尺度、不透明度和 SH 颜色的增量。深度先验用于初始化高斯并约束不可见区域，训练联合最小化颜色、置信度、深度、表面法线和总变差损失；每个时间点先解码变形高斯，再按标准投影、排序和 alpha 混合渲染。为了避免单目内窥镜的颜色重建把器械遮挡和真实组织形变混为一谈，Endo-4DGS 把形变与几何先验放进同一可微管线，而不是让原始 3DGS 以颜色误差独自解释内窥镜运动。
% METHOD: 2403.04926
\paragraph{BAGS（几何、法线、表面与网格）（2024.03）}
BAGS~\cite{2403_04926} 关注的问题是：把 3DGS 用于几何、表面、深度或网格重建 BAGS 是宽召回补全中发现的论文，主要根据标题、摘要和 arXiv 类别判定其与 3DGS 的关系。
相对原始 3DGS，BAGS 的变化不是只替换损失函数，而是把表示、初始化、约束或网络预测路径重新组织到任务流程中：该工作被召回是因为论文文本包含 Gaussian/3D Gaussian 相关表述；当前分类先确定其任务域和与 3DGS 的相关性，后续证据流水线再补代码、企业和引用信息。
\textbf{实现细节：} 具体实现可以按输入、核心模块和渲染/优化路径来读。输入端首先取得任务所需的图像、相机、文本、时间序列、语义或几何线索，并把这些线索转换成可以约束高斯中心、尺度、旋转、不透明度、颜色或附加特征的条件。核心结构中，方法使用的主要机制是：从摘要层面看，BAGS 围绕“把 3DGS 用于几何、表面、深度或网格重建”展开。若其正文或代码后续可读，应进一步核对它具体改动的是表示、初始化、渲染、优化、损失、动态建模、压缩或应用接口中的哪一环。 训练或推理时，这些模块通常先给出初始化或中间特征，再通过可微 Gaussian rasterization、重建损失、感知/语义约束、几何正则或任务损失回传到高斯参数；如果方法包含前馈网络，则网络输出的是高斯属性、变形场、压缩码、语义特征或控制变量，而不是只输出最终 2D 图像。渲染路径仍以显式高斯投影、深度排序和 alpha 合成为基础，只是在采样、属性预测、局部编辑、动态变形、压缩编码或下游任务接口处加入额外约束。这样写入本章时，它与前后方法的关系是：相对原始 3DGS 和本章已有锚点方法，它可以理解为 3DGS extends\_or\_uses 宽召回分类阶段的关系判断，后续需要按论文方法部分复核。
\textbf{证据与边界：} 本轮召回分类尚未逐项核对实验表格；后续证据补全阶段需要读取论文源码/PDF或项目页确认数据集、指标和对比方法。 该条来自宽召回批次；源码下载受 arXiv e-print 网络限制，部分论文暂时只能依据标题、摘要和类别分类，需要在入可信书前做人工/LLM精读复核。 本书只采用本地 paper pack、数据库分类、源码/引用登记和论文摘要中可核验的信息；没有本地证据支撑的速度、排行榜或开源完整性不在正文中强化。\textbf{方法信息：} 代码：\url{https://github.com/snldmt/bags}；引用信号：55 citations (2.0/month, domain-relative cutoff 4.3/month, semantic\_scholar)。
\textbf{设计动机与边界：} 相对原始 3DGS 以独立高斯和图像重建误差为主的设定，BAGS 把问题落在以下可实现的改动上：从摘要层面看，BAGS 围绕“把 3DGS 用于几何、表面、深度或网格重建”展开。若其正文或代码后续可读，应进一步核对它具体改动的是表示、初始化、渲染、优化、损失、动态建模、压缩或应用接口中的哪一环。其设计目的不是单独增加一个模块，而是用上述表示、监督或渲染约束消除原流程在目标任务中的具体失配。该工作被召回是因为论文文本包含 Gaussian/3D Gaussian 相关表述；当前分类先确定其任务域和与 3DGS 的相关性，后续证据流水线再补代码、企业和引用信息。适用时还应保留论文给出的边界：该条来自宽召回批次；源码下载受 arXiv e-print 网络限制，部分论文暂时只能依据标题、摘要和类别分类，需要在入可信书前做人工/LLM精读复核。
% METHOD: 2403.06912
\paragraph{DNGaussian（几何、法线、表面与网格）（2024.03）}
DNGaussian~\cite{2403_06912} 关注的问题是：把 3DGS 用于几何、表面、深度或网格重建 DNGaussian 是宽召回补全中发现的论文，主要根据标题、摘要和 arXiv 类别判定其与 3DGS 的关系。
相对原始 3DGS，DNGaussian 的变化不是只替换损失函数，而是把表示、初始化、约束或网络预测路径重新组织到任务流程中：该工作被召回是因为论文文本包含 Gaussian/3D Gaussian 相关表述；当前分类先确定其任务域和与 3DGS 的相关性，后续证据流水线再补代码、企业和引用信息。
\textbf{实现细节：} 具体实现可以按输入、核心模块和渲染/优化路径来读。输入端首先取得任务所需的图像、相机、文本、时间序列、语义或几何线索，并把这些线索转换成可以约束高斯中心、尺度、旋转、不透明度、颜色或附加特征的条件。核心结构中，方法使用的主要机制是：从摘要层面看，DNGaussian 围绕“把 3DGS 用于几何、表面、深度或网格重建”展开。若其正文或代码后续可读，应进一步核对它具体改动的是表示、初始化、渲染、优化、损失、动态建模、压缩或应用接口中的哪一环。 训练或推理时，这些模块通常先给出初始化或中间特征，再通过可微 Gaussian rasterization、重建损失、感知/语义约束、几何正则或任务损失回传到高斯参数；如果方法包含前馈网络，则网络输出的是高斯属性、变形场、压缩码、语义特征或控制变量，而不是只输出最终 2D 图像。渲染路径仍以显式高斯投影、深度排序和 alpha 合成为基础，只是在采样、属性预测、局部编辑、动态变形、压缩编码或下游任务接口处加入额外约束。这样写入本章时，它与前后方法的关系是：相对原始 3DGS 和本章已有锚点方法，它可以理解为 3DGS extends\_or\_uses 宽召回分类阶段的关系判断，后续需要按论文方法部分复核。
\textbf{证据与边界：} 本轮召回分类尚未逐项核对实验表格；后续证据补全阶段需要读取论文源码/PDF或项目页确认数据集、指标和对比方法。 该条来自宽召回批次；源码下载受 arXiv e-print 网络限制，部分论文暂时只能依据标题、摘要和类别分类，需要在入可信书前做人工/LLM精读复核。 本书只采用本地 paper pack、数据库分类、源码/引用登记和论文摘要中可核验的信息；没有本地证据支撑的速度、排行榜或开源完整性不在正文中强化。\textbf{方法信息：} 代码：\url{https://github.com/graphdeco-inria/gaussian-splatting}；引用信号：325 citations (11.6/month, domain-relative cutoff 4.3/month, semantic\_scholar)；代码关注：fictionarry/dngaussian has 364 stars (13.0/month, domain-relative cutoff 37.1/month)。
\textbf{设计动机与边界：} 相对原始 3DGS 以独立高斯和图像重建误差为主的设定，DNGaussian 把问题落在以下可实现的改动上：从摘要层面看，DNGaussian 围绕“把 3DGS 用于几何、表面、深度或网格重建”展开。若其正文或代码后续可读，应进一步核对它具体改动的是表示、初始化、渲染、优化、损失、动态建模、压缩或应用接口中的哪一环。其设计目的不是单独增加一个模块，而是用上述表示、监督或渲染约束消除原流程在目标任务中的具体失配。该工作被召回是因为论文文本包含 Gaussian/3D Gaussian 相关表述；当前分类先确定其任务域和与 3DGS 的相关性，后续证据流水线再补代码、企业和引用信息。适用时还应保留论文给出的边界：该条来自宽召回批次；源码下载受 arXiv e-print 网络限制，部分论文暂时只能依据标题、摘要和类别分类，需要在入可信书前做人工/LLM精读复核。
% METHOD: 2403.09413
\paragraph{Relaxing Accurate Initialization Constraint for 3D Gaussian Splatting（几何、法线、表面与网格）（2024.03）}
Relaxing Accurate Initialization Constraint for 3D Gaussian Splatting~\cite{2403_09413} 关注的问题是：把 3DGS 用于几何、表面、深度或网格重建 Relaxing Accurate Initialization Constraint for 3D Gaussian Splatting 是宽召回补全中发现的论文，主要根据标题、摘要和 arXiv 类别判定其与 3DGS 的关系。
相对原始 3DGS，Relaxing Accurate Initialization Constraint for 3D Gaussian Splatting 的变化不是只替换损失函数，而是把表示、初始化、约束或网络预测路径重新组织到任务流程中：该工作被召回是因为论文文本包含 Gaussian/3D Gaussian 相关表述；当前分类先确定其任务域和与 3DGS 的相关性，后续证据流水线再补代码、企业和引用信息。
\textbf{实现细节：} 具体实现可以按输入、核心模块和渲染/优化路径来读。输入端首先取得任务所需的图像、相机、文本、时间序列、语义或几何线索，并把这些线索转换成可以约束高斯中心、尺度、旋转、不透明度、颜色或附加特征的条件。核心结构中，方法使用的主要机制是：从摘要层面看，Relaxing Accurate Initialization Constraint for 3D Gaussian Splatting 围绕“把 3DGS 用于几何、表面、深度或网格重建”展开。若其正文或代码后续可读，应进一步核对它具体改动的是表示、初始化、渲染、优化、损失、动态建模、压缩或应用接口中的哪一环。 训练或推理时，这些模块通常先给出初始化或中间特征，再通过可微 Gaussian rasterization、重建损失、感知/语义约束、几何正则或任务损失回传到高斯参数；如果方法包含前馈网络，则网络输出的是高斯属性、变形场、压缩码、语义特征或控制变量，而不是只输出最终 2D 图像。渲染路径仍以显式高斯投影、深度排序和 alpha 合成为基础，只是在采样、属性预测、局部编辑、动态变形、压缩编码或下游任务接口处加入额外约束。这样写入本章时，它与前后方法的关系是：相对原始 3DGS 和本章已有锚点方法，它可以理解为 3DGS extends\_or\_uses 宽召回分类阶段的关系判断，后续需要按论文方法部分复核。
\textbf{证据与边界：} 本轮召回分类尚未逐项核对实验表格；后续证据补全阶段需要读取论文源码/PDF或项目页确认数据集、指标和对比方法。 该条来自宽召回批次；源码下载受 arXiv e-print 网络限制，部分论文暂时只能依据标题、摘要和类别分类，需要在入可信书前做人工/LLM精读复核。 本书只采用本地 paper pack、数据库分类、源码/引用登记和论文摘要中可核验的信息；没有本地证据支撑的速度、排行榜或开源完整性不在正文中强化。\textbf{方法信息：} 代码：\url{https://github.com/cvlab-kaist/rain-gs}；引用信号：57 citations (2.0/month, domain-relative cutoff 4.3/month, semantic\_scholar)；代码关注：ku-cvlab/rain-gs has 363 stars (13.0/month, domain-relative cutoff 37.1/month)。
\textbf{设计动机与边界：} 相对原始 3DGS 以独立高斯和图像重建误差为主的设定，Relaxing Accurate Initialization Constraint for 3D Gaussian Splatting 把问题落在以下可实现的改动上：从摘要层面看，Relaxing Accurate Initialization Constraint for 3D Gaussian Splatting 围绕“把 3DGS 用于几何、表面、深度或网格重建”展开。若其正文或代码后续可读，应进一步核对它具体改动的是表示、初始化、渲染、优化、损失、动态建模、压缩或应用接口中的哪一环。其设计目的不是单独增加一个模块，而是用上述表示、监督或渲染约束消除原流程在目标任务中的具体失配。该工作被召回是因为论文文本包含 Gaussian/3D Gaussian 相关表述；当前分类先确定其任务域和与 3DGS 的相关性，后续证据流水线再补代码、企业和引用信息。适用时还应保留论文给出的边界：该条来自宽召回批次；源码下载受 arXiv e-print 网络限制，部分论文暂时只能依据标题、摘要和类别分类，需要在入可信书前做人工/LLM精读复核。
% METHOD: 2403.09875
\paragraph{Touch-GS（几何、法线、表面与网格）（2024.03）}
Touch-GS~\cite{2403_09875} 关注的问题是：把 3DGS 用于几何、表面、深度或网格重建 Touch-GS 是宽召回补全中发现的论文，主要根据标题、摘要和 arXiv 类别判定其与 3DGS 的关系。
相对原始 3DGS，Touch-GS 的变化不是只替换损失函数，而是把表示、初始化、约束或网络预测路径重新组织到任务流程中：该工作被召回是因为论文文本包含 Gaussian/3D Gaussian 相关表述；当前分类先确定其任务域和与 3DGS 的相关性，后续证据流水线再补代码、企业和引用信息。
\textbf{实现细节：} 具体实现可以按输入、核心模块和渲染/优化路径来读。输入端首先取得任务所需的图像、相机、文本、时间序列、语义或几何线索，并把这些线索转换成可以约束高斯中心、尺度、旋转、不透明度、颜色或附加特征的条件。核心结构中，方法使用的主要机制是：从摘要层面看，Touch-GS 围绕“把 3DGS 用于几何、表面、深度或网格重建”展开。若其正文或代码后续可读，应进一步核对它具体改动的是表示、初始化、渲染、优化、损失、动态建模、压缩或应用接口中的哪一环。 训练或推理时，这些模块通常先给出初始化或中间特征，再通过可微 Gaussian rasterization、重建损失、感知/语义约束、几何正则或任务损失回传到高斯参数；如果方法包含前馈网络，则网络输出的是高斯属性、变形场、压缩码、语义特征或控制变量，而不是只输出最终 2D 图像。渲染路径仍以显式高斯投影、深度排序和 alpha 合成为基础，只是在采样、属性预测、局部编辑、动态变形、压缩编码或下游任务接口处加入额外约束。这样写入本章时，它与前后方法的关系是：相对原始 3DGS 和本章已有锚点方法，它可以理解为 3DGS extends\_or\_uses 宽召回分类阶段的关系判断，后续需要按论文方法部分复核。
\textbf{证据与边界：} 本轮召回分类尚未逐项核对实验表格；后续证据补全阶段需要读取论文源码/PDF或项目页确认数据集、指标和对比方法。 该条来自宽召回批次；源码下载受 arXiv e-print 网络限制，部分论文暂时只能依据标题、摘要和类别分类，需要在入可信书前做人工/LLM精读复核。 本书只采用本地 paper pack、数据库分类、源码/引用登记和论文摘要中可核验的信息；没有本地证据支撑的速度、排行榜或开源完整性不在正文中强化。\textbf{方法信息：} 代码：\url{https://github.com/armlabstanford/touch-gs}。
\textbf{设计动机与边界：} 相对原始 3DGS 以独立高斯和图像重建误差为主的设定，Touch-GS 把问题落在以下可实现的改动上：从摘要层面看，Touch-GS 围绕“把 3DGS 用于几何、表面、深度或网格重建”展开。若其正文或代码后续可读，应进一步核对它具体改动的是表示、初始化、渲染、优化、损失、动态建模、压缩或应用接口中的哪一环。其设计目的不是单独增加一个模块，而是用上述表示、监督或渲染约束消除原流程在目标任务中的具体失配。该工作被召回是因为论文文本包含 Gaussian/3D Gaussian 相关表述；当前分类先确定其任务域和与 3DGS 的相关性，后续证据流水线再补代码、企业和引用信息。适用时还应保留论文给出的边界：该条来自宽召回批次；源码下载受 arXiv e-print 网络限制，部分论文暂时只能依据标题、摘要和类别分类，需要在入可信书前做人工/LLM精读复核。
% METHOD: 2403.10050
\paragraph{Texture-GS（几何、法线、表面与网格）（2024.03）}
Texture-GS~\cite{2403_10050} 关注的问题是：把 3DGS 用于几何、表面、深度或网格重建 Texture-GS 是宽召回补全中发现的论文，主要根据标题、摘要和 arXiv 类别判定其与 3DGS 的关系。
相对原始 3DGS，Texture-GS 的变化不是只替换损失函数，而是把表示、初始化、约束或网络预测路径重新组织到任务流程中：该工作被召回是因为论文文本包含 Gaussian/3D Gaussian 相关表述；当前分类先确定其任务域和与 3DGS 的相关性，后续证据流水线再补代码、企业和引用信息。
\textbf{实现细节：} 具体实现可以按输入、核心模块和渲染/优化路径来读。输入端首先取得任务所需的图像、相机、文本、时间序列、语义或几何线索，并把这些线索转换成可以约束高斯中心、尺度、旋转、不透明度、颜色或附加特征的条件。核心结构中，方法使用的主要机制是：从摘要层面看，Texture-GS 围绕“把 3DGS 用于几何、表面、深度或网格重建”展开。若其正文或代码后续可读，应进一步核对它具体改动的是表示、初始化、渲染、优化、损失、动态建模、压缩或应用接口中的哪一环。 训练或推理时，这些模块通常先给出初始化或中间特征，再通过可微 Gaussian rasterization、重建损失、感知/语义约束、几何正则或任务损失回传到高斯参数；如果方法包含前馈网络，则网络输出的是高斯属性、变形场、压缩码、语义特征或控制变量，而不是只输出最终 2D 图像。渲染路径仍以显式高斯投影、深度排序和 alpha 合成为基础，只是在采样、属性预测、局部编辑、动态变形、压缩编码或下游任务接口处加入额外约束。这样写入本章时，它与前后方法的关系是：相对原始 3DGS 和本章已有锚点方法，它可以理解为 3DGS extends\_or\_uses 宽召回分类阶段的关系判断，后续需要按论文方法部分复核。
\textbf{证据与边界：} 本轮召回分类尚未逐项核对实验表格；后续证据补全阶段需要读取论文源码/PDF或项目页确认数据集、指标和对比方法。 该条来自宽召回批次；源码下载受 arXiv e-print 网络限制，部分论文暂时只能依据标题、摘要和类别分类，需要在入可信书前做人工/LLM精读复核。 本书只采用本地 paper pack、数据库分类、源码/引用登记和论文摘要中可核验的信息；没有本地证据支撑的速度、排行榜或开源完整性不在正文中强化。\textbf{方法信息：} 代码：\url{https://github.com/slothfulxtx/texture-gs}；引用信号：50 citations (1.8/month, domain-relative cutoff 4.3/month, semantic\_scholar)；代码关注：slothfulxtx/texture-gs has 199 stars (7.1/month, domain-relative cutoff 37.1/month)。
\textbf{设计动机与边界：} 相对原始 3DGS 以独立高斯和图像重建误差为主的设定，Texture-GS 把问题落在以下可实现的改动上：从摘要层面看，Texture-GS 围绕“把 3DGS 用于几何、表面、深度或网格重建”展开。若其正文或代码后续可读，应进一步核对它具体改动的是表示、初始化、渲染、优化、损失、动态建模、压缩或应用接口中的哪一环。其设计目的不是单独增加一个模块，而是用上述表示、监督或渲染约束消除原流程在目标任务中的具体失配。该工作被召回是因为论文文本包含 Gaussian/3D Gaussian 相关表述；当前分类先确定其任务域和与 3DGS 的相关性，后续证据流水线再补代码、企业和引用信息。适用时还应保留论文给出的边界：该条来自宽召回批次；源码下载受 arXiv e-print 网络限制，部分论文暂时只能依据标题、摘要和类别分类，需要在入可信书前做人工/LLM精读复核。
% METHOD: 2403.11324
\paragraph{GeoGaussian（几何、法线、表面与网格）（2024.03）}
GeoGaussian~\cite{2403_11324} 关注的问题是：把 3DGS 用于几何、表面、深度或网格重建 GeoGaussian 是宽召回补全中发现的论文，主要根据标题、摘要和 arXiv 类别判定其与 3DGS 的关系。
相对原始 3DGS，GeoGaussian 的变化不是只替换损失函数，而是把表示、初始化、约束或网络预测路径重新组织到任务流程中：该工作被召回是因为论文文本包含 Gaussian/3D Gaussian 相关表述；当前分类先确定其任务域和与 3DGS 的相关性，后续证据流水线再补代码、企业和引用信息。
\textbf{实现细节：} 具体实现可以按输入、核心模块和渲染/优化路径来读。输入端首先取得任务所需的图像、相机、文本、时间序列、语义或几何线索，并把这些线索转换成可以约束高斯中心、尺度、旋转、不透明度、颜色或附加特征的条件。核心结构中，方法使用的主要机制是：从摘要层面看，GeoGaussian 围绕“把 3DGS 用于几何、表面、深度或网格重建”展开。若其正文或代码后续可读，应进一步核对它具体改动的是表示、初始化、渲染、优化、损失、动态建模、压缩或应用接口中的哪一环。 训练或推理时，这些模块通常先给出初始化或中间特征，再通过可微 Gaussian rasterization、重建损失、感知/语义约束、几何正则或任务损失回传到高斯参数；如果方法包含前馈网络，则网络输出的是高斯属性、变形场、压缩码、语义特征或控制变量，而不是只输出最终 2D 图像。渲染路径仍以显式高斯投影、深度排序和 alpha 合成为基础，只是在采样、属性预测、局部编辑、动态变形、压缩编码或下游任务接口处加入额外约束。这样写入本章时，它与前后方法的关系是：相对原始 3DGS 和本章已有锚点方法，它可以理解为 3DGS extends\_or\_uses 宽召回分类阶段的关系判断，后续需要按论文方法部分复核。
\textbf{证据与边界：} 本轮召回分类尚未逐项核对实验表格；后续证据补全阶段需要读取论文源码/PDF或项目页确认数据集、指标和对比方法。 该条来自宽召回批次；源码下载受 arXiv e-print 网络限制，部分论文暂时只能依据标题、摘要和类别分类，需要在入可信书前做人工/LLM精读复核。 本书只采用本地 paper pack、数据库分类、源码/引用登记和论文摘要中可核验的信息；没有本地证据支撑的速度、排行榜或开源完整性不在正文中强化。\textbf{方法信息：} 代码：\url{https://github.com/yanyan-li/GeoGaussian}；引用信号：67 citations (2.4/month, domain-relative cutoff 4.3/month, semantic\_scholar)；代码关注：yanyan-li/geogaussian has 225 stars (8.1/month, domain-relative cutoff 37.1/month)。
\textbf{设计动机与边界：} 相对原始 3DGS 以独立高斯和图像重建误差为主的设定，GeoGaussian 把问题落在以下可实现的改动上：从摘要层面看，GeoGaussian 围绕“把 3DGS 用于几何、表面、深度或网格重建”展开。若其正文或代码后续可读，应进一步核对它具体改动的是表示、初始化、渲染、优化、损失、动态建模、压缩或应用接口中的哪一环。其设计目的不是单独增加一个模块，而是用上述表示、监督或渲染约束消除原流程在目标任务中的具体失配。该工作被召回是因为论文文本包含 Gaussian/3D Gaussian 相关表述；当前分类先确定其任务域和与 3DGS 的相关性，后续证据流水线再补代码、企业和引用信息。适用时还应保留论文给出的边界：该条来自宽召回批次；源码下载受 arXiv e-print 网络限制，部分论文暂时只能依据标题、摘要和类别分类，需要在入可信书前做人工/LLM精读复核。
% METHOD: 2404.17774
\paragraph{High-quality Surface Reconstruction using Gaussian Surfels（几何、法线、表面与网格）（2024.04）}
High-quality Surface Reconstruction using Gaussian Surfels~\cite{2404_17774} 关注的问题是：把 3DGS 用于几何、表面、深度或网格重建 High-quality Surface Reconstruction using Gaussian Surfels 是宽召回补全中发现的论文，主要根据标题、摘要和 arXiv 类别判定其与 3DGS 的关系。
相对原始 3DGS，High-quality Surface Reconstruction using Gaussian Surfels 的变化不是只替换损失函数，而是把表示、初始化、约束或网络预测路径重新组织到任务流程中：该工作被召回是因为论文文本包含 Gaussian/3D Gaussian 相关表述；当前分类先确定其任务域和与 3DGS 的相关性，后续证据流水线再补代码、企业和引用信息。
\textbf{实现细节：} 具体实现可以按输入、核心模块和渲染/优化路径来读。输入端首先取得任务所需的图像、相机、文本、时间序列、语义或几何线索，并把这些线索转换成可以约束高斯中心、尺度、旋转、不透明度、颜色或附加特征的条件。核心结构中，方法使用的主要机制是：从摘要层面看，High-quality Surface Reconstruction using Gaussian Surfels 围绕“把 3DGS 用于几何、表面、深度或网格重建”展开。若其正文或代码后续可读，应进一步核对它具体改动的是表示、初始化、渲染、优化、损失、动态建模、压缩或应用接口中的哪一环。 训练或推理时，这些模块通常先给出初始化或中间特征，再通过可微 Gaussian rasterization、重建损失、感知/语义约束、几何正则或任务损失回传到高斯参数；如果方法包含前馈网络，则网络输出的是高斯属性、变形场、压缩码、语义特征或控制变量，而不是只输出最终 2D 图像。渲染路径仍以显式高斯投影、深度排序和 alpha 合成为基础，只是在采样、属性预测、局部编辑、动态变形、压缩编码或下游任务接口处加入额外约束。这样写入本章时，它与前后方法的关系是：相对原始 3DGS 和本章已有锚点方法，它可以理解为 3DGS extends\_or\_uses 宽召回分类阶段的关系判断，后续需要按论文方法部分复核。
\textbf{证据与边界：} 本轮召回分类尚未逐项核对实验表格；后续证据补全阶段需要读取论文源码/PDF或项目页确认数据集、指标和对比方法。 该条来自宽召回批次；源码下载受 arXiv e-print 网络限制，部分论文暂时只能依据标题、摘要和类别分类，需要在入可信书前做人工/LLM精读复核。 本书只采用本地 paper pack、数据库分类、源码/引用登记和论文摘要中可核验的信息；没有本地证据支撑的速度、排行榜或开源完整性不在正文中强化。\textbf{方法信息：} 代码：\url{https://github.com/turandai/gaussian\_surfels}；引用信号：237 citations (9.0/month, domain-relative cutoff 4.3/month, semantic\_scholar)；代码关注：turandai/gaussian\_surfels has 680 stars (25.7/month, domain-relative cutoff 37.1/month)。
\textbf{设计动机与边界：} 相对原始 3DGS 以独立高斯和图像重建误差为主的设定，High-quality Surface Reconstruction using Gaussian Surfels 把问题落在以下可实现的改动上：从摘要层面看，High-quality Surface Reconstruction using Gaussian Surfels 围绕“把 3DGS 用于几何、表面、深度或网格重建”展开。若其正文或代码后续可读，应进一步核对它具体改动的是表示、初始化、渲染、优化、损失、动态建模、压缩或应用接口中的哪一环。其设计目的不是单独增加一个模块，而是用上述表示、监督或渲染约束消除原流程在目标任务中的具体失配。该工作被召回是因为论文文本包含 Gaussian/3D Gaussian 相关表述；当前分类先确定其任务域和与 3DGS 的相关性，后续证据流水线再补代码、企业和引用信息。适用时还应保留论文给出的边界：该条来自宽召回批次；源码下载受 arXiv e-print 网络限制，部分论文暂时只能依据标题、摘要和类别分类，需要在入可信书前做人工/LLM精读复核。
% METHOD: 2406.06521
\paragraph{PGSR（平面高斯与无偏深度）（2024.06）}
PGSR~\cite{2406_06521} 是平面/多视图几何路线的锚点。它提出无偏深度渲染：直接渲染从相机原点到高斯平面的距离和法线图，再通过比值得到深度，而不是使用高斯中心深度。

PGSR 的关键不是简单给 2DGS 增加损失，而是重新定义平面高斯的几何读数。对每个平面高斯，它渲染相机原点到该高斯平面的射线距离，同时渲染法线图，再用两者关系恢复像素深度；这避免了“把高斯中心深度加权平均”造成的系统偏差。训练中有三类约束：单视图几何正则让同一视图的深度、法线和局部表面一致；多视图光度正则把参考视图和邻近视图按可见性重投影后比较颜色；多视图几何正则进一步约束跨视图深度和法线闭合。方法还为相机学习曝光补偿，降低不同光照条件对几何监督的干扰。网格输出时，PGSR 过滤不可靠深度并用 TSDF fusion 得到表面。相对 2DGS 的通用 surfel 几何，PGSR 更强调无偏深度和多视图一致；后续 TriaGS、VA-GS、2DGS-R 等方法都在不同程度上沿着“深度定义与几何监督”这条线继续改。
\textbf{实现细节：} PGSR（2024.06） 修改了原始 3DGS 的深度渲染与正则化环节，用平面化高斯的无偏深度渲染替代标准 $\alpha$-混合 深度，光度损失和自适应密度控制保持不变。PGSR 是多视角几何约束的里程碑工作。动机。标准 $\alpha$-混合期望深度在高斯不完整累积时系统性偏向相机侧，且椭球形高斯的"平均深度"不对应任何物理表面，导致渲染深度图无法直接用于表面重建。问题：标准深度渲染的偏差。直觉地说，标准 3DGS 的期望深度 $D = \sum_{i \in N} T_i \alpha_i z_i$ 就像对一条射线上的所有高斯做加权平均——但这个平均有两个问题：一是每个高斯本身是椭球而非平面，"平均"出来的深度对应一个弯曲面；二是累积权重 $T_i \alpha_i$ 的总和可能小于 1（射线末端的空白区域不贡献任何高斯），导致深度整体偏向相机侧。这两个偏差使得标准深度图无法直接用于表面重建。方案：平面化与无偏深度。为消除上述偏差，PGSR 先将三维高斯椭球沿最小缩放轴压扁为平面，最小化最短缩放因子 $\mathcal{L}_s = \|\min(s_1, s_2, s_3)\|_1$，使每个高斯局部逼近一个平面片。然后分别渲染法线图和相机到平面距离图： 其中 $R_c$ 是相机到世界坐标系的旋转，$d_i = (R_c^T(\bm{\mu}_i - T_c))^T(R_c^T \bm{n}_i)$ 是第 $i$ 个高斯平面到相机中心的距离。关键一步：通过射线与平面的交点关系将距离和法线转换为无偏深度 $D(\bm{p}) = \mathcal{D} / (\bm{N}(\bm{p}) K^{-1} \tilde{\bm{p}})$，其中 $\tilde{\bm{p}}$ 是像素 $[u,v]^T$ 的齐次坐标。直觉上，这一除法相当于"用法线方向修正距离值"，消除了累积权重不等于 1 带来的偏差，使深度值真正落在高斯平面上而非偏向相机。三层正则化。基于无偏深度图，PGSR 引入由局部到全局的三层约束。
\textbf{锚点细化：} PGSR：针对 $\alpha$-混合 期望深度在累积权重小于 1 时系统性偏向相机侧的局限，用"平面化高斯的无偏深度渲染"替代期望深度，核心思想是先将椭球沿最小缩放轴压扁为平面，再分别渲染法线图和距离图，通过射线-平面交点关系将距离除以法线点积得到无偏深度——这一除法消除了累积权重偏差使深度真正落在高斯平面上，是 3DGS 多视角几何约束的里程碑工作。
\textbf{方法信息：} 代码：\url{https://github.com/zju3dv/PGSR}；GitHub 关注度很高，约 1058 stars；引用很多，Semantic Scholar 约 347 次。



% METHOD: 2412.01931
\paragraph{Planar Gaussian Splatting（解析平面实例）（2024.12）}
Planar Gaussian Splatting~\cite{2412_01931} 直接从多视角 RGB 学习 3D 平面。每个高斯除了常规位置、协方差、颜色和透明度，还带平面描述符和表面法线。2D SAM 掩码被提升到 3D，为平面描述符提供弱语义线索。

方法先用 2D 分割基础模型在训练图像中产生候选平面/物体 mask，再把这些 mask 的一致性提升到三维高斯上，学习每个高斯的 plane descriptor。由于不同视图中的 mask 数量和身份不固定，论文把跨视图 proposal 关联转成线性回归和图划分问题，使属于同一物理平面的高斯获得相近描述符。随后，它构建层次 Gaussian mixture tree：从叶子高斯开始，按位置、法线和描述符相似性概率性合并，逐步得到不同 3D 平面实例。训练仍通过可微渲染优化 RGB 和几何属性，不需要 3D 平面标签或深度监督。相对 PGSR，它不只是用平面改善深度，而是输出解析平面结构，适合室内布局理解和机器人/AR 场景解析。
\textbf{实现细节：} PGS（2024.12） 提出完全无监督的 3D 平面重建方法，核心技术是高斯混合树（Gaussian Mixture Tree, GMT）和平面描述子。PGS 为每个高斯额外学习法线 $\mathbf{n} \in \mathbb{R}^{}3$ 和平面描述子 $\mathbf{z} \in \mathbb{R}^k$（$\|\mathbf{z}\|_2 = 1$）。描述子通过将 SAM 2D 分割掩码提升到 3D 来学习：对每张训练图，用线性回归将渲染的描述子图映射到 SAM 的 one-hot 分割标签，闭式求解 $W = (Z^{}\top Z)^{-1} Z^{}\top Y$，无需跨视角的掩码关联。GMT 自底向上构建：先用 KNN（K-Nearest Neighbors，K 近邻）聚类局部高斯作为叶节点，再用 Bhattacharyya 距离（衡量两个概率分布相似度的度量，值越小越相似）$D_B$ 和描述子余弦相似度作为合并准则，递归地将相邻高斯节点合并为父节点，形成树结构——根节点的子节点即为各独立平面实例。此外，PGS 通过局部平面对齐将高斯中心投影到 KNN 邻域的局部切平面（由协方差矩阵最大两个特征值对应的特征向量定义），迫使高斯精确贴合物体表面。训练中周期性施加 recurrent mean-shift 更新 $\mathbf{Z} \leftarrow \mathbf{Z}(\alpha K D^{-1} + (1-\alpha) I)$（基于 von Mises-Fisher 分布），确保所有平面描述子在整个场景中保持全局可分性。
\textbf{方法信息：} Qualcomm AI Research 参与了这篇论文；当前未确认完整公开代码。
\textbf{设计动机与边界：} 相对原始 3DGS 以独立高斯和图像重建误差为主的设定，Planar Gaussian Splatting 把问题落在以下可实现的改动上：PGS方法的核心是将场景建模为高斯基元集合，每个基元除了常规的位置、协方差、颜色、透明度外，还包含平面描述符和表面法线。方法采用层次高斯混合的树状结构，从叶子节点开始概率性地将相似高斯合并，逐步识别出不同的3D平面实例。平面描述符通过2D分割模型（如SAM）的掩码提升到3D获得。整个框架通过可微渲染进行端到端优化，同时实现新视角合成和平面重建。关键优势在于不依赖任何3D平面标注或深度监督，完全通过多视角图像自监督学习。其设计目的不是单独增加一个模块，而是用上述表示、监督或渲染约束消除原流程在目标任务中的具体失配。提出PGS，首个利用高斯泼溅基元直接从多视角RGB图像学习3D几何和解析3D平面的方法。通过层次高斯混合将相似高斯概率性合并以识别不同3D平面实例，无需3D平面标签或深度监督。相比现有监督方法具有更好的跨数据集泛化能力，且比优化方法快得多。适用时还应保留论文给出的边界：主要针对室内场景中的平面结构；对曲面和非平面几何的处理能力有限；高度依赖2D分割质量。
% METHOD: 2503.04079
\paragraph{Surgical Gaussian Surfels（几何、法线、表面与网格）（2025.03）}
Surgical Gaussian Surfels~\cite{2503_04079} 关注的问题是：把 3DGS 用于几何、表面、深度或网格重建 Surgical Gaussian Surfels 是宽召回补全中发现的论文，主要根据标题、摘要和 arXiv 类别判定其与 3DGS 的关系。
相对原始 3DGS，Surgical Gaussian Surfels 的变化不是只替换损失函数，而是把表示、初始化、约束或网络预测路径重新组织到任务流程中：该工作被召回是因为论文文本包含 Gaussian/3D Gaussian 相关表述；当前分类先确定其任务域和与 3DGS 的相关性，后续证据流水线再补代码、企业和引用信息。
\textbf{实现细节：} 具体实现可以按输入、核心模块和渲染/优化路径来读。输入端首先取得任务所需的图像、相机、文本、时间序列、语义或几何线索，并把这些线索转换成可以约束高斯中心、尺度、旋转、不透明度、颜色或附加特征的条件。核心结构中，方法使用的主要机制是：从摘要层面看，Surgical Gaussian Surfels 围绕“把 3DGS 用于几何、表面、深度或网格重建”展开。若其正文或代码后续可读，应进一步核对它具体改动的是表示、初始化、渲染、优化、损失、动态建模、压缩或应用接口中的哪一环。 训练或推理时，这些模块通常先给出初始化或中间特征，再通过可微 Gaussian rasterization、重建损失、感知/语义约束、几何正则或任务损失回传到高斯参数；如果方法包含前馈网络，则网络输出的是高斯属性、变形场、压缩码、语义特征或控制变量，而不是只输出最终 2D 图像。渲染路径仍以显式高斯投影、深度排序和 alpha 合成为基础，只是在采样、属性预测、局部编辑、动态变形、压缩编码或下游任务接口处加入额外约束。这样写入本章时，它与前后方法的关系是：相对原始 3DGS 和本章已有锚点方法，它可以理解为 3DGS extends\_or\_uses 宽召回分类阶段的关系判断，后续需要按论文方法部分复核。
\textbf{证据与边界：} 本轮召回分类尚未逐项核对实验表格；后续证据补全阶段需要读取论文源码/PDF或项目页确认数据集、指标和对比方法。 该条来自宽召回批次；源码下载受 arXiv e-print 网络限制，部分论文暂时只能依据标题、摘要和类别分类，需要在入可信书前做人工/LLM精读复核。 本书只采用本地 paper pack、数据库分类、源码/引用登记和论文摘要中可核验的信息；没有本地证据支撑的速度、排行榜或开源完整性不在正文中强化。\textbf{方法信息：} 代码：\url{https://github.com/aloma85/surgicalgaussiansurfels}。
\textbf{设计动机与边界：} 相对原始 3DGS 以独立高斯和图像重建误差为主的设定，Surgical Gaussian Surfels 把问题落在以下可实现的改动上：从摘要层面看，Surgical Gaussian Surfels 围绕“把 3DGS 用于几何、表面、深度或网格重建”展开。若其正文或代码后续可读，应进一步核对它具体改动的是表示、初始化、渲染、优化、损失、动态建模、压缩或应用接口中的哪一环。其设计目的不是单独增加一个模块，而是用上述表示、监督或渲染约束消除原流程在目标任务中的具体失配。该工作被召回是因为论文文本包含 Gaussian/3D Gaussian 相关表述；当前分类先确定其任务域和与 3DGS 的相关性，后续证据流水线再补代码、企业和引用信息。适用时还应保留论文给出的边界：该条来自宽召回批次；源码下载受 arXiv e-print 网络限制，部分论文暂时只能依据标题、摘要和类别分类，需要在入可信书前做人工/LLM精读复核。
% METHOD: 2511.03950
\paragraph{Improving Multi-View Reconstruction via Texture-Guided Gaussian-Mesh Joint Optimization（几何、法线、表面与网格）（2025.11）}
Improving Multi-View Reconstruction via Texture-Guided Gaussian-Mesh Joint Optimization~\cite{2511_03950} 关注的问题是：把 3DGS 用于几何、表面、深度或网格重建 Improving Multi-View Reconstruction via Texture-Guided Gaussian-Mesh Joint Optim 是宽召回补全中发现的论文，主要根据标题、摘要和 arXiv 类别判定其与 3DGS 的关系。
作为 supplementary 类条目，Improving Multi-View Reconstruction via Texture-Guided Gaussian-Mesh Joint Optimization 不应被写成单纯的新渲染算子。它更像是把 3DGS 放进 几何与表面重建 的协议、系统或应用边界里：该工作被召回是因为论文文本包含 Gaussian/3D Gaussian 相关表述；当前分类先确定其任务域和与 3DGS 的相关性，后续证据流水线再补代码、企业和引用信息。
\textbf{实现细节：} 具体实现可以按输入、核心模块和渲染/优化路径来读。输入端首先取得任务所需的图像、相机、文本、时间序列、语义或几何线索，并把这些线索转换成可以约束高斯中心、尺度、旋转、不透明度、颜色或附加特征的条件。核心结构中，方法使用的主要机制是：从摘要层面看，Improving Multi-View Reconstruction via Texture-Guided Gaussian-Mesh Joint Optim 围绕“把 3DGS 用于几何、表面、深度或网格重建”展开。若其正文或代码后续可读，应进一步核对它具体改动的是表示、初始化、渲染、优化、损失、动态建模、压缩或应用接口中的哪一环。 训练或推理时，这些模块通常先给出初始化或中间特征，再通过可微 Gaussian rasterization、重建损失、感知/语义约束、几何正则或任务损失回传到高斯参数；如果方法包含前馈网络，则网络输出的是高斯属性、变形场、压缩码、语义特征或控制变量，而不是只输出最终 2D 图像。渲染路径仍以显式高斯投影、深度排序和 alpha 合成为基础，只是在采样、属性预测、局部编辑、动态变形、压缩编码或下游任务接口处加入额外约束。这样写入本章时，它与前后方法的关系是：相对原始 3DGS 和本章已有锚点方法，它可以理解为 3DGS extends\_or\_uses 宽召回分类阶段的关系判断，后续需要按论文方法部分复核。
\textbf{证据与边界：} 本轮召回分类尚未逐项核对实验表格；后续证据补全阶段需要读取论文源码/PDF或项目页确认数据集、指标和对比方法。 该条来自宽召回批次；源码下载受 arXiv e-print 网络限制，部分论文暂时只能依据标题、摘要和类别分类，需要在入可信书前做人工/LLM精读复核。 本书只采用本地 paper pack、数据库分类、源码/引用登记和论文摘要中可核验的信息；没有本地证据支撑的速度、排行榜或开源完整性不在正文中强化。\textbf{方法信息：} 代码：\url{https://github.com/zhejia01/texguided-gs2mesh}。
\textbf{设计动机与边界：} 相对原始 3DGS 以独立高斯和图像重建误差为主的设定，Improving Multi-View Reconstruction via Texture-Guided Gaussian-Mesh Joint Optimization 把问题落在以下可实现的改动上：从摘要层面看，Improving Multi-View Reconstruction via Texture-Guided Gaussian-Mesh Joint Optim 围绕“把 3DGS 用于几何、表面、深度或网格重建”展开。若其正文或代码后续可读，应进一步核对它具体改动的是表示、初始化、渲染、优化、损失、动态建模、压缩或应用接口中的哪一环。其设计目的不是单独增加一个模块，而是用上述表示、监督或渲染约束消除原流程在目标任务中的具体失配。该工作被召回是因为论文文本包含 Gaussian/3D Gaussian 相关表述；当前分类先确定其任务域和与 3DGS 的相关性，后续证据流水线再补代码、企业和引用信息。适用时还应保留论文给出的边界：该条来自宽召回批次；源码下载受 arXiv e-print 网络限制，部分论文暂时只能依据标题、摘要和类别分类，需要在入可信书前做人工/LLM精读复核。
% METHOD: 2512.06269
\paragraph{TriaGS（多视角三角测量一致性）（2025.12）}
TriaGS~\cite{2512_06269} 把经典多视角三角测量引入 3DGS 优化。PGSR、GaussianSurf 这类方法多在成对视图之间做深度或法线一致性，约束是局部的；一旦某个视图深度有偏，误差可能沿视图对累积。TriaGS 的改变是为每个渲染出的 3D 点，从多个邻近视图重新计算一个三角测量共识点。

流程先需要可靠的 3D 点。为此它借鉴 RaDe-GS 的几何一致深度，不用普通 alpha 混合深度，而是考虑射线与 3D 高斯椭球的交点，在局部仿射近似下得到更接近真实表面的深度。然后，对一个参考视图中的渲染点，方法在多个相邻视图中寻找对应观测并做三角测量，得到共识 3D 点；优化时最小化当前高斯渲染点到该共识点的距离。多视角三角测量天然具备离群点拒绝和歧义消解能力，因此约束比成对一致性更全局。相对 PGSR 的平面深度一致，TriaGS 把几何监督推进到多视角共识三维点。
\textbf{实现细节：} TriaGS（2025.12） 将经典多视角三角测量（triangulation，通过多条射线的交汇确定三维点位置）以可微形式融入 3DGS 训练循环，用全局自一致性信号替代 PGSR 等方法的成对视角比较。其核心洞察是：成对比较会累积误差导致全局几何漂移，而多视角三角测量让所有视角同时参与对一个三维点的共识投票，从根本上消除漂移。几何一致的深度-法线渲染。标准 3DGS 的 $\alpha$-混合 深度忽略了高斯的形状和朝向，斜盘状高斯的深度值可能不对应任何真实表面点，产生噪声深度和浮游伪影。TriaGS 基于 RaDe-GS 的光栅化深度公式，将射线与高斯椭球的交点定义为沿射线密度最大处的距离——直觉上，这等价于"射线穿过高斯最稠密的点"，即高斯中心在射线方向上的投影。
\textbf{方法信息：} 代码：\url{https://github.com/cogniboticslab/triags}。
\textbf{设计动机与边界：} 相对原始 3DGS 以独立高斯和图像重建误差为主的设定，TriaGS 把问题落在以下可实现的改动上：TriaGS框架核心是将经典多视角三角测量融入3DGS优化管道：（1）问题分析：3DGS仅由光度损失引导导致几何不一致，现有成对视图一致性方法（如GaussianSurf、PGSR）是严格局部的，容易误差累积和全局几何漂移。（2）基于RaDe-GS的几何一致深度渲染：标准3DGS的alpha混合深度忽略高斯形状和朝向，产生噪声深度图。TriaGS利用RaDe-GS方法确定视线与3D高斯椭球的精确交点（最大密度位置），通过局部仿射近似将计算转换到射线空间实现可光栅化。（3）多视角三角测量一致性：对于每个渲染的3D点，从一组邻近视角重新三角测量其位置获得鲁棒共识点，最小化初始估计与共识点的距离作为可微损失。这实现了超越成对约束的全局几何一致性，利用多视角同时聚合约束实现离群点拒绝和歧义消解。其设计目的不是单独增加一个模块，而是用上述表示、监督或渲染约束消除原流程在目标任务中的具体失配。提出首个基于多视角三角测量的全局几何一致性约束，超越成对视图局部一致性限制，在DTU上达到0.50mm Chamfer Distance，超过可比显式方法，完全自监督无需额外标注。适用时还应保留论文给出的边界：多视角三角测量增加计算开销；需要足够的邻近视角覆盖；对纹理缺失区域可能仍有限制。
% METHOD: 2601.20331
\paragraph{GVGS（高斯可见性感知多视图几何）（2026.01）}
GVGS~\cite{2601_20331} 认为多视图深度重投影存在循环依赖：深度差时，共可见区域判断也会错。它从高斯基元级别显式建模可见性，用渲染贡献估计每个高斯在不同视图中的可见程度。

GVMV 首先在邻近视图渲染时累计每个高斯对所有像素的 $\alpha$-混合贡献，把这个贡献归一化为该高斯在该视图中的可见概率；同一高斯在参考视图和邻近视图中都可见，才进入共可见集合。几何一致性损失随后只在这些高斯支撑的区域上比较深度和法线，因此可见性不再由可能错误的重投影深度决定。单目深度先验方面，QDC 不直接把 Depth Anything 一类深度图按全局尺度对齐，而是在可见性感知区域内做粗到细四叉树分块仿射校准：大块先校准整体尺度，小块再保留局部结构。相对 TriaGS 的多视角三角测量，GVGS 的锚点价值在于把“哪些约束可信”提前到高斯可见性层面，再用 GVMV 和 QDC 共同监督表面。
\textbf{实现细节：} GVGS（2026.01） 修改了 PGSR 的可见性判断机制——用高斯基元级可见性替代像素级深度重投影，多视角几何约束的框架保持不变。此前 PGSR 的多视角几何一致性损失依赖前-后向重投影误差 $\phi(\bm{p}_r)$ 估计置信度，$\phi \geq 1$ 的像素被视为遮挡而排除——但这恰恰在纹理匮乏区域（如建筑立面）造成大面积监督空白，形成"深度不可靠导致不可见、不可见导致无监督"的恶性循环。GVGS 的核心思路是将可见性建模从像素级别的深度重投影切换到高斯基元级别，利用体积渲染合成公式本身推导可见性，从而解耦可见性与深度质量。高斯可见性感知的多视角一致性（GVMV）。在差分光栅化过程中，GVGS 直接计算每个高斯在邻近视角 $v_n$ 中的渲染贡献 $W_i = \sum_{x \in \Omega_n} \alpha_i(x) T_i(x)$，将其归一化为可见性概率 $p_i = W_i / \sum_j W_j$，实际采用阈值化近似 $\hat{v}_i = \mathbb{I}(W_i > \epsilon)$ 以兼顾效率。直觉上，这相当于问"这个高斯在另一个视角中是否被看到"，而非"这个像素的深度在另一视角中是否一致"。然后将可见性投影回参考视角，构建选择性累积不透明度图 $O_r(x) = \sum_i \hat{v}_i \cdot \alpha_i(x) \prod_{j<i}(1-\alpha_j(x))$，该图仅聚合在两个视角中均确认可见的高斯贡献。在此基础上，GVGS 扩展 PGSR 的多视角几何一致性损失为 $\mathcal{L}_{\mathrm{gvmvgeom}} = \frac{1}{|V|}\sum_{x \in V} (\exp(-\phi(x)) + \beta \cdot O_r(x)) \cdot \phi(x)$，其中 $V$ 既包含传统深度一致性像素，也包含通过可见性图新发现的共视区域像素，$\beta$ 动态平衡两项贡献。四叉树校准深度约束（QDC）。GVGS 将单目深度先验（来自 DepthAnything 等大模型）作为粗结构引导，但全局尺度-偏移校准无法捕捉空间变化的深度畸变。
\textbf{方法信息：} 代码：\url{https://github.com/gvgscode/gvgs}。
\textbf{设计动机与边界：} 相对原始 3DGS 以独立高斯和图像重建误差为主的设定，GVGS 把问题落在以下可实现的改动上：GVVS包含两个核心组件：(1) 高斯可见性感知多视图一致性(GVMV)：基于高斯渲染贡献估计每个高斯的跨视图可见性，定义共可见高斯集合，构建可靠的几何监督信号。具体包括：高斯级可见性估计模块（通过体积合成推导可见性，解耦与深度质量的依赖）、基于可见性的深度一致性损失、基于可见性的法线一致性损失。(2) 渐进式四叉树校准深度约束(QDC)：通过粗到细的分块仿射校准在可见性感知指导下整合单目深度先验，缓解尺度模糊性同时保留局部几何保真度。整体框架通过联合优化GVMV多视图约束和QDC单目约束实现高质量表面重建。其设计目的不是单独增加一个模块，而是用上述表示、监督或渲染约束消除原流程在目标任务中的具体失配。提出GVVS，从高斯基元级别显式建模可见性，打破深度重投影中的循环依赖。引入高斯可见性感知多视图一致性(GVMV)公式，通过聚合共享高斯的跨视图可见性构建可靠的共可见区域监督；提出渐进式四叉树校准深度对齐(QDC)策略，在可见性感知指导下执行分块仿射校准，有效缓解尺度模糊性同时保留局部几何结构。适用时还应保留论文给出的边界：依赖多视图校准相机；在弱纹理区域仍受限于单目深度先验的质量；四叉树分块策略的超参数需调优；计算开销较标准3DGS增加。
% METHOD: 2601.06997
\paragraph{ObjSplat（几何、法线、表面与网格）（2026.01）}
ObjSplat~\cite{2601_06997} 关注的问题是：把 3DGS 用于几何、表面、深度或网格重建 ObjSplat 是宽召回补全中发现的论文，主要根据标题、摘要和 arXiv 类别判定其与 3DGS 的关系。
作为 supplementary 类条目，ObjSplat 不应被写成单纯的新渲染算子。它更像是把 3DGS 放进 几何与表面重建 的协议、系统或应用边界里：该工作被召回是因为论文文本包含 Gaussian/3D Gaussian 相关表述；当前分类先确定其任务域和与 3DGS 的相关性，后续证据流水线再补代码、企业和引用信息。
\textbf{实现细节：} 具体实现可以按输入、核心模块和渲染/优化路径来读。输入端首先取得任务所需的图像、相机、文本、时间序列、语义或几何线索，并把这些线索转换成可以约束高斯中心、尺度、旋转、不透明度、颜色或附加特征的条件。核心结构中，方法使用的主要机制是：从摘要层面看，ObjSplat 围绕“把 3DGS 用于几何、表面、深度或网格重建”展开。若其正文或代码后续可读，应进一步核对它具体改动的是表示、初始化、渲染、优化、损失、动态建模、压缩或应用接口中的哪一环。 训练或推理时，这些模块通常先给出初始化或中间特征，再通过可微 Gaussian rasterization、重建损失、感知/语义约束、几何正则或任务损失回传到高斯参数；如果方法包含前馈网络，则网络输出的是高斯属性、变形场、压缩码、语义特征或控制变量，而不是只输出最终 2D 图像。渲染路径仍以显式高斯投影、深度排序和 alpha 合成为基础，只是在采样、属性预测、局部编辑、动态变形、压缩编码或下游任务接口处加入额外约束。这样写入本章时，它与前后方法的关系是：相对原始 3DGS 和本章已有锚点方法，它可以理解为 3DGS extends\_or\_uses 宽召回分类阶段的关系判断，后续需要按论文方法部分复核。
\textbf{证据与边界：} 本轮召回分类尚未逐项核对实验表格；后续证据补全阶段需要读取论文源码/PDF或项目页确认数据集、指标和对比方法。 该条来自宽召回批次；源码下载受 arXiv e-print 网络限制，部分论文暂时只能依据标题、摘要和类别分类，需要在入可信书前做人工/LLM精读复核。 本书只采用本地 paper pack、数据库分类、源码/引用登记和论文摘要中可核验的信息；没有本地证据支撑的速度、排行榜或开源完整性不在正文中强化。\textbf{方法信息：} 代码：\url{https://github.com/zabayone/objsplat}。
\textbf{设计动机与边界：} 相对原始 3DGS 以独立高斯和图像重建误差为主的设定，ObjSplat 把问题落在以下可实现的改动上：从摘要层面看，ObjSplat 围绕“把 3DGS 用于几何、表面、深度或网格重建”展开。若其正文或代码后续可读，应进一步核对它具体改动的是表示、初始化、渲染、优化、损失、动态建模、压缩或应用接口中的哪一环。其设计目的不是单独增加一个模块，而是用上述表示、监督或渲染约束消除原流程在目标任务中的具体失配。该工作被召回是因为论文文本包含 Gaussian/3D Gaussian 相关表述；当前分类先确定其任务域和与 3DGS 的相关性，后续证据流水线再补代码、企业和引用信息。适用时还应保留论文给出的边界：该条来自宽召回批次；源码下载受 arXiv e-print 网络限制，部分论文暂时只能依据标题、摘要和类别分类，需要在入可信书前做人工/LLM精读复核。
% METHOD: 2602.21552
\paragraph{Generalizing Visual Geometry Priors to Sparse Gaussian Occupancy Prediction（几何、法线、表面与网格）（2026.02）}
Generalizing Visual Geometry Priors to Sparse Gaussian Occupancy Prediction~\cite{2602_21552} 关注的问题是：把 3DGS 用于几何、表面、深度或网格重建 Generalizing Visual Geometry Priors to Sparse Gaussian Occupancy Prediction 是宽召回补全中发现的论文，主要根据标题、摘要和 arXiv 类别判定其与 3DGS 的关系。
作为 supplementary 类条目，Generalizing Visual Geometry Priors to Sparse Gaussian Occupancy Prediction 不应被写成单纯的新渲染算子。它更像是把 3DGS 放进 几何与表面重建 的协议、系统或应用边界里：该工作被召回是因为论文文本包含 Gaussian/3D Gaussian 相关表述；当前分类先确定其任务域和与 3DGS 的相关性，后续证据流水线再补代码、企业和引用信息。
\textbf{实现细节：} 具体实现可以按输入、核心模块和渲染/优化路径来读。输入端首先取得任务所需的图像、相机、文本、时间序列、语义或几何线索，并把这些线索转换成可以约束高斯中心、尺度、旋转、不透明度、颜色或附加特征的条件。核心结构中，方法使用的主要机制是：从摘要层面看，Generalizing Visual Geometry Priors to Sparse Gaussian Occupancy Prediction 围绕“把 3DGS 用于几何、表面、深度或网格重建”展开。若其正文或代码后续可读，应进一步核对它具体改动的是表示、初始化、渲染、优化、损失、动态建模、压缩或应用接口中的哪一环。 训练或推理时，这些模块通常先给出初始化或中间特征，再通过可微 Gaussian rasterization、重建损失、感知/语义约束、几何正则或任务损失回传到高斯参数；如果方法包含前馈网络，则网络输出的是高斯属性、变形场、压缩码、语义特征或控制变量，而不是只输出最终 2D 图像。渲染路径仍以显式高斯投影、深度排序和 alpha 合成为基础，只是在采样、属性预测、局部编辑、动态变形、压缩编码或下游任务接口处加入额外约束。这样写入本章时，它与前后方法的关系是：相对原始 3DGS 和本章已有锚点方法，它可以理解为 3DGS extends\_or\_uses 宽召回分类阶段的关系判断，后续需要按论文方法部分复核。
\textbf{证据与边界：} 本轮召回分类尚未逐项核对实验表格；后续证据补全阶段需要读取论文源码/PDF或项目页确认数据集、指标和对比方法。 该条来自宽召回批次；源码下载受 arXiv e-print 网络限制，部分论文暂时只能依据标题、摘要和类别分类，需要在入可信书前做人工/LLM精读复核。 本书只采用本地 paper pack、数据库分类、源码/引用登记和论文摘要中可核验的信息；没有本地证据支撑的速度、排行榜或开源完整性不在正文中强化。\textbf{方法信息：} 代码：\url{https://github.com/juivyy/gpocc}。
\textbf{设计动机与边界：} 相对原始 3DGS 以独立高斯和图像重建误差为主的设定，Generalizing Visual Geometry Priors to Sparse Gaussian Occupancy Prediction 把问题落在以下可实现的改动上：从摘要层面看，Generalizing Visual Geometry Priors to Sparse Gaussian Occupancy Prediction 围绕“把 3DGS 用于几何、表面、深度或网格重建”展开。若其正文或代码后续可读，应进一步核对它具体改动的是表示、初始化、渲染、优化、损失、动态建模、压缩或应用接口中的哪一环。其设计目的不是单独增加一个模块，而是用上述表示、监督或渲染约束消除原流程在目标任务中的具体失配。该工作被召回是因为论文文本包含 Gaussian/3D Gaussian 相关表述；当前分类先确定其任务域和与 3DGS 的相关性，后续证据流水线再补代码、企业和引用信息。适用时还应保留论文给出的边界：该条来自宽召回批次；源码下载受 arXiv e-print 网络限制，部分论文暂时只能依据标题、摘要和类别分类，需要在入可信书前做人工/LLM精读复核。
% METHOD: 2603.27228
\paragraph{NimbusGS（几何、法线、表面与网格）（2026.03）}
NimbusGS~\cite{2603_27228} 关注的问题是：把 3DGS 用于几何、表面、深度或网格重建 NimbusGS 是宽召回补全中发现的论文，主要根据标题、摘要和 arXiv 类别判定其与 3DGS 的关系。
相对原始 3DGS，NimbusGS 的变化不是只替换损失函数，而是把表示、初始化、约束或网络预测路径重新组织到任务流程中：该工作被召回是因为论文文本包含 Gaussian/3D Gaussian 相关表述；当前分类先确定其任务域和与 3DGS 的相关性，后续证据流水线再补代码、企业和引用信息。
\textbf{实现细节：} 具体实现可以按输入、核心模块和渲染/优化路径来读。输入端首先取得任务所需的图像、相机、文本、时间序列、语义或几何线索，并把这些线索转换成可以约束高斯中心、尺度、旋转、不透明度、颜色或附加特征的条件。核心结构中，方法使用的主要机制是：从摘要层面看，NimbusGS 围绕“把 3DGS 用于几何、表面、深度或网格重建”展开。若其正文或代码后续可读，应进一步核对它具体改动的是表示、初始化、渲染、优化、损失、动态建模、压缩或应用接口中的哪一环。 训练或推理时，这些模块通常先给出初始化或中间特征，再通过可微 Gaussian rasterization、重建损失、感知/语义约束、几何正则或任务损失回传到高斯参数；如果方法包含前馈网络，则网络输出的是高斯属性、变形场、压缩码、语义特征或控制变量，而不是只输出最终 2D 图像。渲染路径仍以显式高斯投影、深度排序和 alpha 合成为基础，只是在采样、属性预测、局部编辑、动态变形、压缩编码或下游任务接口处加入额外约束。这样写入本章时，它与前后方法的关系是：相对原始 3DGS 和本章已有锚点方法，它可以理解为 3DGS extends\_or\_uses 宽召回分类阶段的关系判断，后续需要按论文方法部分复核。
\textbf{证据与边界：} 本轮召回分类尚未逐项核对实验表格；后续证据补全阶段需要读取论文源码/PDF或项目页确认数据集、指标和对比方法。 该条来自宽召回批次；源码下载受 arXiv e-print 网络限制，部分论文暂时只能依据标题、摘要和类别分类，需要在入可信书前做人工/LLM精读复核。 本书只采用本地 paper pack、数据库分类、源码/引用登记和论文摘要中可核验的信息；没有本地证据支撑的速度、排行榜或开源完整性不在正文中强化。\textbf{方法信息：} 代码：\url{https://github.com/lyy-ovo/nimbusgs}。
\textbf{设计动机与边界：} 相对原始 3DGS 以独立高斯和图像重建误差为主的设定，NimbusGS 把问题落在以下可实现的改动上：从摘要层面看，NimbusGS 围绕“把 3DGS 用于几何、表面、深度或网格重建”展开。若其正文或代码后续可读，应进一步核对它具体改动的是表示、初始化、渲染、优化、损失、动态建模、压缩或应用接口中的哪一环。其设计目的不是单独增加一个模块，而是用上述表示、监督或渲染约束消除原流程在目标任务中的具体失配。该工作被召回是因为论文文本包含 Gaussian/3D Gaussian 相关表述；当前分类先确定其任务域和与 3DGS 的相关性，后续证据流水线再补代码、企业和引用信息。适用时还应保留论文给出的边界：该条来自宽召回批次；源码下载受 arXiv e-print 网络限制，部分论文暂时只能依据标题、摘要和类别分类，需要在入可信书前做人工/LLM精读复核。
% METHOD: 2607.04923
\paragraph{UniSpine-GS（几何、法线、表面与网格）（2026.07）}
UniSpine-GS~\cite{2607_04923} 关注的问题是：把 3DGS 用于几何、表面、深度或网格重建 UniSpine-GS 是宽召回补全中发现的论文，主要根据标题、摘要和 arXiv 类别判定其与 3DGS 的关系。
作为 supplementary 类条目，UniSpine-GS 不应被写成单纯的新渲染算子。它更像是把 3DGS 放进 几何与表面重建 的协议、系统或应用边界里：该工作被召回是因为论文文本包含 Gaussian/3D Gaussian 相关表述；当前分类先确定其任务域和与 3DGS 的相关性，后续证据流水线再补代码、企业和引用信息。
\textbf{实现细节：} 具体实现可以按输入、核心模块和渲染/优化路径来读。输入端首先取得任务所需的图像、相机、文本、时间序列、语义或几何线索，并把这些线索转换成可以约束高斯中心、尺度、旋转、不透明度、颜色或附加特征的条件。核心结构中，方法使用的主要机制是：从摘要层面看，UniSpine-GS 围绕“把 3DGS 用于几何、表面、深度或网格重建”展开。若其正文或代码后续可读，应进一步核对它具体改动的是表示、初始化、渲染、优化、损失、动态建模、压缩或应用接口中的哪一环。 训练或推理时，这些模块通常先给出初始化或中间特征，再通过可微 Gaussian rasterization、重建损失、感知/语义约束、几何正则或任务损失回传到高斯参数；如果方法包含前馈网络，则网络输出的是高斯属性、变形场、压缩码、语义特征或控制变量，而不是只输出最终 2D 图像。渲染路径仍以显式高斯投影、深度排序和 alpha 合成为基础，只是在采样、属性预测、局部编辑、动态变形、压缩编码或下游任务接口处加入额外约束。这样写入本章时，它与前后方法的关系是：相对原始 3DGS 和本章已有锚点方法，它可以理解为 3DGS extends\_or\_uses 宽召回分类阶段的关系判断，后续需要按论文方法部分复核。
\textbf{证据与边界：} 本轮召回分类尚未逐项核对实验表格；后续证据补全阶段需要读取论文源码/PDF或项目页确认数据集、指标和对比方法。 该条来自宽召回批次；源码下载受 arXiv e-print 网络限制，部分论文暂时只能依据标题、摘要和类别分类，需要在入可信书前做人工/LLM精读复核。 本书只采用本地 paper pack、数据库分类、源码/引用登记和论文摘要中可核验的信息；没有本地证据支撑的速度、排行榜或开源完整性不在正文中强化。\textbf{方法信息：} 代码：\url{https://github.com/orangeisland66/unispine-gs}。


几何章节的主线是从“可渲染高斯”走向“可验证表面”。DN-Splatter、CDGS、VCR-GauS 和 VA-GS 用深度/法线/置信度损失补光度监督；2DGS、PGSR 和 RaDe-GS 重新定义或改进深度与法线渲染；SuGaR、GOF、GS2Mesh 和 MeshSplatting 把高斯推进到点云/网格/三角形输出；GSDF、3DGSR、GSurf 和 GaussianUDF 用连续隐式场补高斯的离散缺陷；TrimGS、Spiking GS、CurveGaussian 和 Planar GS 则分别处理冗余高斯、尾部重叠、曲线和平面结构。读这些方法时，关键不是只看 Chamfer 是否下降，而是看它改变的是初始化、基元形状、深度定义、正则化信号、可见性建模，还是最终表面提取方式。

\textbf{设计动机与边界：} 相对原始 3DGS 以独立高斯和图像重建误差为主的设定，UniSpine-GS 把问题落在以下可实现的改动上：从摘要层面看，UniSpine-GS 围绕“把 3DGS 用于几何、表面、深度或网格重建”展开。若其正文或代码后续可读，应进一步核对它具体改动的是表示、初始化、渲染、优化、损失、动态建模、压缩或应用接口中的哪一环。其设计目的不是单独增加一个模块，而是用上述表示、监督或渲染约束消除原流程在目标任务中的具体失配。该工作被召回是因为论文文本包含 Gaussian/3D Gaussian 相关表述；当前分类先确定其任务域和与 3DGS 的相关性，后续证据流水线再补代码、企业和引用信息。适用时还应保留论文给出的边界：该条来自宽召回批次；源码下载受 arXiv e-print 网络限制，部分论文暂时只能依据标题、摘要和类别分类，需要在入可信书前做人工/LLM精读复核。
\textbf{尾项实现澄清：} UniSpine-GS 的动机是避免高成本三维医学采集以及 CT、超声等模态间物理差异造成的泛化失配。它不以显式体重建为目标，而是学习几何感知高斯表示进行多视图投影渲染；SPWM 按脊柱结构重加权监督，使边界和局部解剖细节得到更强约束。摘要未公开完整网络和损失权重，因此可信版只确认其结构引导重加权与跨视图解剖一致性目标，不虚构优化细节。
